% Generated by maint/texbuild/theory_tex.py from theory/workbooks/engine/engine_table.md.
% Edits here are overwritten. Edit the markdown, then run the script again.
\chapter{\texorpdfstring{The engine table}{The engine table}}
\label{chap:engine-table}

\textbf{Purpose:} The machine, part by part: what each part computes, the exact algebra it holds to, what it does today, what it wants, every hypothesis tried there with its result, and the next move.

\textbf{Layout (Doug, 23 September):} \texttt{engine/\allowbreak{}base/\allowbreak{}} never sees a body, \texttt{engine/\allowbreak{}nbody/\allowbreak{}} is defined over n ≥ 1 bodies, \texttt{engine/\allowbreak{}prg\_\allowbreak{}sch/\allowbreak{}} composes the runs (answer\_\allowbreak{}key, cfg, nbody\_\allowbreak{}program, run\_\allowbreak{}cfg, run\_\allowbreak{}log), \texttt{engine/\allowbreak{}render/\allowbreak{}} draws a step's state, and \texttt{engine/\allowbreak{}sims/\allowbreak{}} holds the sims (M19). The benches are the root \texttt{bench/\allowbreak{}}, the tests the root \texttt{test/\allowbreak{}}, and every build writes only into the root \texttt{build/\allowbreak{}}: each build into its own \texttt{build/\allowbreak{}<date>\_\allowbreak{}<time>\_\allowbreak{}<name>/\allowbreak{}}, removed when a later build starts more than a day on, with the finished program copied up into \texttt{build/\allowbreak{}} and \texttt{build/\allowbreak{}logs/\allowbreak{}} and \texttt{build/\allowbreak{}tsv/\allowbreak{}} never removed. The six steps (ingest, seal, lift, measure, partition, relate) and every file's step are in engine/README.md.

\textbf{Scope:} \texttt{engine/\allowbreak{}}: exact arithmetic, the residual operator, the component tree, the moments, the correlation operator, the marginal, matching, the entropy history, the files (their size against the floor is tracked in compression\_\allowbreak{}table.md), the record machine, errors, the seal (integrity; its theory is \hyperref[chap:obsignatio-seal]{obsignatio\_\allowbreak{}seal.md}), the scheduler, the sift, render, the memory operations, the period reading, the sims, the root universal, and the ask and the state. The programs that run on it are separate concerns with their own tables: cell tracking in cell\_\allowbreak{}tracking\_\allowbreak{}table.md, and RSNA knee MRI (its own session). Numbers are in \hyperref[chap:ledger]{ledger.md}. Statuses follow \hyperref[chap:readme]{README.md}.

\section{\texorpdfstring{The rule the machine is held to}{The rule the machine is held to}}

\textbf{The engine is optimized for no scale} (Doug, 23 September). A size, a spacing, an order, a window, a width, a cell or a voxel count is never written into the machine. Each comes in with the request or is read from the data. The machine is exact at every scale its words can hold, and where a word is too narrow it says so (a request error, or \texttt{needed\_\allowbreak{}bits}) and never rounds. Tuning the machine to one program's scale is the error this table exists to catch: it has been put in and ripped out repeatedly, and the audit below lists what is still in.

The machine and each program are optimized against each other only through the interface at the end of this table. Merged, neither can be optimized on its own.

\section{\texorpdfstring{The engine in math}{The engine in math}}

Doug, 24 September: “write the whole fucking engine algo out in math so I CAN SEE WHERE YOU LACK UNDERSTANDING”, then “ask your theorist, read the workbook, make sure you miss zero terms”. Sources: every workbook file (this table, \hyperref[chap:vertical-time-compression]{vertical\_\allowbreak{}time\_\allowbreak{}compression.md}, \hyperref[chap:two-crystals]{two\_\allowbreak{}crystals.md}, \hyperref[chap:imprint-key-cycle]{imprint\_\allowbreak{}key\_\allowbreak{}cycle.md}, \hyperref[chap:keys-explained]{keys\_\allowbreak{}explained.md}, \hyperref[chap:kolmogorov-arnold]{kolmogorov\_\allowbreak{}arnold.md}, \hyperref[chap:noise-sieve-tower]{noise\_\allowbreak{}sieve\_\allowbreak{}tower.md}, \hyperref[chap:obsignatio-seal]{obsignatio\_\allowbreak{}seal.md}, \hyperref[chap:tessera-scheduler]{tessera\_\allowbreak{}scheduler.md}, compression\_\allowbreak{}table.md) and the code read this session (\texttt{tower.\allowbreak{}cu}, \texttt{keymath.\allowbreak{}cu}, \texttt{cycle.\allowbreak{}c}, \texttt{pi\_\allowbreak{}tower.\allowbreak{}cu}). Each line carries the status its source gives it. What is not tied to the code is listed at the end.

\textbf{Objects.} A sample is u : Λ → [0, 2\textasciicircum{}16), Λ = [0,T)×[0,Z)×[0,Y)×[0,X), axes (t, z, y, x). A frame is u\_\allowbreak{}t = u(t, ·). The atom is the whole of the thing under inspection, held as one integer: A = Σ\_\allowbreak{}i v\_\allowbreak{}i·2\textasciicircum{}\{16i\}, lane i at place 2\textasciicircum{}\{16i\}. Every value is in ℤ; nothing is rounded.

\textbf{E0. Numbers} (M1, A0).

\begin{itemize}
\item Width: 2\textasciicircum{}k·32 bits, k ≥ 0, no ceiling (256 limbs in the engine build, fitted to the record file).
\item a·b: long multiplication, then Karatsuba from 32 limbs, then Schönhage–Strassen mod 2\textasciicircum{}n + 1 from 8,192 limbs (every twiddle a shift).
\item a = q·b + r, q = trunc(a/b), |r| < |b|, sgn r = sgn a: Knuth's D, and Newton's reciprocal once divisor and quotient reach 8,192 limbs. b = 0 is refused.
\item Exact a/b: b = 2\textasciicircum{}s·o with o odd. The quotient is ((a ≫ s)·o⁻¹) mod 2\textasciicircum{}\{32·limbs\}, with o⁻¹ from x ← x(2 − o·x); it is refused unless q·b = a.
\item gcd: Lehmer's (Knuth L); gcd(0, x) = |x|, gcd(0, 0) = 0.
\item A decimal (−1)\textasciicircum{}neg·digits·10\textasciicircum{}ex maps to ⌊|d|·2\textasciicircum{}\{−e2\}⌋ in W bits, with e2, \texttt{terminates}, \texttt{fits} and \texttt{needed\_\allowbreak{}bits}.
\end{itemize}

\textbf{E1. The record machine} (M10, A13): the evaluator every engine program runs on.

\begin{itemize}
\item A program is P = (s\_\allowbreak{}1, …, s\_\allowbreak{}n), s\_\allowbreak{}i = (op\_\allowbreak{}i, l\_\allowbreak{}i, r\_\allowbreak{}i, μ\_\allowbreak{}i), with l\_\allowbreak{}i, r\_\allowbreak{}i < i, and outputs O ⊆ \{1..n\}, no output twice.
\item A lane ℓ reads up to 3 records ρ\_\allowbreak{}μ = in\_\allowbreak{}μ[idx\_\allowbreak{}μ(ℓ)] (idx the identity when none is given). Registers v\_\allowbreak{}1..v\_\allowbreak{}n ∈ ℤ, v\_\allowbreak{}i = op\_\allowbreak{}i(v\_\allowbreak{}\{l\_\allowbreak{}i\}, v\_\allowbreak{}\{r\_\allowbreak{}i\}):

{\small
\setlength{\tabcolsep}{6pt}
\begin{longtable}{>{\RaggedRight\arraybackslash}p{\dimexpr0.2896\linewidth-2\tabcolsep\relax}>{\RaggedRight\arraybackslash}p{\dimexpr0.3313\linewidth-2\tabcolsep\relax}>{\RaggedRight\arraybackslash}p{\dimexpr0.3791\linewidth-2\tabcolsep\relax}}
\toprule
op & v\_\allowbreak{}i & width w\_\allowbreak{}i (keymath) \\
\midrule
\endhead
FIELD & bits [o, o+b) of ρ\_\allowbreak{}μ, unsigned & b \\
FIELD\_\allowbreak{}SIGNED & the same, signed & b \\
CONSTANT & c < 2\textasciicircum{}64 & bitlen(c) \\
SUM, DIFFERENCE & a ± b & the fewer of max(w\_\allowbreak{}a, w\_\allowbreak{}b) + 1 and the linear form's bound \\
PRODUCT & a·b & w\_\allowbreak{}a + w\_\allowbreak{}b; by a constant, the fewer of that and the linear form's bound \\
ABSOLUTE & |a| & w\_\allowbreak{}a \\
COMPARE & sgn(a − b) & 1 \\
LADDER (b > 0) & sgn(a)·\#\{k ∈ [1, 91] : F\_\allowbreak{}k·b ≤ |a|\}, F\_\allowbreak{}1 = F\_\allowbreak{}2 = 1 & bitlen(91) \\
QUOTIENT & trunc(a/b) & w\_\allowbreak{}a − ⌊log2 c⌋ (≥ 1) for a constant c, else w\_\allowbreak{}a \\
REMAINDER & a − b·trunc(a/b) & min(w\_\allowbreak{}a, w\_\allowbreak{}b) \\
GCD & gcd(a, b) ≥ 0 & max(w\_\allowbreak{}a, w\_\allowbreak{}b) \\
EXACT\_\allowbreak{}QUOTIENT & a/b, b | a, else the lane is refused & as QUOTIENT \\
XOR, AND & bitwise on two's complement without end & max + 1; AND with one never-negative operand = that operand's width, with two = the narrower; XOR of two never-negative = max \\
WRAP (w ≥ 4) & ((a + 2\textasciicircum{}\{w−1\}) mod 2\textasciicircum{}w) − 2\textasciicircum{}\{w−1\} & min(w\_\allowbreak{}a, w) \\
TABLE & T[|a| mod 2\textasciicircum{}b], 1 ≤ b ≤ min(32, w\_\allowbreak{}a) & out\_\allowbreak{}bits \\
\bottomrule
\end{longtable}
}
\item The linear form (A16, 25 September): every register is c + Σ c\_\allowbreak{}i·v\_\allowbreak{}\{a\_\allowbreak{}i\} over atoms a\_\allowbreak{}i. SUM and DIFFERENCE add their operands' forms, PRODUCT by a constant (a form with no atoms) scales the other's, CONSTANT is its value, and a WRAP that passes its register through keeps that register's form. Every other register, FIELD included, is an atom. The bound is bitlen(|c| + Σ |c\_\allowbreak{}i|(2\textasciicircum{}\{w\_\allowbreak{}\{a\_\allowbreak{}i\}\} − 1)). A coefficient or constant past 2\textasciicircum{}62 makes the register an atom.
\item Never negative (N): FIELD, CONSTANT, ABSOLUTE, GCD and TABLE; SUM, PRODUCT, QUOTIENT, EXACT\_\allowbreak{}QUOTIENT or XOR of two in N; REMAINDER whose numerator is in N; AND with either in N; a WRAP that passes one in N through.
\item Refused at imprint: an operand that is not earlier, w\_\allowbreak{}i > 32·256, a WRAP below 4, a table index wider than its source. Refused per lane: b = 0, or an inexact EXACT\_\allowbreak{}QUOTIENT. A refused lane adds 1 to a device count, and a count ≠ 0 refuses the whole sweep as a request error (\texttt{cycle.\allowbreak{}cu}).
\item Out: ⊕\_\allowbreak{}\{i∈O\} v\_\allowbreak{}i, each as two's complement over w\_\allowbreak{}i + 1 bits, the sum below 2\textasciicircum{}31 bits.
\item Layout: register i lives on [i, last reader of i]. With reuse, first fit in the file, ≤ 256 limbs; without it, one register per step at the running total. The kernel's file is 64 limbs when that holds the layout, else 256. A dividing program carries 5·WIDE + 4 limbs of scratch.
\item Sweep: R\_\allowbreak{}ℓ = P(ρ(ℓ)) for every ℓ, grid-stride, ≤ 4096 × 256 threads, one launch. The host runs P on the exact integer, and device = host word for word is the port's proof.
\item A member's field lies in its first 2\textasciicircum{}32 bits.
\item Laws:

\begin{itemize}
\item (h∘g)∘f = h∘(g∘f), g∘f ≠ f∘g. A stack of floors is one program. Time per lane ∝ n, not depth. Stacked against chained measured 13 to 22 times faster.
\item WRAP to w is π\_\allowbreak{}w : ℤ₂ → ℤ/2\textasciicircum{}w, window W\_\allowbreak{}w = [−2\textasciicircum{}\{w−1\}, 2\textasciicircum{}\{w−1\}); π\_\allowbreak{}w ∘ π\_\allowbreak{}\{w+1\} = π\_\allowbreak{}w.
\item For F over \{SUM, DIFFERENCE, PRODUCT, XOR, AND\}: π\_\allowbreak{}w ∘ F = F\_\allowbreak{}w ∘ π\_\allowbreak{}w (proved).
\item These do not factor through π\_\allowbreak{}w: ABSOLUTE, COMPARE, QUOTIENT, REMAINDER, GCD, LADDER, TABLE, ⌊v/2\textasciicircum{}k⌋ (through π\_\allowbreak{}\{w+k\}), and EXACT\_\allowbreak{}QUOTIENT by an even c.
\item For c odd with c | v: WRAP(v/c, w) = WRAP(WRAP(v, w)·(c⁻¹ mod 2\textasciicircum{}w), w). The kernel grows c⁻¹ by Newton: 6, 12, 24, 48 bits in a word, then doubling.
\item Mod an odd p\textasciicircum{}v, REMAINDER is the projection ℤ\_\allowbreak{}p → ℤ/p\textasciicircum{}v, and SUM, DIFFERENCE and PRODUCT commute with it. The CRT join y\_\allowbreak{}2 + 2\textasciicircum{}8·(((y\_\allowbreak{}p − y\_\allowbreak{}2)·2\textasciicircum{}\{−8\}) rem m) is the run mod 2\textasciicircum{}8·m. XOR breaks mod 3.
\item τ\_\allowbreak{}k(x) = trunc(x/2\textasciicircum{}k): τ\_\allowbreak{}j∘τ\_\allowbreak{}k = τ\_\allowbreak{}\{j+k\}; order is kept weakly; τ\_\allowbreak{}k(x + y) − τ\_\allowbreak{}k(x) − τ\_\allowbreak{}k(y) ∈ \{−1, 0, 1\}.
\end{itemize}
\item Keys, the linear programs (imprint\_\allowbreak{}key\_\allowbreak{}cycle, keys\_\allowbreak{}explained):

\begin{itemize}
\item key = P(δ), δ the impulse. For linear shift-invariant P, P(x) = key ∗ x.
\item w·x = Σ over the set bits b of w of x ≪ b.
\item With lanes wide enough that none carries, A·K holds the smoothed lanes, and K = K\_\allowbreak{}z·K\_\allowbreak{}y·K\_\allowbreak{}x.
\item A lane below 2\textasciicircum{}m under a row of sum S stays below 2\textasciicircum{}\{m + bitlen(S − 1)\}.
\item The CRC-64 key: the advance over 2\textasciicircum{}k zero bytes is A\textasciicircum{}\{2\textasciicircum{}k\}, 48 matrices (24,576 bytes), and crc(a ⊕ b) = crc(a) ⊕ crc(b) ⊕ crc(0).
\item The mirror fold repeats with period 2N, so one period is laid as a table of offsets.
\end{itemize}
\item Tables: (g∘f)[x] = g[f[x]]. Two never-negative u, v with b\_\allowbreak{}u + b\_\allowbreak{}v ≤ 32 index as u·2\textasciicircum{}\{b\_\allowbreak{}v\} + v. A table is 2\textasciicircum{}b·⌈out\_\allowbreak{}bits/32⌉·4 bytes.
\item Not built: the lane index as a register; the latch, first lane = min\{ℓ : cond(ℓ)\}.
\end{itemize}

\textbf{E2. Ingest} (M9). Source → (u, σ). A signed DICOM pixel p ↦ p + 2\textasciicircum{}15. Slices are ordered by an exact key from their positions (E0). σ = the headers, deflated by our LZ77 and Huffman.

\textbf{E3. The seal} (M13, A10).

\begin{itemize}
\item K\_\allowbreak{}ℓ = BLAKE3-derive-key(“obsignatio aeterna 2026-09-23 “ ‖ ℓ), H\_\allowbreak{}ℓ(m) = BLAKE3-keyed(K\_\allowbreak{}ℓ, m). The levels ℓ: row, plane, volume, lanes, chunk, stream, side stored, side inflated, side, members, sample, set, file, universal.
\item ρ(t, z, y) = H\_\allowbreak{}row(u[t, z, y, ·]), π(t, z) = H\_\allowbreak{}plane(ρ(t, z, ·)), υ(t) = H\_\allowbreak{}volume(π(t, ·)), Λ = H\_\allowbreak{}lanes(υ(·)).
\item κ(c) = H\_\allowbreak{}chunk(LE64(ℓ\_\allowbreak{}c) ‖ bits\_\allowbreak{}c); Σ = H\_\allowbreak{}stream(LE64(o\_\allowbreak{}0 … o\_\allowbreak{}\{C−1\}) ‖ LE64(B) ‖ κ(0) ‖ … ‖ κ(C−1)).
\item σ = H\_\allowbreak{}side(H\_\allowbreak{}side stored(deflated) ‖ H\_\allowbreak{}side inflated(inflated)), μ = H\_\allowbreak{}members(the six tables ‖ the names).
\item Ω = H\_\allowbreak{}sample(LE64(extent, C, B, lane offset, members, side bytes, deflated bytes, name bytes, lane nodes) ‖ Λ ‖ Σ ‖ σ ‖ μ). Θ = H\_\allowbreak{}set(Ω\_\allowbreak{}1 ‖ … ‖ Ω\_\allowbreak{}n), in the order named.
\item Check order: κ, Σ, σ\_\allowbreak{}s, μ, then σ\_\allowbreak{}i, σ, then ρ … Λ, then Ω.
\item Locality: a change at (t, z, y, x) changes exactly ρ, π, υ, Λ, Ω and Θ on its path; a change in chunk c changes κ(c), Σ, Ω and Θ. Any other node agrees except with probability 2\textasciicircum{}−256.
\item The device fold (pair each level, carry the odd node up) is BLAKE3's left-heavy tree. H is not associative, so the shape is fixed.
\item The witness: the decoded voxels against the source on the device, then a second independent read, pixel for pixel.
\item Ruled, not built: H\_\allowbreak{}universal over every file the engine touches plus the engine's own hash.
\end{itemize}

\textbf{E4. The lift} (M9, M20, A8, A14).

\begin{itemize}
\item One level along one axis of length N, on the active block:

\begin{itemize}
\item d\_\allowbreak{}j = x\_\allowbreak{}\{2j+1\} − ⌊(x\_\allowbreak{}\{2j\} + x\_\allowbreak{}\{2j+2\})/2⌋, j < ⌊N/2⌋;
\item s\_\allowbreak{}i = x\_\allowbreak{}\{2i\} + ⌊(d\_\allowbreak{}\{i−1\} + d\_\allowbreak{}i + 2)/4⌋, i < ⌈N/2⌉;
\item an edge repeats its neighbor: x\_\allowbreak{}\{2j+2\} := x\_\allowbreak{}\{2j\} at the end, d\_\allowbreak{}\{−1\} := d\_\allowbreak{}0, d\_\allowbreak{}\{⌊N/2⌋\} := d\_\allowbreak{}\{i−1\};
\item lows stored on [0, ⌈N/2⌉), highs on [⌈N/2⌉, N);
\item ⌊v/2\textasciicircum{}k⌋ is taken toward −∞ (\texttt{tower\_\allowbreak{}floor\_\allowbreak{}shift}).
\end{itemize}
\item A level runs the axes 0..3 whose extent is ≥ 2, in order. The extents then go e ↦ ⌈e/2⌉, and the next level lifts only the all-low corner. There are L levels, until every extent is 1.
\item T⁻¹: x\_\allowbreak{}\{2i\} = s\_\allowbreak{}i − ⌊(d\_\allowbreak{}\{i−1\} + d\_\allowbreak{}i + 2)/4⌋, then x\_\allowbreak{}\{2j+1\} = d\_\allowbreak{}j + ⌊(x\_\allowbreak{}\{2j\} + x\_\allowbreak{}\{2j+2\})/2⌋. Levels and axes run in reverse.
\item Edges: E\_\allowbreak{}k(w) = w − (w mod 2\textasciicircum{}b) + π\_\allowbreak{}k(w mod 2\textasciicircum{}b), π\_\allowbreak{}k a permutation of [0, 2\textasciicircum{}b), 1 ≤ b ≤ 20.

\begin{itemize}
\item Lift = E\_\allowbreak{}L ∘ T\_\allowbreak{}L ∘ … ∘ E\_\allowbreak{}1 ∘ T\_\allowbreak{}1 ∘ E\_\allowbreak{}0. Lower is the inverse, last laid first.
\item Two edges on one floor fold into π\_\allowbreak{}b ∘ π\_\allowbreak{}a; a level between them blocks the fold.
\item A non-bijective f rides as (x, y) ↦ (x, y ⊕ f(x)) (Bennett), at the price of b\_\allowbreak{}x + b\_\allowbreak{}y ≤ 20 bits.
\end{itemize}
\item |c| ≥ 2\textasciicircum{}30 refuses the lift.
\item Stored (\texttt{compression.\allowbreak{}cu}):

\begin{itemize}
\item zigzag: z(c) = 2c for c ≥ 0, 2|c| − 1 for c < 0;
\item per block of 64: g = bitlen(⌊Σz/64⌋); k is the least-cost k ∈ [max(g − 2, 0), min(g + 1, 31)], with cost(k) = 5 + Σ\_\allowbreak{}z (⌊z/2\textasciicircum{}k⌋ < 24 ? ⌊z/2\textasciicircum{}k⌋ + 1 + k : 24 + 32);
\item written as k in 5 bits, then per value q = ⌊z/2\textasciicircum{}k⌋ ones and a zero and the low k bits, or 24 ones and the 32-bit value;
\item 64 blocks a chunk, and chunks sit at stored bit offsets, so any floor and any chunk is read directly (the elevator E(f, t)).
\end{itemize}
\item Held when T⁻¹(decode(stream)) = u lane for lane, and u equals a second read of the source.
\item Properties (proved along one line on 64 samples at L = 4 unless marked):

\begin{itemize}
\item T is a bijection of ℤ\textasciicircum{}n, a homeomorphism of ℤ₂\textasciicircum{}n (derived), and keeps Haar measure.
\item Reach: T's level-ℓ lows read 3ℓ bits down, its highs 3ℓ − 2; T⁻¹ reads L + 2. So |T(x) − T(y)|₂ ≤ 2\textasciicircum{}\{3L\}|x − y|₂.
\item Linear part M, entries in ℤ[1/2]: T(x + 2\textasciicircum{}\{3L\}z) = T(x) + M·2\textasciicircum{}\{3L\}z. A constant c moves only the level-L lows, by c. Negation and doubling do not pass. det M = +1.
\item Haar: every output quantum at level w holds exactly 2\textasciicircum{}\{3Ln\} input quanta at level w + 3L.
\item ℝ: 2\textasciicircum{}\{−k\}·T(trunc(2\textasciicircum{}k x)) → Mx. ℤ\_\allowbreak{}p, p odd: T does not extend; M does, in SL\_\allowbreak{}n(ℤ\_\allowbreak{}p).
\item Widths: from 25 September keymath's linear forms (A16) give every register at level ℓ, low and high, w + ℓ + 1 bits: the first level adds 2, each after it 1 (derived by hand; the ring below matches). Before, a low gained 4 bits a level and a high 2. Values obey |s| ≤ 1.5B + 1 and |d| ≤ 2B (+0.585 and +1 bits a level).
\item Ring: ring\_\allowbreak{}ℓ = ring\_\allowbreak{}0 + n + 2n(1 − 2\textasciicircum{}\{−ℓ\}) for ℓ ≥ 1 (measured, \texttt{build/\allowbreak{}20260925\_\allowbreak{}013426\_\allowbreak{}record\_\allowbreak{}boundary\_\allowbreak{}test}: 1,536, 1,664, 1,696, 1,712 and 1,720 at n = 64, w = 24; it was ring\_\allowbreak{}0 + 6n(1 − 2\textasciicircum{}\{−ℓ\}) under the old rule). The mirror floor still adds n/2\textasciicircum{}ℓ (1,720, 1,712, 1,696 and 1,600), and the heap mirrors exactly across the crystal.
\item Code length: Σ\_\allowbreak{}r ⌈log₂(W\_\allowbreak{}r + 1)⌉ + Σ\_\allowbreak{}r m\_\allowbreak{}r ≥ K(x | program) − c.
\item Identity by null permutation: a lane is identified when heap(T(x)) < heap(T(σ\_\allowbreak{}i x)) for every i ≤ d. Noise passes at a rate ≤ 1/(d + 1). T(σx) = T(x) ⇔ σx = x.
\item The KA step: x ← x + Φ(Σ\_\allowbreak{}p φ\_\allowbreak{}p(x\_\allowbreak{}p)), with Φ(s) = ⌊s/2\textasciicircum{}k⌋.
\item Matching on a floor (Doug, 25 September: “say we want to find x and x is on floor 2, we have the knf, we only need to \&\&“). Floor 2 of a 64³ frame is its 16³ corner after two levels, m = n/64 values. Laid once as 16 bit planes, it gives every position holding x as ∧\_\allowbreak{}j (x\_\allowbreak{}j ? P\_\allowbreak{}j : ¬P\_\allowbreak{}j), 32 positions a word, reading no value and at most 16m/32 = n/128 plane words; a word stops once its mask is empty. \textbf{Proved} (\texttt{floor\_\allowbreak{}match}, 11 checks, 0 failed, \texttt{build/\allowbreak{}20260925\_\allowbreak{}015942\_\allowbreak{}sim\_\allowbreak{}floor\_\allowbreak{}match}): the engine's floor 2, the crystal's corner lowered as its own tower, equals the host's two-level lifting at 4,096 of 4,096; 1,024 values on floor 2 and 1,024 not, every match set equal to the host's scan; the most any query read was 1,340 plane words, against n/2 = 131,072 samples. The index costs one read of a, the lift, paid once.
\end{itemize}
\item As record floors (only in \texttt{test/\allowbreak{}record\_\allowbreak{}bitwise\_\allowbreak{}test} and \texttt{test/\allowbreak{}record\_\allowbreak{}boundary\_\allowbreak{}test}; no engine module emits them):

\begin{itemize}
\item ⌊v/2\textasciicircum{}k⌋ = EXACT\_\allowbreak{}QUOTIENT(v − AND(v, 2\textasciicircum{}k − 1), 2\textasciicircum{}k);
\item E(v) = (v − AND(v, 2\textasciicircum{}k − 1)) + TABLE\_\allowbreak{}π(AND(v, 2\textasciicircum{}k − 1));
\item F ∘ T⁻¹ is one stack, and T⁻¹ ∘ T = id;
\item T⁻¹GT is G conjugated, (T⁻¹G₂T)(T⁻¹G₁T) = T⁻¹G₂G₁T, and a G on level-ℓ highs costs ℓ levels each way;
\item a linear F: |F(T⁻¹c) − F(W⁻¹c)| ≤ L·‖F‖∞ along one line;
\item the wrap at the mirror: WRAP(v, b + 1) on each rebuilt register, b its forward twin's width.
\end{itemize}
\end{itemize}

\textbf{E5. Measure} (M2, M8, M18), per frame.

\begin{itemize}
\item Residual: L\_\allowbreak{}n = [1 1]\textasciicircum{}\{∗n\} per axis, undivided; B\_\allowbreak{}s ∗ B\_\allowbreak{}b = B\_\allowbreak{}\{s+b\}; R = 2\textasciicircum{}\{|b|\}·L\_\allowbreak{}s u\_\allowbreak{}t − L\_\allowbreak{}\{s+b\} u\_\allowbreak{}t.

\begin{itemize}
\item ΣB\_\allowbreak{}s = 2\textasciicircum{}\{|s|\}, so a constant cancels.
\item The heat identity: 4\textasciicircum{}m δ − H\textasciicircum{}m = Σ\_\allowbreak{}\{k<m\} 4\textasciicircum{}\{m−1−k\} H\textasciicircum{}k (4δ − H), with H = [1 2 1] and 4δ − H = −Δ².
\item Parity: 2δ − [1 1] = [1 −1], half a voxel off, so an odd b is refused and an odd s is held with its offset (ruling (a)).
\item The key is 2\textasciicircum{}g·B\_\allowbreak{}narrow − B\_\allowbreak{}wide·B\_\allowbreak{}narrow, in 9 limbs.
\item Readings of any declared width w (the unit sweep's planes, 25 September): R leaves in ⌈(w + Σs + Σb + 1)/32⌉ limbs, and a reading ≥ 2\textasciicircum{}w is refused.
\item The positive set: P(x) = [R(x) > 0].
\end{itemize}
\item The history:

\begin{itemize}
\item f\_\allowbreak{}j(x, W) = Σ\_\allowbreak{}\{t∈W\} bit\_\allowbreak{}j(u\_\allowbreak{}t(x) ⊕ u\_\allowbreak{}\{t−1\}(x)), |W| = 11 transitions, ≤ 16 windows;
\item held: Σ\_\allowbreak{}W f\_\allowbreak{}j ≡ bit\_\allowbreak{}j(u\_\allowbreak{}0(x) ⊕ u\_\allowbreak{}\{F−1\}(x)) (mod 2) at every x and j;
\item the anchor count a\_\allowbreak{}j(x) = Σ\_\allowbreak{}t bit\_\allowbreak{}j(u\_\allowbreak{}t(x)); the anchor bits are ∧\_\allowbreak{}t u\_\allowbreak{}t and ∧\_\allowbreak{}t ¬u\_\allowbreak{}t;
\item the cloud: D\_\allowbreak{}W(x) = Σ\_\allowbreak{}j f\_\allowbreak{}j(x, W)·2\textasciicircum{}j, and O(W, W′) = Σ\_\allowbreak{}x D\_\allowbreak{}W(x)·D\_\allowbreak{}\{W′\}(x);
\item the knf section s(x) = (D\_\allowbreak{}W(x))\_\allowbreak{}W less its mean over W; E = Σ\_\allowbreak{}x Σ\_\allowbreak{}\{a∈\{z,y,x\}\} ⟨s(x), s(x + e\_\allowbreak{}a)⟩ on the torus; the departure is d·E − Σ\_\allowbreak{}i E(σ\_\allowbreak{}i x).
\end{itemize}
\item The period, per axis of extent n, M = N/n:

\begin{itemize}
\item L = ⌊n/2⌋; the starts are the positions below n − L on the axis, Q = (n − L)·M;
\item a(ℓ) = \#\{x : u(x) = u(x + ℓ·e)\};
\item a candidate p has 2 ≤ p, 2p + 1 ≤ L, a(p) > max(a(p ± 1)) and a(2p) > max(a(2p ± 1));
\item h(p) = min(a(p) − max(a(p ± 1)), a(2p) − max(a(2p ± 1)));
\item the null draw k: each line along the axis permuted by a keyed Fisher–Yates (swap s with h mod (s + 1), 4 keys from the content signum and c = draw·rank + axis), h\_\allowbreak{}k = max\_\allowbreak{}p h;
\item P = min\{p : h(p)/Q > max\_\allowbreak{}k h\_\allowbreak{}k/Q\}, a 128-bit cross-multiply; P = 0 when none clears.
\item False-period rate ≤ 1/(d + 1).
\end{itemize}
\end{itemize}

\textbf{E6. Partition} (M3, M4, A2, A3).

\begin{itemize}
\item Admitted: R(x) > 0. The code c(x) = (dense rank of R(x) over the admitted) + 1, and 0 elsewhere.
\item Adjacency: 6 (a face joins x and x + e\_\allowbreak{}a, a ∈ \{z, y, x\}, both admitted).
\item The face key is (c of the weaker end, then ¬(3x + a) in the lowest limb), so a stronger weaker end wins and then the lower face name. The tree is the maximum spanning forest by these keys (\texttt{max\_\allowbreak{}tree\_\allowbreak{}contract}, \texttt{max\_\allowbreak{}tree.\allowbreak{}cu}):

\begin{itemize}
\item Borůvka rounds: each component takes its strongest outgoing face;
\item of a mutual pair the higher root hooks onto the lower;
\item pointer jumping between rounds;
\item the chosen faces are the forest.
\end{itemize}
\item Components at every level from the forest: count(r) = \#\{x : c(x) ≥ r\} − \#\{forest faces with min(c) ≥ r\}, one scan from the top. The cut is r* = min\{r : count(r) = max count\}. The bodies are the components of \{c ≥ r*\}, labeled by min-label spreading along the forest. The count found must equal count(r*), else it is a logic error.
\item A node n has a level λ(n), and ν(x) is x's own node.

\begin{itemize}
\item The component at level r holding x is the ancestor a of ν(x) with λ(a) ≥ r > λ(parent(a)).
\item A subtree is a DFS range [in(a), out(a)).
\item x and x′ share a component at r exactly when r ≤ λ(LCA(ν(x), ν(x′))).
\end{itemize}
\item Ties: every face's key carries the face's name in its lowest limb, so no two keys tie.
\item Moments: μ(n) = Σ\_\allowbreak{}\{ν(x)=n\} (1, x, x xᵀ), and each node's are Σ over sub(n). The centred K = m·M − S·Sᵀ is never divided. ℓ ≤ 2 harmonics = (m, S, M).
\item Theory, not built: the probe partition from k − 1 LCA levels; J(v)[n, n′] with C\_\allowbreak{}\{r,r′\} a rectangle sum over two DFS ranges; moments by one DFS prefix.
\end{itemize}

\textbf{E7. Relate} (M5, M6, M7).

\begin{itemize}
\item Agreement: A(v) = Σ\_\allowbreak{}x P\_\allowbreak{}t(x)·P\_\allowbreak{}\{t′\}(x + v), by an NTT mod 998,244,353.

\begin{itemize}
\item Each axis of extent n is padded to the least power of 2 ≥ 2n − 1, so every lag −(n − 1)..(n − 1) is linear, never wrapped.
\item It is refused when |Λ\_\allowbreak{}t| ≥ the prime or the padded total > 2\textasciicircum{}31 − 1, and the longest padded axis is ≤ 2\textasciicircum{}23.
\item v* = argmax A, ties to the least Σ\_\allowbreak{}a w\_\allowbreak{}a·v\_\allowbreak{}a², then the least index.
\end{itemize}
\item \texttt{C\_\allowbreak{}\{t,t′\}[a, b](v)} = Σ\_\allowbreak{}x [L\_\allowbreak{}t(x) = a]·[P\_\allowbreak{}\{t′\}(x + v)]·[L\_\allowbreak{}\{t′\}(x + v) = b]. Its reductions:

\begin{itemize}
\item the argmax over v;
\item the climb (\texttt{climb\_\allowbreak{}machine.\allowbreak{}cu}): S(v′) is the climber's count at lag v′, over v′ ∈ v + \{−1, 0, 1\}³. It moves to the v′ ≠ v with S(v′) > S(v), the largest, ties to the least w\_\allowbreak{}z·z² + y² + x². It stops when none exceeds S(v), and holds S(v). The landing is a majority vote (one pass picks, one counts) of the leaf under each of the climber's voxels at the lag and along a spiral around it, the most agreed kept;
\item the transpose, \texttt{C\_\allowbreak{}\{t′,t\}[b, a](−v)} = \texttt{C\_\allowbreak{}\{t,t′\}[a, b](v)};
\item box sums at nested widths from one summed-area table;
\item the tent;
\item the moments, at v = 0 with weight φ ∈ \{1, x, x xᵀ\}.
\end{itemize}
\item Marginal:

\begin{itemize}
\item Z = Σ\_\allowbreak{}α Π\_\allowbreak{}s ω(s, α(s)), α injective with ∅ allowed, is the permanent of [ω(s, t) | diag ω(s, ∅)];
\item P(s → t) = Z\_\allowbreak{}\{s→t\}/Z, and it factors over the support's components;
\item bounds: ≤ 2\textasciicircum{}16 arrangements, ≤ 64 targets, 16 limbs.
\end{itemize}
\item Matching: M* = argmax over one-to-one M of Σ\_\allowbreak{}\{(a,b)∈M\} w(a, b), in exact integers.
\item Body programs (velocity, division, contact side, fingerprint, print pair) are E1 programs over packed body records, ≤ 3 members.
\item Null draws on the identity line: frame f stands on floor ⌊f/11⌋ at place f mod 11. Its draw on floor g is 11g + (f mod 11), the floors taken by least O.
\end{itemize}

\textbf{E8. The scheduler} (M14, A11).

\begin{itemize}
\item The quantities: device C, in use U(τ); job j has d\_\allowbreak{}j, r\_\allowbreak{}j (from d\_\allowbreak{}j, only growing) and u\_\allowbreak{}j(τ) by pid.
\item O(τ) = U(τ) − Σ\_\allowbreak{}j u\_\allowbreak{}j(τ); H(τ) = C − O(τ) − Σ\_\allowbreak{}j max(r\_\allowbreak{}j, u\_\allowbreak{}j(τ)).
\item Admit k when d\_\allowbreak{}k ≤ H, with r\_\allowbreak{}k = d\_\allowbreak{}k (Doug's rule; the working tree's max(d\_\allowbreak{}k, p\_\allowbreak{}σ) is unapproved).
\item Growth: u\_\allowbreak{}j > r\_\allowbreak{}j sets r\_\allowbreak{}j = u\_\allowbreak{}j and warns. On release, p\_\allowbreak{}σ = max\_\allowbreak{}τ u\_\allowbreak{}j(τ) under the signum σ = BLAKE3(request).
\item Asked: d\_\allowbreak{}k > p\_\allowbreak{}σ. The job is held, overridable, and lost after its holding time.
\item Backfill with the head kept:

\begin{itemize}
\item the shadow S = least time with H + Σ\_\allowbreak{}\{e\_\allowbreak{}j ≤ S\} max(r\_\allowbreak{}j, u\_\allowbreak{}j) ≥ d\_\allowbreak{}h, e\_\allowbreak{}j = start\_\allowbreak{}j + t\_\allowbreak{}σj;
\item the spare X = that sum − d\_\allowbreak{}h;
\item k goes now if d\_\allowbreak{}k ≤ H and (now + t\_\allowbreak{}σk ≤ S or d\_\allowbreak{}k ≤ X);
\item a signum with no history is never backfilled.
\end{itemize}
\item Every deadline sits in one min-heap on time to change.
\item A ticket is (pid, start time).
\end{itemize}

\textbf{E9. Errors} (M12). Every entry returns a value or refuses with (kind ∈ \{request, resource, logic\}, module, site, status, execaddr, evacaddr, imagebase, ≤ 16 frames). First write wins. A logic error withholds the result.

\textbf{E10. The sift, the ladder, memory, render} (M11, M15, M16, M17).

\begin{itemize}
\item Sift dispatch: in order when 100·total² against 85·distinct·Σcount² says so, ≥ 143 bits. Anchors are 1 when a period P is found, else 4. A periodic corpus filters no lower than 1/P. H₂ is the collision entropy.
\item β(n) = \#\{Fibonacci rungs ≤ n after the first\}, 92 rungs below 2\textasciicircum{}64. LADDER(a, b) reads β of |a|/b.
\item The memory arms answer exactly what a byte loop answers; render's host bytes equal the device's.
\end{itemize}

\textbf{E11. The ask and the state} (M21, A15).

\begin{itemize}
\item p\_\allowbreak{}i = Tr(ρE\_\allowbreak{}i), E\_\allowbreak{}i = (I + v\_\allowbreak{}i·σ)/4, so p\_\allowbreak{}i = (1 + r·v\_\allowbreak{}i)/4.
\item With Σv\_\allowbreak{}i = 0 and Σv\_\allowbreak{}i v\_\allowbreak{}iᵀ = c·I, r = (4/c)·Σ p\_\allowbreak{}i v\_\allowbreak{}i.
\item Crossing: (1 + r\_\allowbreak{}a)/2 = Σ p\_\allowbreak{}i (½ + (2/c)v\_\allowbreak{}\{i,a\}).
\item Valid ⇔ |r|² ≤ 1.
\item ℚ(√3): (a + b√3)(c + d√3) = (ac + 3bd) + (ad + bc)√3, the sign from a² against 3b².
\end{itemize}

\textbf{E12. The limits and the ordered machine} (two\_\allowbreak{}crystals).

\begin{itemize}
\item The inverse limit ℤ₂ = lim← ℤ/2\textasciicircum{}w; the direct limit of ℤ/2\textasciicircum{}w under x ↦ 2x is ℤ[1/2]/ℤ, its dual. ℤ is dense in ℤ₂, and ℤ₂ \textbackslash{} ℤ is uncountable.
\item The solenoid Σ₂ = (ℝ × ℤ₂)/ℤ. The adeles 𝔸/ℚ ≅ (ℝ × Ẑ)/ℤ, Ẑ = Π\_\allowbreak{}p ℤ\_\allowbreak{}p, and |x|\_\allowbreak{}∞·Π\_\allowbreak{}p |x|\_\allowbreak{}p = 1.
\item Orders: F\textasciicircum{}k for k ∈ ℤ, F\textasciicircum{}\{−k\} = (F⁻¹)\textasciicircum{}k. The counter n ↦ n + 1 is a shear.
\item ±ω: each bit ← lim sup. On a finite window that is the OR over the cycle, and +ω = −ω.
\end{itemize}

\textbf{E13. Programs on the engine} (M19).

\begin{itemize}
\item The camera law:

\begin{itemize}
\item S = background + ramp·x + wave + Σ bodies' brightness;
\item shot noise Binomial(4S, ½) − S, read noise Binomial(4r², ½) − 2r²;
\item value = offset + fixed pattern + gain·S + read;
\item every draw is splitmix64(key ⊕ counter).
\end{itemize}
\item Ω (\texttt{chaitin\_\allowbreak{}omega}): Ω = Σ over closed BLC terms t with a normal form of 2\textasciicircum{}\{−|t|\}.

\begin{itemize}
\item The bracket through L bits: [Σ\_\allowbreak{}halted 2\textasciicircum{}\{−|t|\}, that + Σ\_\allowbreak{}open 2\textasciicircum{}\{−|t|\} + the mass past L, directed].
\item Halt: a normal form. Loop: Brent's watcher. Forever: the spine grows. Halts: simply typed.
\item Codes, de Bruijn: 00M is λM, 01MN is MN, 1\textasciicircum{}k0 is variable k. Terms are enumerated by rank: unrank(length, depth, index) over the counts c(n, k), the variable, then the lambda, then applications by the left part's length.
\item On the engine (\texttt{chaitin\_\allowbreak{}omega.\allowbreak{}cu}, the run by the engine): one normal-order step of every live term is one sweep, a term is one record a token, and there is no step or token budget. Per round:

\begin{itemize}
\item phase one, a thread a term: find the leftmost outermost redex (λM)N; a normal form settles;
\item phase two: each reduct token is a plan (depth D, bound k, body, argument) gathered from its source token by the index. The record program computes (λM)N → M[N], with a body variable i > D + 1 ↦ i − 1, an argument variable i > k ↦ i + D, and every other token copied. It is imprinted at the round's widths;
\item phase three: the same decisions as \texttt{omega\_\allowbreak{}run}, namely halt at a normal form, loop by Brent's watcher, grows forever by the spine proof;
\item settled terms leave as one record; kept terms move into the next round's rooms;
\item a job (≤ 2\textasciicircum{}22 terms of one length) is admitted once fewer than 2\textasciicircum{}21 are live;
\item a term the device cannot hold is parked as outgrown, and a stop leaves the live terms open;
\item the run is one tessera job.
\end{itemize}
\end{itemize}
\item π (\texttt{pi\_\allowbreak{}tower}), BBP at hex position d:

\begin{itemize}
\item term lane i ≤ d, for j ∈ \{1, 4, 5, 6\} with m = 8i + j: r = 16\textasciicircum{}\{d−i\} mod m, by r ← r\textasciicircum{}16·16\textasciicircum{}\{nibble\} mod m one nibble at a time;
\item f\_\allowbreak{}j = ⌊r·2\textasciicircum{}W/m⌋, and F\_\allowbreak{}i = (4f\_\allowbreak{}1 − 2f\_\allowbreak{}4 − f\_\allowbreak{}5 − f\_\allowbreak{}6) mod 2\textasciicircum{}W;
\item tail t < T: ⌊2\textasciicircum{}W/(m·16\textasciicircum{}t)⌋, the same combination; pairs F ← (F\_\allowbreak{}\{2k\} + F\_\allowbreak{}\{2k+1\}) mod 2\textasciicircum{}W;
\item S = ΣF mod 2\textasciicircum{}W, error ≤ 4(N + T) + 1 units, and the leading bits where S ± E agree are certified;
\item cost: d + 1 lanes, linear in d, ⌈(d + 1)/70,656⌉ sweeps on the RTX 3070;
\item neither tower is in it.
\end{itemize}
\item Host only: ψ (\texttt{ka\_\allowbreak{}psi}, the exponent β(r) = (nʳ − 1)/(n − 1) held symbolically) and Goodstein (hereditary trees, ordinals below ε₀).
\end{itemize}

\textbf{E14. The floor the coder answers to} is its own theory, the compression book: compression\_\allowbreak{}table.md, C0 to C13.

\textbf{E15. The governing functional} (noise\_\allowbreak{}sieve\_\allowbreak{}tower). F[U] = Π₋₄(⊕\_\allowbreak{}\{∂Φ\_\allowbreak{}new\} LUT(Key(W\_\allowbreak{}local \& L↓))):

\begin{itemize}
\item L↓ is the program;
\item W\_\allowbreak{}local is the atom;
\item \& applies the key's masks;
\item Key is the imprint;
\item LUT is the table step;
\item ⊕ is the cycle's fold;
\item Π₋₄ is the proof on the product.
\end{itemize}

\textbf{Tied to the code since the first draft:} the tree's adjacency and forest (\texttt{max\_\allowbreak{}tree.\allowbreak{}cu}), the Rice coder (\texttt{compression.\allowbreak{}cu}), the climb (\texttt{climb\_\allowbreak{}machine.\allowbreak{}cu}), the refused lane (\texttt{cycle.\allowbreak{}cu}), Ω's engine rounds (\texttt{chaitin\_\allowbreak{}omega.\allowbreak{}cu}) and the NTT's padding (\texttt{shift\_\allowbreak{}agreement.\allowbreak{}cu}). Nothing listed is left untied; the theorist's audit against the code is pending.

\section{\texorpdfstring{The table}{The table}}

\tablerecord{\textbf{M1. Exact numbers}}
\begin{description}
\item[{the algebra it holds to}] Every quantity is an exact integer or an exact rational held as a pair; nothing is rounded and no float is formed.
\item[{does today}] \texttt{Anchor\allowbreak{}Exact\allowbreak{}Integer} (anchor\_\allowbreak{}sift's) is the integer, in \texttt{engine/\allowbreak{}base/\allowbreak{}no\_\allowbreak{}rounding/\allowbreak{}exact\_\allowbreak{}integer.\allowbreak{}\{c,h\}}. Since 23 September the engine build compiles it from this tree by default (\texttt{cell\_\allowbreak{}tracking/\allowbreak{}build\_\allowbreak{}engine.\allowbreak{}sh} and \texttt{build\_\allowbreak{}driver.\allowbreak{}sh}; \texttt{ANCHOR\_\allowbreak{}EXACT\_\allowbreak{}ROOT} still overrides), and this tree is the source anchor\_\allowbreak{}sift's engine is to be copied from (anchor\_\allowbreak{}sift). The header carries their full doc text. The integer grew on 23 September (Doug, in order: “You have division”, “2s compliment”, “Mulmask exact”, “add that to the exact arithmetic”, “It makes no sense to not have infinite division with infinite mul”, a link to Schönhage–Strassen, “We use karstsuba”, “Expand the header to accept any n bits power of 2”, “Static assert”). The width is any power of two from one limb (32 bits) up, with no ceiling, named in limbs (\texttt{ANCHOR\_\allowbreak{}EXACT\_\allowbreak{}LIMBS}) or in bits (\texttt{ANCHOR\_\allowbreak{}EXACT\_\allowbreak{}BITS}). The default is 128 limbs, 4,096 bits. Static asserts refuse a width that is not a power of two, is below 32 bits, or where the two names disagree. The declared floor \texttt{ANCHOR\_\allowbreak{}EXACT\_\allowbreak{}DIGITS} stays 1,024 digits, so a width below 4,096 bits declares its own, and one that does not is refused at compile time. Width-sized working copies sit on the stack up to \texttt{ANCHOR\_\allowbreak{}EXACT\_\allowbreak{}STACK\_\allowbreak{}LIMBS} (4,096 limbs) and come from the heap past it. A copy the heap cannot hold returns \texttt{ANCHOR\_\allowbreak{}EXACT\_\allowbreak{}WILL\_\allowbreak{}NOT\_\allowbreak{}FIT}, so no width is bounded by a stack. Multiplication is a ladder: long multiplication, then Karatsuba from 32 limbs (\texttt{ANCHOR\_\allowbreak{}EXACT\_\allowbreak{}KARATSUBA\_\allowbreak{}LIMBS}), then the Schönhage–Strassen transform once the shorter factor reaches 8,192 limbs (\texttt{ANCHOR\_\allowbreak{}EXACT\_\allowbreak{}TRANSFORM\_\allowbreak{}LIMBS}). The transform works modulo 2\textasciicircum{}n + 1, where every twiddle is a shift, so nothing is rounded. An integer cost model picks each split's depth, and a rung whose workspace cannot be held steps down to the one below. \texttt{anchor\_\allowbreak{}exact\_\allowbreak{}divide} is Knuth's algorithm D in base 2\textasciicircum{}32. Once the divisor and the quotient both reach 8,192 limbs (\texttt{ANCHOR\_\allowbreak{}EXACT\_\allowbreak{}NEWTON\_\allowbreak{}LIMBS}), it switches to Newton's reciprocal: the reciprocal is grown at half precision recursively, taken one Newton step and corrected exactly, and the quotient is then one product on the ladder. Either way the quotient rounds toward zero and the remainder carries the numerator's sign, so numerator = quotient · divisor + remainder. \texttt{anchor\_\allowbreak{}exact\_\allowbreak{}divide\_\allowbreak{}exact} is for a division known to leave nothing. The divisor's twos shift out, and the rest multiplies by the inverse of the odd part modulo 2\textasciicircum{}(32 · limbs), then is masked. The inverse comes from Newton's step x(2 − dx), 2 − t being t's two's complement plus two, and each step works only to the precision it doubles to. The quotient is multiplied back, and a remainder is refused as \texttt{ANCHOR\_\allowbreak{}EXACT\_\allowbreak{}NOT\_\allowbreak{}EXACT}. \texttt{anchor\_\allowbreak{}exact\_\allowbreak{}gcd} is Lehmer's (Knuth's algorithm L). It replaced Stein's binary gcd, whose bit-at-a-time shifts stalled the 4,194,304-bit test. A zero divisor returns \texttt{ANCHOR\_\allowbreak{}EXACT\_\allowbreak{}BY\_\allowbreak{}ZERO}. \texttt{anchor\_\allowbreak{}exact\_\allowbreak{}multiply\_\allowbreak{}transform} and \texttt{anchor\_\allowbreak{}exact\_\allowbreak{}divide\_\allowbreak{}newton} run their rungs at any size, for tests and timing. The sims' rationals (\texttt{sim\_\allowbreak{}rational.\allowbreak{}h}) use the gcd and the exact division to stay in lowest terms at any width. The key machine divides too, as four record operations (M10). \texttt{decimal\_\allowbreak{}double} (23 September) expands a decimal string, digits × 10\textasciicircum{}ex, into the exact leading bits of its binary value in the candidate's full width, with its binary exponent. It says whether the value fits and terminates, and how many bits it needs when it does not (A0).
\item[{wants}] Every decimal the machine is handed (a spacing, a position, a cosine) enters through \texttt{decimal\_\allowbreak{}double}, so geometry is exact.
\item[{tried, and what it gave}] \textbf{measured} (the scratchpad \texttt{decimal\_\allowbreak{}double\_\allowbreak{}test}, a checker in exact arithmetic): 0 failures. 1e2200 needs 5,109 bits and reports fits = 0 at 4,096; the machine says “grow this” and does not round. \textbf{Proved}, the ladder and Newton's division (\texttt{test/\allowbreak{}exact\_\allowbreak{}transform\_\allowbreak{}test}, \texttt{build/\allowbreak{}20260923\_\allowbreak{}234836\_\allowbreak{}exact\_\allowbreak{}transform\_\allowbreak{}test}, 4 checks, 0 failed at each of 128 limbs, 4,096 limbs and 4,194,304 bits). The ladder and the transform alone each equal a long multiplication kept in the test, on 24, 54 and 63 products: balanced and unbalanced, keyed and all ones, across every rung boundary up to half the width. Each product divides back to its factors, exactly and with remainder, and the gcd holds a factor. Newton's division rebuilds its numerator, with the remainder below the divisor and signed as the numerator, and equals the dispatched division on 24, 54 and 63 divisions with divisors up to half the width. \textbf{Measured}, the rungs' timings (the ledger's “The exact integer: width, ladder and division”). \textbf{Proved}, the width's guards: 1, 8 and 64 limbs compile with declared floors (9, 76 and 616 digits), and 8 limbs with no floor is refused at compile time. \textbf{Proved}, the division on its own (\texttt{test/\allowbreak{}exact\_\allowbreak{}divide\_\allowbreak{}test}, \texttt{build/\allowbreak{}20260923\_\allowbreak{}234731\_\allowbreak{}exact\_\allowbreak{}divide\_\allowbreak{}test}, 9 checks, 0 failed, at 128 limbs, 4,096 bits). Long division rebuilds its numerator, with the remainder below the divisor and the signs of truncation, on 200,000 trials of edge-shaped limbs up to 8 wide and on 4,000 trials at widths drawn up to 127 limbs. It also holds on a case built to run the add-back step (anchor\_\allowbreak{}sift confirmed the branch runs, with an instrumented copy). Exact division returns the quotient of every product on 20,000 trials, odd and even divisors, and refuses one past each product. The gcd divides both values, holds their shared factor and leaves coprime cofactors on 4,000 trials, and it equals Euclid's gcd run on the long division on 4,000 more; gcd(0, x) = |x| and gcd(0, 0) = 0. A zero divisor is refused. With lowest terms, \texttt{ka\_\allowbreak{}psi} is still 152 of 152 and \texttt{ask\_\allowbreak{}state} 33 of 33 on the merged source.
\item[{status}] built; the engine build takes it from this tree
\item[{next}] Read the DICOM spacings and positions through it (M9). anchor\_\allowbreak{}sift's engine copied from this tree (item 7 below).
\end{description}

\tablerecord{\textbf{M2. The residual operator}}
\begin{description}
\item[{the algebra it holds to}] R = 2\textasciicircum{}\{|b|\} L\_\allowbreak{}s − L\_\allowbreak{}\{s+b\} over one binomial ladder L\_\allowbreak{}n = [1 1]\textasciicircum{}n ∗ I, separable per axis (A1). The kernel sums to zero exactly, so a constant added to I leaves R unchanged.
\item[{does today}] \texttt{residual\_\allowbreak{}program} builds the four steps; \texttt{keymath} imprints one key, \texttt{key\_\allowbreak{}schedule} lays it out, \texttt{cycle} runs it in one launch. \texttt{unit\_\allowbreak{}sweep} runs the same as centred [1 2 1] steps per axis in shared memory, working only the live limbs. \texttt{engine\_\allowbreak{}residual} runs either or both; both proved compares every lane, and \textbf{a disagreement is a logic error} (23 September): the result is withheld. The orders come in with the request; every step is error-wired (M12). Since 25 September the unit sweep also takes its input as limb planes of a declared width (\texttt{device\_\allowbreak{}planes}, \texttt{input\_\allowbreak{}bits}). A value past that width is refused as a request error, and the output's limbs follow from the width and the orders. \texttt{engine\_\allowbreak{}residual\_\allowbreak{}planes} hands it host planes and returns the residual with its limbs. The 16-bit volume path and the key are unchanged.
\item[{wants}] Orders from the data, per request, with nothing rounded (Doug, 23 September): a per-axis coherence reading gives an exact integer period P, and order = period. Parity by the algebra, not by a blanket guard: an odd background order misaligns the two terms by half a voxel and is refused; an odd smooth order shifts both terms by the same half voxel and is held with that offset recorded (Doug's ruling (a)). Every (s, b) from one sweep that keeps its ladder levels (A1).
\item[{tried, and what it gave}] The key against the step-by-step residual: \textbf{proved}, 0 of 10,485,760,000 lanes differ; 31 ms a frame against 24.5 for the passes (\textbf{measured}). The unit sweeps against the key: \textbf{proved}, 0 lanes differ, and about half the key's time a frame (\textbf{measured}; numbers in the cell book's on\_\allowbreak{}the\_\allowbreak{}engine.md, since they were taken on a program's samples). The planes (\texttt{test/\allowbreak{}unit\_\allowbreak{}sweep\_\allowbreak{}planes\_\allowbreak{}test}, 25 September, a tessera job): \textbf{proved}, 336 checks, 0 failed, against the 16-bit volume path by linearity. The readings are built from 16-bit parts at chosen shifts, at 16, 34, 64 and 116 bits in. On every lane their residual equals the parts' residuals shifted and summed. That is 4,357 lanes a case, 4,356 of them not zero, over three sets of orders, \{2,34,34\}/\{6,96,96\} at 10 and 13 limbs among them. Sixteen-bit planes equal the volume path. A constant added moves no lane. A bit past the width refuses, and the sweep gives the same lanes after it. Eight malformed requests refuse.
\item[{status}] proved (exact); the sweeps 2× faster; any input width proved (the planes)
\item[{next}] Ruling (a) at all three guards (\texttt{residual.\allowbreak{}cu:24}, \texttt{keymath.\allowbreak{}cu:118}, \texttt{unit\_\allowbreak{}sweep.\allowbreak{}cu:254}, and \texttt{narrow\_\allowbreak{}bits} at :259), with the half-voxel offset as a field of the result. Taking the orders from the per-axis coherence reading, now built as M18 (\texttt{period\_\allowbreak{}read}, A12); the knee pulls it. The ladder-keeping sweep.
\end{description}

\tablerecord{\textbf{M3. The component tree}}
\begin{description}
\item[{the algebra it holds to}] The dense level code, the components of every upper level set, nested as a tree. Every question about levels is a question about tree nodes: membership at level r is an ancestor query, two voxels share a component at r exactly when r ≤ level(LCA), and every subtree is one range of the DFS order (A2).
\item[{does today}] \texttt{max\_\allowbreak{}tree\_\allowbreak{}objects}: the level code, the components at one cut (the level with the most, proved against the forest's own count), the per-component moments (M4). \texttt{max\_\allowbreak{}tree\_\allowbreak{}slide} (23 September, the split): the caller's probe voxels, and each recorded level's probe partition, proved against the tree's count at that level; a disagreement is a logic error. \texttt{max\_\allowbreak{}tree\_\allowbreak{}overlap}: two frames' trees at one residual level each, the caller's probes and probe links, and the census of the bipartite support at the lag: one partner, unpaired, mutual, forked, per side. \texttt{max\_\allowbreak{}tree\_\allowbreak{}keep\_\allowbreak{}frames} holds a frame's code for the next. No cell is known to it: probes are points, links are point pairs.
\item[{wants}] Every level at once instead of one cut or sampled cuts: the probe partition at every level from k − 1 LCA levels, the census at every level from one own-node pair table summed over DFS ranges (A2).
\item[{tried, and what it gave}] The slide's counts: \textbf{proved}, 0 of 26,454 differ from the tree (before the split). The overlap and census on 6bba\_\allowbreak{}48816121 and 44b6\_\allowbreak{}0113de3b: \textbf{measured} (numbers in the cell table, row 1, since they are graded by a program's key).
\item[{status}] built, error-wired
\item[{next}] The LCA and DFS-range forms (A2), which remove the bisection and the per-level recount.
\end{description}

\tablerecord{\textbf{M4. The moments}}
\begin{description}
\item[{the algebra it holds to}] Per component, exact: the mass, the three index sums, the six second moments. Together these are the harmonics up to ℓ = 2 (ℓ = 0 the mass, ℓ = 1 the sums, ℓ = 2 the moments). Every node's moments are its subtree's own-node moments summed (A3).
\item[{does today}] \texttt{max\_\allowbreak{}tree.\allowbreak{}cu} accumulates them on the device for the components at the cut; \texttt{grow\_\allowbreak{}leaves} carries them out (\texttt{ENGINE\_\allowbreak{}BODY\_\allowbreak{}WORDS} 11: peak, mass, 3 sums, 6 moments).
\item[{wants}] Every node's moments at every level, by a prefix over the DFS order (A3).
\item[{tried, and what it gave}] Carried into every record program and the viewer; proofs are the programs' (device equal to host).
\item[{status}] built, error-wired
\item[{next}] A3's prefix, which a level-free program (the slide's chains) needs.
\end{description}

\tablerecord{\textbf{M5. The correlation operator C}}
\begin{description}
\item[{the algebra it holds to}] \texttt{C\_\allowbreak{}\{t,t′\}[a, b](v)} = Σ\_\allowbreak{}x [L\_\allowbreak{}t(x) = a] · [P\_\allowbreak{}\{t′\}(x + v)] · [L\_\allowbreak{}\{t′\}(x + v) = b]: one count of labeled positive voxels under a shift, read by different reductions (A4).
\item[{does today}] \texttt{shift\_\allowbreak{}agreement}: the argmax of the unlabeled sum over every v, by an NTT mod 998,244,353, exact because every count is below the prime. \texttt{climb\_\allowbreak{}machine}: a local argmax per label from a start lag, 26 neighbours a step; \texttt{-\allowbreak{}-\allowbreak{}box} reads C at the 27 shifts around every climbed lag; \texttt{-\allowbreak{}-\allowbreak{}core} the nested widths; extents per label. \texttt{body\_\allowbreak{}overlap}: C at one lag.
\item[{wants}] One C per frame pair over a window of shifts, then only reads: sums, local argmax, transpose, box sums by one summed-area table (A4). The window is the one open design choice, and it must come from the request, not a scale. The search read against the integrated noise floor (M8, Doug, 23 September): a count at a shift is weighed by how far it departs from what the noise functionals predict there, so the climb and the null draws spend work only where the field departs from the floor.
\item[{tried, and what it gave}] The box (\texttt{-\allowbreak{}-\allowbreak{}box}, 23 September): \textbf{proved}. C is read at the 27 shifts around every climbed lag and at the drift, for every forward pair and every null draw. A cell meeting more labels than the kernel's table writes its raw pieces and the host merges them, with no cap. At the climbed lag, C sums to the climb's held count on every climber; the climbed lag is the highest of its 27 neighbours on every climber; at the drift, C equals \texttt{body\_\allowbreak{}overlap} on every label (numbers in the cell book's on\_\allowbreak{}the\_\allowbreak{}engine.md, since they were taken on a program's sample). The core per width: \textbf{proved} at width 0 against C at the climbed lag, 35,755 of 35,755 (numbers in the cell table, row 13).
\item[{status}] built; the box proved
\item[{next}] Make C the one pass, with the window from the request.
\end{description}

\tablerecord{\textbf{M6. The exact marginal}}
\begin{description}
\item[{the algebra it holds to}] Z = the sum over injective arrangements of sources onto targets (or a null each) of the product of their integer weights: a permanent. It factors exactly over the connected components of the bipartite support graph, so a source's probability over its component is its probability over everything (A5).
\item[{does today}] \texttt{marginal}: Z and Z\_\allowbreak{}\{b→·\} for each source over a local set, device equal to the host, compared by cross-multiplying. The weights are the caller's.
\item[{wants}] The set taken as the source's whole support component, so the probability is the global one exactly; the component's size reported, not truncated.
\item[{tried, and what it gave}] Timings and outcomes are the program's (cell table, row 7).
\item[{status}] built
\item[{next}] The component as the set (A5); the limits below become measured costs.
\end{description}

\tablerecord{\textbf{M7. One-to-one matching}}
\begin{description}
\item[{the algebra it holds to}] The heaviest one-to-one set on exact integer costs.
\item[{does today}] \texttt{heaviest\_\allowbreak{}matching}.
\item[{wants}] none
\item[{tried, and what it gave}] Used by the cell program's \texttt{-\allowbreak{}-\allowbreak{}print-\allowbreak{}match} (cell table, row 7).
\item[{status}] built
\item[{next}] none
\end{description}

\tablerecord{\textbf{M8. The entropy history}}
\begin{description}
\item[{the algebra it holds to}] Per voxel and bit, the flips in each window of transitions; the parity of a bit's total flips equals that bit of the net change, at every voxel (A7). The window-overlap cloud folds into the same pass.
\item[{does today}] \texttt{entropy\_\allowbreak{}history}: the flips, the parity proof, the cloud, one pass on the device, error-wired (23 September; broken parity is a logic error). A sample of fewer than two frames is a request error (no transition to count).
\item[{wants}] The window from the request or the data, not the machine (see the audit). The noise categories as functionals (Doug, 23 September; the plan in compression\_\allowbreak{}table.md): fixed pattern from the bits that never flip, shot and read from the photon transfer curve (Σ (I\_\allowbreak{}t − I\_\allowbreak{}\{t−1\})² against Σ I\_\allowbreak{}t per signal level, exact sums), quantization as one step of the lane, and the spatially correlated terms per row, column and place. Integrated, they give every voxel's floor and so every departure from it.
\item[{tried, and what it gave}] \textbf{proved} exact on 2,000 voxels × 9 windows. Measurements of what it shows are the cell program's (cell table, row 12). Its identity by spatial null permutation (\texttt{knf\_\allowbreak{}identity}, M19; A13, the two crystals): \textbf{proved}, the host's walk equals the engine's history at all 262,144 voxels of a 64³ lattice over 16 windows, and of 8,192 single flipped bits 1,500 leave it unchanged, each where the window rule says, so the history is many-to-one; the 48 exact cube motions carry it and keep its cloud.
\item[{status}] built, error-wired
\item[{next}] Take the window as a request field. The pairwise test between sections: one body's departure curve against another's (A13, the two crystals).
\end{description}

\tablerecord{\textbf{M9. The files}}
\begin{description}
\item[{the algebra it holds to}] Every file the machine writes is proved by reading it back: pixels voxel for voxel against a second read of the source, and every part by its seal (M13; the crystal's since 23 September, the others' CRC-64 until they are sealed too). Integer lifting is exactly invertible (A8).
\item[{does today}] \textbf{\texttt{.\allowbreak{}kcr}}, the crystal (Doug named it on 23 September, the Kolmogorov crystal; \texttt{.\allowbreak{}iapx} crystals are not read, sets are re-ingested). Layout, read front to back: head (12 words), the seal (6 roots, every lane node, every chunk leaf), offsets and stream, the deflated side bytes, the member tables and names, EOF. \textbf{dicom} is a source kind: one sample is one SeriesUID, slices ordered by an exact key from their positions (M1), signed pixels sign-extended and lifted by 2\textasciicircum{}15 into the u16 lane (a slice whose raw words differ keeps its pixel bytes), every header kept whole in the side bytes (deflated by our own LZ77 and Huffman coder). The tower (\texttt{tower\_\allowbreak{}lift}, integer lifting to the floors), \texttt{compression} (a Rice coder per block), \texttt{apxrep} (the container), and the rebuild proof on every ingest. \texttt{.\allowbreak{}oapx}: the entropy history. \texttt{.\allowbreak{}bapx}: the bodies' table. \texttt{flatten}: the set as one magnitude vector per label in one file. The ingest readers: zarr v2/v3/N5 with every codec, TIFF, HDF5, npy/npz, NRRD, NIfTI, \texttt{.\allowbreak{}stack}. \textbf{zip} (23 September): its own module (\texttt{engine/\allowbreak{}base/\allowbreak{}zip}), the ZIP64 walk moved out of \texttt{npy.\allowbreak{}c}, the directory read whole with no size cap, members read and CRC-32 checked, an index by containing folder, and one archive held across samples. \texttt{npy} reads through it. Every part is error-wired: an IO failure carries errno, bad contents in the machine's own files are logic errors, a malformed source is a request error.
\item[{wants}] The export (Doug): the crystal written back out as any source format (the dicom zip byte for byte from the side bytes, npy, NIfTI, TIFF, zarr). Names for the other files (flattened, history, bodies, key) are open with Doug. The floor (Doug, 23 September): the residue coded against the integrated noise functionals (M8) over exact counts. Exact arithmetic coding ends within 2 bits a stream of the length those functionals give, and enumerative coding in the engine's exact integers hits it exactly; a fixed-width range coder rounds each interval and is only the fast approximation. That length is the most any coder can take once the functionals are the real categories, with the model's own bits counted as a two-part code (F5 and F3′ in compression\_\allowbreak{}table.md). It replaces the Rice coder's chosen block, k and escape.
\item[{tried, and what it gave}] \texttt{.\allowbreak{}kcr} on the RSNA knee: \textbf{proved}, 1 × 34 × 960 × 960 unsigned, 25,218,496 bytes (40.2\% of raw) sealed \texttt{7ac2cb89…12eb2a}; 1 × 24 × 640 × 640 signed, 11,319,416 bytes (57.5\%) sealed \texttt{3e70b8cb…8d9966}; rebuilt voxel for voxel, pixel for pixel and node for node. The RSNA session's test\_\allowbreak{}series (23 September): 15 of 15 held, 192,020,272 of 599,191,552 bytes (32.0\%), set root \texttt{98b25a42…d7af979}, prove 696 ms with every seal node held. Deflate: within about 1\% of zlib-9 on DICOM headers, round-tripped through our inflate and zlib. Before the seal, \texttt{.\allowbreak{}iapx}: \textbf{proved}, 25 of 25 samples, 42.0\% of raw, set CRC \texttt{091daa41e1aceb7e}; iapx-prove on 44b6\_\allowbreak{}0113de3b after the tower and compression wiring: 1 of 1, CRC-64 \texttt{faa6d04b0c716352} over the set. A missing sample reports module 10, site 178, status 2 (errno ENOENT). \texttt{flatten}: \textbf{measured}, slowed by its residual slab sized to half the free device memory, with the same bytes out one frame at a time (numbers in the cell book's on\_\allowbreak{}the\_\allowbreak{}engine.md, since they were taken on a program's sample). Codec variants: the residue table in the compression book's ledger.
\item[{status}] proved: the field, and the RSNA knee as \texttt{.\allowbreak{}kcr}
\item[{next}] The floor: F2 per floor on 44b6\_\allowbreak{}0113de3b against the Rice coder, then the 25 re-ingested as \texttt{.\allowbreak{}kcr} (compression\_\allowbreak{}table.md). The prose and identifiers still say iapx (\texttt{engine\_\allowbreak{}iapx\_\allowbreak{}head}, \texttt{entry\_\allowbreak{}iapx\_\allowbreak{}encode}, the cell README) where the file is \texttt{.\allowbreak{}kcr} (\texttt{ENTRY\_\allowbreak{}CRYSTAL\_\allowbreak{}SUFFIX}, engine.cu:127). The export. The engine's opens need a long-path form (a set path past 260 characters fails today). Flatten: per-request orders per sample (ruled), and a slab not sized as a fraction of memory.
\end{description}

\tablerecord{\textbf{M10. The record machine}}
\begin{description}
\item[{the algebra it holds to}] A record program is steps over limbs, run for every record in one sweep; truthy verdicts are fields (zero false, nonzero true; AND a product, OR a sum). A step may read a lookup table: a general one-variable function of one register's low bits, so the machine's edges are no longer only linear (Doug, 23 September; A13). The step count is the scheduler's n, held apart from the register file, which is the binding resource (Doug, 23 September; A13).
\item[{does today}] \texttt{keymath} imprints record programs, \texttt{key\_\allowbreak{}schedule} lays them out, \texttt{cycle} runs them on the device (with a host reference); every entry is error-wired. The operations are field, constant, product, sum, difference, ladder, absolute, compare, signed field, and table (\texttt{ENGINE\_\allowbreak{}RECORD\_\allowbreak{}TABLE}, op 10, 23 September), threaded through the config, \texttt{keymath}, \texttt{key\_\allowbreak{}schedule} and \texttt{cycle}, on the device and the host port. A table maps the low \texttt{index\_\allowbreak{}bits} (at most 32) of one register to an \texttt{out\_\allowbreak{}bits} value, laid out as 2\textasciicircum{}index\_\allowbreak{}bits rows of ⌈out\_\allowbreak{}bits/32⌉ limbs. The exact integer's division joined as four operations on 23 September (M1): quotient (\texttt{ENGINE\_\allowbreak{}RECORD\_\allowbreak{}QUOTIENT}, op 11), remainder (12), gcd (13) and exact quotient (14). The quotient rounds toward zero and the remainder carries the numerator's sign. The gcd is never negative. A zero divisor, or an exact quotient that leaves a remainder, refuses the lane. The widths are read at imprint: a quotient is as wide as its numerator, a remainder as the narrower operand and a gcd as the wider. On the device, Knuth's algorithm D divides, the gcd is Euclid's on it and the exact quotient is the multiply and mask. Only a program that divides carries the scratch, 5 · WIDE + 4 limbs a lane. The host's steps are the exact integer library itself. Xor (\texttt{ENGINE\_\allowbreak{}RECORD\_\allowbreak{}XOR}, op 15), and (16) and the two's complement wrap (17, to 4 bits or more) joined on 24 September, bitwise on each operand's two's complement sign-extended without end. The step count has no bound since 24 September: each sign sits beside its register, so \texttt{ENGINE\_\allowbreak{}RECORD\_\allowbreak{}LIMBS\_\allowbreak{}MOST} = 256 limbs of registers is the only bound, and floors of steps stack into one program. Register reuse (\texttt{reuse} on the key-schedule and record requests, off by default) frees a register once its last reader has run, so a chain needs only its live registers; off, every step keeps its own register as before.
\item[{wants}] Limits measured, not written in (see the audit). The table's functions chosen without fitting (item 3 of \hyperref[chap:kolmogorov-arnold]{kolmogorov\_\allowbreak{}arnold.md}).
\item[{tried, and what it gave}] Device equal to host on every program run (the programs' proofs). The table (\texttt{test/\allowbreak{}record\_\allowbreak{}table\_\allowbreak{}test.\allowbreak{}cu}, run by \texttt{test/\allowbreak{}record\_\allowbreak{}table\_\allowbreak{}test.\allowbreak{}sh}): \textbf{proved}, 15 checks, 0 failed, every table filled by running an ordinary-ops program over the whole 16-bit alphabet, so the ops are the oracle. x² (product) and |x − 30000| + 100 (constant, difference, absolute, sum): the table equals the ops on 4,096 lanes, device, host and table alike. Composition, f(x) = 65535 − x and g(y) = |y − 30000|: a table read through a table equals the ops for |35535 − x|, a single table composed on the host (g read through f) equals the two, and the reversed order f(g(x)) differs on some lane, so the chain keeps its order. Register reuse: a 21-step additive chain gives the same output from 3 limbs as from 21. Two refusals: an index wider than its source register, and a table index with no table behind it. The division (\texttt{test/\allowbreak{}record\_\allowbreak{}divide\_\allowbreak{}test}, \texttt{build/\allowbreak{}20260923\_\allowbreak{}234739\_\allowbreak{}record\_\allowbreak{}divide\_\allowbreak{}test}): \textbf{proved}, 15 checks, 0 failed. The device equals the host word for word on 4,096 lanes of signed 160-bit numerators in the 64-limb register file and on 512 lanes of 2,048-bit numerators in the 256-limb file. On every lane numerator = quotient · divisor + remainder, the remainder is below the divisor and signed as the numerator, the gcd divides both, and the exact quotient returns the factor it was built from. A zero divisor and an inexact division refuse the lane, host and device. 3\textasciicircum{}40 divides exactly by 3\textasciicircum{}20 to 3\textasciicircum{}20. A quotient reading a later step is refused at imprint.
\item[{status}] built, error-wired; the table, register reuse and the division proved (A13)
\item[{next}] \texttt{cycle\_\allowbreak{}record\_\allowbreak{}run\_\allowbreak{}host} (cycle.c) is not error-wired yet.
\end{description}

\tablerecord{\textbf{M11. The golden ladder}}
\begin{description}
\item[{the algebra it holds to}] β(n): n's rung among the Fibonacci numbers that fit the word.
\item[{does today}] \texttt{golden\_\allowbreak{}bands}, \texttt{band\_\allowbreak{}of}, 92 rungs (every Fibonacci number below 2⁶⁴: a word's width, not a scale).
\item[{wants}] none
\item[{tried, and what it gave}] none
\item[{status}] built
\item[{next}] none
\end{description}

\tablerecord{\textbf{M12. Errors and integrity}}
\begin{description}
\item[{the algebra it holds to}] The only acceptable error is a stateful one (Doug): request, resource (a failed CUDA call, allocation or read, with its status or errno) and logic (unforgivable; the result is withheld). First write wins; each request carries its record and the engine keeps its own. The record holds kind, module, site, status, execaddr, evacaddr, imagebase and up to 16 return frames, and the builds emit a \texttt{.\allowbreak{}pdb}, so execaddr − imagebase steps to the line.
\item[{does today}] Wired, 23 September: the engine's own entries, max\_\allowbreak{}tree, flatten, decimal\_\allowbreak{}double, unit\_\allowbreak{}sweep, cycle, keymath, key\_\allowbreak{}schedule, residual, grow, apxrep, compression, tower, entropy\_\allowbreak{}history, zip, npy, dicom, obsignatio (modules 0 to 17). Output (Doug, 23 September): no print in any loop; each functional fills its result struct and one print format renders it on demand through \texttt{scriptura} (our own port of MMgr's formatting and SWAR memory scans, its memory operations in M17); errors always print. The run log's default place is beside the program (\texttt{engine\_\allowbreak{}program\_\allowbreak{}directory}, 23 September): \texttt{build/\allowbreak{}logs/\allowbreak{}track\_\allowbreak{}driver.\allowbreak{}log} for the published driver, whatever folder it runs from; \texttt{-\allowbreak{}-\allowbreak{}log} still names any other. A failed ingest names its stage: no place, source unread, crystal unwritten, or not rebuilt.
\item[{wants}] Every module wired: stack, schedule, golden\_\allowbreak{}bands, the other ingest readers and codecs, run\_\allowbreak{}cfg, run\_\allowbreak{}log, cfg\_\allowbreak{}json, relate\_\allowbreak{}frames, double\_\allowbreak{}fields, \texttt{cycle\_\allowbreak{}record\_\allowbreak{}run\_\allowbreak{}host}. Integrity is M13; device scheduling is M14.
\item[{tried, and what it gave}] Every wiring build checked against a program's run, 0 lanes differ in both modes (numbers in the cell book's on\_\allowbreak{}the\_\allowbreak{}engine.md). Fixed 23 September (Doug: “you should fix any defect found”): a series the archive does not hold is refused as no such member (the entry layer asks the archive; a folder search that finds nothing no longer counts as found); a failed ingest removes the directories it made (the maker counts what it created; its 1,024-byte walk bound is gone); a failed prove prints no set root; a structurally refused crystal prints as not read; every path buffer is sized by the platform's own limit (\texttt{ENGINE\_\allowbreak{}PATH\_\allowbreak{}ROOM}: 32,768 on Windows with a long-path-aware manifest, \texttt{PATH\_\allowbreak{}MAX} elsewhere), and seven sites that truncated a long path silently now refuse it; the npy and nrrd header caps are gone; the host exact integer is fitted to the record path's widest step (256 limbs), so host and device accept the same steps. A 395-character crystal path ingested and proved to the same root.
\item[{status}] wiring in progress
\item[{next}] The next unwired module.
\end{description}

\tablerecord{\textbf{M13. The seal}}
\begin{description}
\item[{the algebra it holds to}] Every file, set and state is sealed by a dimensional Merkle DAG of keyed BLAKE3 nodes, one derived key per level, 32 bytes a node, every node stored. A change anywhere changes exactly its path of nodes up to the root, so a mismatch names its place (A10; the theory is \hyperref[chap:obsignatio-seal]{obsignatio\_\allowbreak{}seal.md}).
\item[{does today}] \texttt{obsignatio} (module 17), our own BLAKE3, host and device sharing one compression core: \texttt{obsignatio\_\allowbreak{}signum} (any length, extended output), \texttt{obsignatio\_\allowbreak{}many} (equal messages at a stride, level-parallel), \texttt{obsignatio\_\allowbreak{}bits} (bit ranges of a limb stream), \texttt{obsignatio\_\allowbreak{}lanes} (rows → planes → volumes → lanes), and the level keys. The \texttt{.\allowbreak{}kcr} stores its six roots, every lane node and every chunk leaf. Ingest seals from the source lanes and proves the file against them. Prove and load check in order: stored bytes, then the inflated side, then the decoded rows, then the root. The set root seals every sample's root in the order named. The tower's CRC is gone. The seal sits after the head (Doug, 23 September): both placements measured alike on the knee (the same bytes, prove median 76 ms each), and the seal exists before the stream is written. Ingest witnesses twice (Doug: keep both): the decoded voxels against the source on the device, then a second, independent read of the source pixel for pixel, so a flip in the first read is not sealed as the truth.
\item[{wants}] Stage B: the history, bodies, flattened, key and schedule files sealed the same way. The universal root over every file the engine touches plus the engine itself (its hash injected at compile time), and the anchor that makes it tamper-proof and not only accident-proof (a signature or a published root).
\item[{tried, and what it gave}] The suite (\texttt{obsignatio/\allowbreak{}test}, the official BLAKE3 vectors): \textbf{proved}, 2,182,450 cases, 0 failures (rerun with the 14 levels, 23 September). That includes every length 0 to 102,400 host against device, every single-bit flip of every vector input, lane locality 9,234 of 9,234 and bit locality 57,216 of 57,216. Five planted bugs (a dropped byte, a wrong merge, a lost carry, no root flag, a stuck counter) each fail it. On a sealed knee crystal: 8 single-bit tamperings each refused and placed (head shape, a root, row t 0 z 8 y 0, a plane node, chunk 1,200, chunk 1,234 from a stream limb, the deflated side, a member name), the clean copy proved. Cost: 5.4\% of the crystal, 2.1\% of raw at 960 columns (16/X of raw from the rows).
\item[{status}] built and proved: the crystal and the set
\item[{next}] Stage B, then the universal root.
\end{description}

\tablerecord{\textbf{M14. The scheduler}}
\begin{description}
\item[{the algebra it holds to}] Jobs from separate processes share one device by memory: admit while the declaration fits the measured headroom, grow and warn when a job outgrows its reservation, keep each request signum's peak (A11; the theory is \hyperref[chap:tessera-scheduler]{tessera\_\allowbreak{}scheduler.md}).
\item[{does today}] Nothing arbitrates between processes; schedule.cu plans one process against free/3·2, a chosen share.
\item[{wants}] \texttt{tessera} (Doug, 23 September): a daemon per host and device over a local socket (named pipe, Unix socket, systemd socket activation, a docker socket mount); per-pid measurement (NVML on Linux, the PDH GPU Process Memory counter under WDDM, where NVML reports none); dead sweeps and lost and found; over budget asked at admission against the signum's last peak, held, overridable; the budget and every time a request field; the driver submitting through it; service files; a test suite.
\item[{tried, and what it gave}] Measured on this machine: the PDH counter lists each process's dedicated device bytes by pid. NVML per-process memory is unavailable under WDDM by NVIDIA's documentation. The ledger (\texttt{tessera\_\allowbreak{}ledger}: the accounting, backfill with the head kept, one min-heap of deadlines): \textbf{proved}, 143,808 cases, 0 failed, rerun from \texttt{engine/\allowbreak{}daemon/\allowbreak{}test} (the module moved from \texttt{engine/\allowbreak{}base/\allowbreak{}tessera}, 24 September). The frame (\texttt{tessera\_\allowbreak{}frame.\allowbreak{}c}, 128 little-endian bytes): \textbf{proved}, 100,000 random round trips byte for byte, the layout at fixed places, and 131 single-bit flips of the head, each refused or held exactly as the rule says. The measure (\texttt{tessera\_\allowbreak{}measure.\allowbreak{}c}: NVML loaded at run time, the PDH counter matched on pid and LUID, a process held by a handle or pidfd): \textbf{measured} on the RTX 3070 under WDDM, 8 of 8 cases. This process read 24,576 bytes, then 412,254,208 after allocating 268,435,456 (the rest is its CUDA context); the device's in-use rose by the same 412,070,912; a pid with no device work read 0; a child held before its exit read live, then ended. NVML's total is 655,360 bytes above what the runtime reports (WDDM keeps them), and C is NVML's figure. The client (\texttt{tessera\_\allowbreak{}client.\allowbreak{}c}: connect, spawn the daemon if absent, submit, override, wait, precalc kept, release) and the daemon (\texttt{tessera\_\allowbreak{}daemon.\allowbreak{}c}, linked with the ledger, the measure, the paths, the frame, scriptura and obsignatio) compile clean at /W4. The job test (\texttt{tessera\_\allowbreak{}job\_\allowbreak{}test}, the client against the built daemon, which the client starts through \texttt{daemon\_\allowbreak{}path}): \textbf{proved} on the RTX 3070 under WDDM, 24 checks, 0 failed, on two runs in a row (\texttt{build/\allowbreak{}20260924\_\allowbreak{}014618\_\allowbreak{}tessera\_\allowbreak{}test} and \texttt{build/\allowbreak{}20260924\_\allowbreak{}014737\_\allowbreak{}tessera\_\allowbreak{}test}, with the ledger, frame and measure tests, each exit 0; before the seal, 15 checks, 0 failed, \texttt{build/\allowbreak{}20260924\_\allowbreak{}013052\_\allowbreak{}tessera\_\allowbreak{}test}). Its first four paths, in order. A new signum declaring 64 MiB is admitted with 64 MiB granted; it takes 256 MiB, and its release reports it grew to 412,254,208 bytes, with a peak of 412,254,208, the process's whole dedicated memory, its CUDA context included. The same signum declaring four times that peak is asked, the ask naming the kept peak, and the override admits it on its declaration. Declaring exactly the kept peak is admitted at once. Declaring four times the peak with no answer is held for its 0.4 s holding time, then lost; its ticket names a lost and found directory, and the precalc kept releases it. The fifth part runs in a scratch \texttt{TESSERA\_\allowbreak{}STATE} inside the run's build directory. The lost ticket's seal holds, with “precalc kept” as its body's last line. The daemon ends once idle, and the history it leaves is whole records and the seal (80 bytes here: one record and the 32-byte seal). A byte flipped, a byte cut and the seal stripped are each refused: the submit fails, and no daemon starts. The history restored, the same signum at four times the peak is asked with the kept peak, which proves the history was read. The daemon seals the history and the tickets with \texttt{obsignatio\_\allowbreak{}seal} (keyed BLAKE3, on the host), and it makes no CUDA call, so it holds no device context (\hyperref[chap:tessera-scheduler]{tessera\_\allowbreak{}scheduler.md}). The daemon the last check starts ended itself within 5 s of the test. The history is loaded before any endpoint exists, so a refused history ends the daemon before a client can reach it (24/0 again, \texttt{build/\allowbreak{}20260924\_\allowbreak{}015203\_\allowbreak{}tessera\_\allowbreak{}test}). A program's driver submits through it (24 September; its jobs and the sample they ran on are in the cell book's on\_\allowbreak{}the\_\allowbreak{}engine.md). \textbf{Measured} there: two proves of one crystal in a row, each granted its declaration of 838,860,800 bytes, peaked at 3,958,566,912 and 3,962,761,216. Anchor\_\allowbreak{}sift's run on the real state gave the same declarations, and the two prove peaks in the other order: identical requests measured 4,194,304 bytes (2\textasciicircum{}22) apart. The second prove was admitted on its declaration, below its kept peak, as A11 says, so until its sweeps grew it, its reservation stood 3,119,706,112 bytes below what its signum is known to need, and another job could have been admitted into that room. Doug's rule stands (24 September): “We reserve what they ask for and then if it cost less we remember that for next time”. A job that takes more than it declared is grown and warned by the sweeps, as A11 says. The same day, Anchor\_\allowbreak{}sift changed the working tree's code, uncommitted and not approved by Doug, to reserve the larger of the declaration and the kept peak (\texttt{tessera\_\allowbreak{}ledger\_\allowbreak{}wants}). Its ledger suite passed (\texttt{build/\allowbreak{}20260924\_\allowbreak{}025121\_\allowbreak{}tessera\_\allowbreak{}test}, budget 22 checks, 0 failed), and with it two proves each reserved 3,962,761,216 bytes, the kept peak. The sims submit too (\texttt{engine/\allowbreak{}sims/\allowbreak{}sim\_\allowbreak{}job.\allowbreak{}cu}, 24 September): each device sim is one job, with a signum over its name and arguments, the bytes of its first allocation declared before it makes it, the same three times, and \texttt{TESSERA\_\allowbreak{}OVERRIDE=1} to override. \textbf{Measured} here in a scratch \texttt{TESSERA\_\allowbreak{}STATE}: all six device sims and Ω's engine run were admitted and released, each peak above its declaration (the figures are in M19), and the history was then 368 bytes, seven records and the seal. Two peaks differ from Anchor\_\allowbreak{}sift's runs by 2,097,152 bytes (2\textasciicircum{}21): nbody\_\allowbreak{}lattice read 175,276,032 here against 173,178,880, and fixed\_\allowbreak{}pattern 145,915,904 against 143,818,752. No program in the repo loads the engine DLL, whose only builder is \texttt{build\_\allowbreak{}engine.\allowbreak{}sh}, so no other caller is left to submit. \textbf{Reported by Anchor\_\allowbreak{}sift, not rerun here:} the Linux build and suite under WSL 2 (gcc 13.3.0, CUDA 13.3), with 0 warnings after fixes to \texttt{PATH\_\allowbreak{}MAX} under strict C11, the noinline helpers under gcc, \texttt{\_\allowbreak{}GNU\_\allowbreak{}SOURCE}, the timer thread's missing return and the Windows-only strings. Under WSL, NVML read a process as 0 bytes before and after it allocated 256 MiB, so a pid can't be measured there. The measure refuses to open on a paravirtual device (\texttt{tessera\_\allowbreak{}measure\_\allowbreak{}paravirtual}), and the daemon refuses to run. The measure test and the job test both check exactly that refusal. The systemd units are in \texttt{engine/\allowbreak{}daemon/\allowbreak{}service} (\texttt{tessera@.\allowbreak{}socket} and \texttt{tessera@.\allowbreak{}service}, which this session read). Under WSL, socket activation started the daemon on a client's connection, and the measure refusal ended it. The docker socket mount, also reported: a static client with no CUDA (\texttt{engine/\allowbreak{}daemon/\allowbreak{}test/\allowbreak{}tessera\_\allowbreak{}socket\_\allowbreak{}probe.\allowbreak{}c}, which this session read) ran in a container under Docker Engine 29.1.3 inside the WSL distro, the host's socket bind-mounted and named by \texttt{TESSERA\_\allowbreak{}RUNTIME}. Its submit reached the socket-activated daemon on the host, and the daemon's WSL refusal ended it, so the submit was refused. How to submit a job and start the daemon: engine/daemon/README.md.
\item[{status}] ledger, frame, measure and the daemon end to end proved on Windows, the history and the tickets sealed; a program's driver and the device sims submit every run, measured; the working tree reserves the larger of the declaration and the kept peak, against Doug's rule; Linux built under WSL, socket activation and the docker mount reached, all reported
\item[{next}] Native Linux, where NVML measures a pid, a daemon kept running on the activated socket and a job admitted through the docker mount (this machine has no native NVIDIA Linux driver); WSL's jobs, which no daemon there can measure; the unsealed history kept aside as \texttt{history.\allowbreak{}unsealed}.
\end{description}

\tablerecord{\textbf{M15. The sift}}
\begin{description}
\item[{the algebra it holds to}] A pattern of p points is found in a lattice by probing anchors, then verifying the whole pattern at the alignments that pass. What the anchors filter is set before the search by the corpus's collision entropy H2 and its coherence (a period P found against its own multiples): k anchors on one residue class collapse to one probe, and a periodic corpus cannot be filtered below 1/P by anchors. The dispatch between the in-order and free-order arms is one exact comparison, 100·total² against 85·distinct·Σcount², which needs 143 bits and so at least 8 limbs.
\item[{does today}] \texttt{nbody/\allowbreak{}anchor\_\allowbreak{}sift/\allowbreak{}} is a byte-identical copy of anchor\_\allowbreak{}sift's \texttt{src/\allowbreak{}engine/\allowbreak{}c/\allowbreak{}engine}: the search, the steering, the portable scan and one scan arm per instruction set. \texttt{anchor\_\allowbreak{}sift\_\allowbreak{}anchors\_\allowbreak{}for} returns 1 where a period is found and 4 otherwise. The driver's build does not compile it: \texttt{build\_\allowbreak{}engine.\allowbreak{}sh} takes only the exact integer, and since 23 September it takes that from this tree's \texttt{base/\allowbreak{}no\_\allowbreak{}rounding} (\texttt{ANCHOR\_\allowbreak{}EXACT\_\allowbreak{}ROOT} still overrides). \texttt{engine/\allowbreak{}CMake\allowbreak{}Lists.\allowbreak{}txt} does (23 September, its paths fixed to this tree): the sift, \texttt{base/\allowbreak{}no\_\allowbreak{}rounding} with every arm, \texttt{render} and every root bench, with the CUDA arms where a CUDA compiler is found.
\item[{wants}] One source for the sift and the exact integer, taken from anchor\_\allowbreak{}sift by a pinned git dependency, not a copy (anchor\_\allowbreak{}sift's port review, 23 September). An anchor count between 1 and 4 on a partially coherent corpus. The readers (M9) feed it real lattices of 1 to 8 dims.
\item[{tried, and what it gave}] anchor\_\allowbreak{}sift's benches, reported in engine/README.md: soundness, 3,421 rows and 465,546 planted occurrences, none refused; the cost per alignment flat over 4 KB to 4 MB; periodic16 off the histogram by exactly 4096 = 2\textasciicircum{}((k−1)·H2), found before the search by coherence; dispatch on flatness alone, 41 of 42 fastest. Built and run here (23 September) on four configurations: MinGW gcc (host), MSVC with CUDA 13.3, WSL gcc 13 with CUDA, WSL clang 18 with CUDA, all at sm\_\allowbreak{}86 and 256 limbs. Every bench exited 0 on every configuration (bench\_\allowbreak{}lattice is not built under MSVC, which has no C99 \texttt{\_\allowbreak{}Complex}). \texttt{bench\_\allowbreak{}exact}, \texttt{bench\_\allowbreak{}lattice} (3,253 rows, every one \texttt{hold}), \texttt{bench\_\allowbreak{}scaling\_\allowbreak{}reads} and \texttt{bench\_\allowbreak{}coherence} are byte-identical across every configuration that builds them. \texttt{bench\_\allowbreak{}exact\_\allowbreak{}arms} at 4,096: every arm agrees with portable, AVX2 1.11× to 1.31×, CUDA 0.77× to 1.83×; the planted period read 4,092 against arithmetic's 4,092, and the device itself answered 2,046 against portable's 2,046. \texttt{bench\_\allowbreak{}steer\_\allowbreak{}arms}: 0 checks failed on each (\textbf{proved}, per configuration).
\item[{status}] anchor\_\allowbreak{}sift's; copied; built and run here on four configurations
\item[{next}] Doug's rulings on the port (canonical copy, n\_\allowbreak{}body or nbody, the error model).
\end{description}

\tablerecord{\textbf{M16. Render}}
\begin{description}
\item[{the algebra it holds to}] A step's state drawn as a sheet or a volume straight from the state, with no export between. The host arm in C and the device arm in CUDA make the same bytes.
\item[{does today}] \texttt{render/\allowbreak{}} is anchor\_\allowbreak{}sift's renderer with its comments stripped: \texttt{anchor\_\allowbreak{}raster\_\allowbreak{}render} (a sheet) and \texttt{anchor\_\allowbreak{}volume\_\allowbreak{}render} (a volume), each preferring the device and falling back to the host. The engine's own drawing today is \texttt{view/\allowbreak{}vis\_\allowbreak{}png}.
\item[{wants}] One renderer for the engine's state, taken from anchor\_\allowbreak{}sift by a pinned dependency.
\item[{tried, and what it gave}] anchor\_\allowbreak{}sift's \texttt{bench\_\allowbreak{}raster}: twenty sheet and twenty volume configurations, host against device byte for byte. Here (23 September, moved from \texttt{base/\allowbreak{}render} to \texttt{engine/\allowbreak{}render}): 20 of 20 \texttt{ok} on each of the four configurations, the device arm present and graded against the host on the three CUDA builds, the host graded alone on MinGW (\textbf{proved}, per configuration).
\item[{status}] anchor\_\allowbreak{}sift's; copied; built and run here
\item[{next}] As M15.
\end{description}

\tablerecord{\textbf{M17. Memory operations}}
\begin{description}
\item[{the algebra it holds to}] length, find, copy, move up (target above source, overlapping), compare (−1, 0, 1) and fill, each an exact byte operation; every arm answers exactly what a byte loop answers. Buffers are aligned and terminated by the caller (Doug's rule), so an arm may read an aligned block past the end and mask it.
\item[{does today}] \texttt{scriptura}'s arms (23 September), one file an instruction set, chosen once at first use; below 32 bytes a call goes straight to portable. \texttt{portable}: 64-bit SWAR (MMgr's), overlapping head and tail words to 32 bytes, 32-byte blocks and an overlapping last block above. \texttt{avx2}: four blocks a branch (\texttt{vpminub} for length, OR for find, AND for compare), overlapping tails. \texttt{avx2-\allowbreak{}erms}: the same, with \texttt{rep movsb} and \texttt{rep stosb} for copy and fill from 16 KiB. \texttt{avx512-\allowbreak{}unrun} (masked loads and stores), \texttt{neon-\allowbreak{}unrun}, \texttt{sve-\allowbreak{}unrun} (predicated).
\item[{wants}] Every arm graded on its own hardware (AVX-512 and SVE have none here; NEON has the Pi 5). The small sizes at the C runtime's speed.
\item[{tried, and what it gave}] The grade (\texttt{test/\allowbreak{}scriptura\_\allowbreak{}arm\_\allowbreak{}test.\allowbreak{}c}, every operation over 483 lengths from 0 to 4,160, the terminator, match or difference at four places): portable and avx2 agree with a byte loop on MinGW gcc 16, MSVC /W4, WSL gcc 13 and WSL clang 18, with 0 warnings (\textbf{proved}, per configuration). NEON and SVE build and link on aarch64 gcc 13 and clang 18 and emit \texttt{cmeq}, \texttt{shrn}; \texttt{whilelo}, \texttt{ld1b}, \texttt{brkb}, \texttt{cntp}; AVX-512 emits \texttt{zmm} masked byte compares and moves. Against the C runtime (\texttt{bench/\allowbreak{}bench\_\allowbreak{}scriptura.\allowbreak{}c}, time-stamp ticks, best of 7, aligned buffers, an i7-5960X with ERMS), reading the assembly found two defects: compare computed its next address from the last compare, a serial chain (2 to 3.8× memcmp; after, 0.31 to 0.77× on MinGW), and length and find tested 32 bytes a branch against glibc's 128 (after, 0.25 to 0.31× the MinGW runtime at 1 to 16 KiB, 0.93 to 1.25× glibc from 4 KiB). Copy and fill at 256 KiB and 1 MiB ran 1.13 to 1.40× every runtime; string moves from 16 KiB brought them to 1.00×, and from 4 KiB they lost (1.28×), so 16 KiB is the measured crossover on this part. Portable's short path was a word loop and a byte loop with saved registers; overlapping words and a long path kept out of line left it a leaf with no frame (\textbf{measured}).
\item[{status}] built; the x86 arms proved, the aarch64 arms built and emitting
\item[{next}] The small sizes against glibc (16 to 64 bytes at 1.06× to 1.7×, one call through the arm table); the string crossover measured per part; the NEON arm run on the Pi.
\end{description}

\tablerecord{\textbf{M18. The period reading}}
\begin{description}
\item[{the algebra it holds to}] Per axis of a lattice, the agreement count at every lag 1 to ⌊n/2⌋ over one set of start positions, so every lag shares one denominator. A candidate is a local peak: the agreement at p and at 2p each stands above both neighbour lags, and the height is the smaller rise (Doug, 23 September). It replaces the margin against the lags that are not multiples of p, which read smoothness as a period (superseded). The floor is a drawn null per axis, not chance (Doug, 23 September): for each axis, every line along it is shuffled on its own, so that axis loses its order while every other axis keeps its structure, and each shuffle reads its strongest height. P is the fundamental, the smallest candidate whose height stands strictly above every one of them, or above a set-wide top the program passes (Doug, 23 September). It replaces the strongest peak, which read a multiple of P, and the global shuffle, which read a period on an axis that carries none whenever another axis carries one (both superseded). Every comparison is an exact cross-multiply (A12).
\item[{does today}] \texttt{period\_\allowbreak{}read} and \texttt{period\_\allowbreak{}draw} (engine/base/period, module 19, 23 September; Doug: its own base module, it returns P and never sets an order, and the program decides which order P sets): a device lattice of u16 lanes, rank 1 to 8. The request names \texttt{draws}, the number of shuffles; \texttt{content}, the lattice's content signum, from which the shuffles are seeded (Doug); \texttt{null\_\allowbreak{}top}, one given top per axis for the set-wide band; and optionally \texttt{band} with \texttt{band\_\allowbreak{}room}, at least rank · draws. It takes draws or a given top, never both: no draws without a top, or a top with draws of its own, is a request error (Doug). The caller passes the signum, for example the sample's lane root from obsignatio, because no module reaches another. Each draw shuffles each axis once (\texttt{period\_\allowbreak{}line\_\allowbreak{}shuffle\_\allowbreak{}kernel}): every line along the axis is permuted alone by a keyed Fisher–Yates, keyed by the content, the draw, the axis, the line and the step; the four keys come from the signum's four words mixed with the counter draw · rank + axis. A draw's histogram must equal the original's; a mismatch is a logic error. The global Feistel shuffle it replaces is removed. Per axis it returns P (0 when none), the candidate (P itself, or the strongest peak when none clears), the agreement at the candidate and at twice it with the higher neighbour of each (\texttt{agreement\_\allowbreak{}beside\_\allowbreak{}candidate}, \texttt{agreement\_\allowbreak{}beside\_\allowbreak{}double}), the pairs per lag, the exact \texttt{margin} (the height over the pairs), and the band's \texttt{band\_\allowbreak{}count} (the draws that reached a peak), \texttt{band\_\allowbreak{}bottom} and \texttt{band\_\allowbreak{}top} (both the given top when one is passed). It also returns N, Σc² as \texttt{collisions} (reported, no longer a floor), the band's heights sorted into \texttt{band} when it is given, and optionally every lag's count. \texttt{period\_\allowbreak{}draw(request, draw, heights)} shuffles each axis once along its own lines, keyed by the content, the draw number and the axis, and returns each axis's strongest height (0 over Q where no peak is reached); the program keeps the top per axis across its series and passes the tops as \texttt{null\_\allowbreak{}top}. The counts come from two kernels (a histogram and one agreement row per lag); the choice is made on the host. A volume of 2\textasciicircum{}32 voxels or more is a request error, because N² must fit 64 bits. The engine builds with it; no entry calls it yet.
\item[{wants}] The knee's orders read from it (the knee program: s = b = P, an odd P a request error, P = 0 order 0), with one band for the set: \texttt{period\_\allowbreak{}draw} across its series, the top per axis passed as \texttt{null\_\allowbreak{}top}. Flatten taking per-sample orders from it.
\item[{tried, and what it gave}] The test (\texttt{test/\allowbreak{}period\_\allowbreak{}test.\allowbreak{}cu}, run by \texttt{test/\allowbreak{}period\_\allowbreak{}test.\allowbreak{}sh}): \textbf{proved}, 700 checks, 0 failed. Every lag's device count equals a host loop. P, the candidate and its height equal an independent host reference of the same rule (the smallest peak above the top; the strongest peak when none clears). Every band is sorted, and its ends are its bottom and top. A second read of the same content gives an identical reading and band. \texttt{period\_\allowbreak{}draw} for draws 0 to 7 reproduces \texttt{period\_\allowbreak{}read}'s band exactly, and the band's top given as \texttt{null\_\allowbreak{}top} judges every axis as the band did. Nine refusals each give a request error from module 19: no lanes, rank 0, rank 9, an extent of 0, 2\textasciicircum{}32 voxels, agreement room short, no draws, band room short, a given top with draws of its own. Planted periods still read exactly and clear their bands under the per-axis null: 7 on one axis of 4,096, height 2,048/2,048 against a band top of 39/2,048; on 24 × 48 × 60, 5 in z (height 34,560/34,560, band top 109/34,560), none in y (no draw reached a peak), and 12 in x (height 34,560/34,560, band top 64/34,560). z is 24 deep, not 20, because the peak rule needs 2P + 1 ≤ ⌊n/2⌋ and 20 slices allow P ≤ 4. \textbf{Measured} at 8 draws: a smooth ramp with noise (x + y + 2z plus uniform 0 to 63), 16 volumes of 24 × 64 × 64, a period on 3 of 48 axes (6.3\%); on one such lattice, y's strongest peak, 15, has height 14/49,152 against a band top of 35/49,152 and is refused, and x reads 5, height 24/49,152 against a band top of 16/49,152. Under the margin rule the same kind of lattice read P = 2 or 3 on every axis (the 34 × 960 × 960 timing run). Uniform random u16, 64 volumes of 32 × 64 × 64: a period on 5 of 192 axes (2.6\%), below the 1/(draws + 1) = 1/9 a permutation test allows (A12); the caller sets the false-period rate with draws (99 draws, 1\%). \textbf{Measured} by \texttt{period\_\allowbreak{}power} (M19), a period 5 planted along x of 32 × 64 × 64 under the camera law, 599 checks, 0 failed: at 8 draws x reads 5 on 32 of 32 at 64, 128 and 1,024 e, and on 31 with 1 multiple at 256 e; at 19 draws, on 32 of 32 from 64 e up. With the plant at 1,024 e the unplanted z and y read a period on 5 of 64 at 8 draws (against 64/9) and on 2 of 64 at 19 draws (against 64/20) (A12). One draw of 34 × 960 × 960 took 152 ms (RTX 3070) under the global shuffle, which is why the band is set-wide (A12). Under the per-axis null (\textbf{measured} 23 September, the same lattice and device): \texttt{period\_\allowbreak{}draw} 455.9 ms a call, and 450.1 ms a marginal draw in \texttt{period\_\allowbreak{}read}, 2.96 times the global shuffle's; one band for the set holds the more (A12). Superseded: the strongest peak under the global shuffle (700 checks, 0 failed; 6 of 192 random axes and 2 of 48 smooth-ramp axes; \texttt{period\_\allowbreak{}power} 599 checks, 1 failed on purpose), which read a multiple, 10 or 15, 103 times in 156 at 8 draws from 32 e up, and with the plant at 1,024 e read a period on the unplanted z and y 23 times in 64 at 8 draws (36\%, against 1/9) and 10 in 64 at 19 draws (15.6\%, against 1/20); the margin rule under the drawn null, 502 checks, 0 failed, 21 of 192 random axes (10.9\%); the chance floor Σc²/N² before it, 431 checks, 0 failed, 63 of 192 (33\%), since Σc² is unchanged by a rearrangement and cannot tell structure from one.
\item[{status}] built, proved against its reference; the local-peak rule, the fundamental above a per-axis drawn null or a set-wide top
\item[{next}] The knee's set-wide band from \texttt{period\_\allowbreak{}draw}; per-sample orders in flatten; the knee pulls it.
\end{description}

\tablerecord{\textbf{M19. The sims}}
\begin{description}
\item[{the algebra it holds to}] A sim builds a lattice whose answer is known and grades a measurement against it, on the device, in exact integers. The camera law (\texttt{sim\_\allowbreak{}camera.\allowbreak{}h}, the same in every sim): electrons S = the background + a ramp · x + an optional planted wave by phase + the brightness of every body whose ellipsoid holds the voxel, each body an integer centre, velocity and semi-axes with a born and an ended frame and a parent. Shot noise is Binomial(4S, ½) − S, which has Poisson's mean and variance exactly and is not Poisson. Read noise is Binomial(4r², ½) − 2r². The value is the offset + the fixed pattern at the voxel + the gain · the electrons + the read noise. Every draw is a keyed counter draw (splitmix64 of the key mixed with the counter), and nothing is a float.
\item[{does today}] \texttt{engine/\allowbreak{}sims/\allowbreak{}} (23 September): twelve sims, each an exact-integer program (on the device, except \texttt{ask\_\allowbreak{}state}, \texttt{ka\_\allowbreak{}psi} and \texttt{goodstein}, which run on the host, and \texttt{pi\_\allowbreak{}tower}, whose turn runs on the host and whose BBP digits run on the engine's record machine) that \texttt{bash engine/\allowbreak{}sims/\allowbreak{}run.\allowbreak{}sh <sim>} builds and runs, exiting 0 only when every check holds. \texttt{nbody\_\allowbreak{}lattice}: moving and dividing bodies with their truth (the mass and the three coordinate sums a frame a body), and \texttt{-\allowbreak{}-\allowbreak{}out <dir>} for \texttt{lattice.\allowbreak{}npy} (u16, T, Z, Y, X) and \texttt{truth.\allowbreak{}tsv}. \texttt{noise\_\allowbreak{}floor}: the Kolmogorov mock (C11), the linear complexity of every bit plane, the photon transfer curve (C3, C7) with the motion removed by the truth and by the data alone, and the neighbour coherence of the frame differences (C5, C6). \texttt{period\_\allowbreak{}power}: \texttt{period\_\allowbreak{}read} (M18) on a planted period from no plant to a strong one, at 8 and 19 draws. \texttt{fixed\_\allowbreak{}pattern} and \texttt{classify\_\allowbreak{}reject\_\allowbreak{}recover}: anchor\_\allowbreak{}sift's ART-4-005 and ART-4-007 on the device, the phase dispersion ratio exact against a drawn null band. \texttt{root\_\allowbreak{}universal}: the tower with lookup edges between its floors, on camera-law volumes (M20). \texttt{ask\_\allowbreak{}state}: a qubit carried exactly as the answers of a complete ask, on the host in exact arithmetic (M21). \texttt{ka\_\allowbreak{}psi}: the superposition theorem's inner function ψ, Sprecher's and Köppen's, held exactly to depth 60 with the exponent never expanded, and the separation of its inner sums (A15; \hyperref[chap:kolmogorov-arnold]{kolmogorov\_\allowbreak{}arnold.md}). \texttt{chaitin\_\allowbreak{}omega}: Chaitin's Ω for Tromp's binary lambda calculus, where a closed term halts when it has a normal form, bracketed between two exact dyadics. Every closed term through L bits is run by normal-order reduction on the device, one term a thread (\texttt{omega\_\allowbreak{}kernel}), or on host threads with \texttt{cpu}. The device takes the host's decisions in the host's order under budgets of at most 256 steps and tokens, and hands every term that reaches a smaller budget or outgrows its room back to the host, which runs it under the full budgets. The host's run of every term through 30 bits must give the device's fates and busy beavers, champions included. \texttt{-\allowbreak{}-\allowbreak{}ledger <path>} plans the run as jobs of 2\textasciicircum{}28 terms, each finished job's tally one flushed line, and a rerun runs only the jobs its ledger lacks. A normal form proves a halt, and Brent's watcher proves a loop when a term returns to a state it held. A term is proven to grow forever when a part of its spine, reduced only by head steps at or below its own position, comes back deeper on the spine (\texttt{omega\_\allowbreak{}grows\_\allowbreak{}forever}). A term still unsettled is proven to halt when it has a simple type (\texttt{omega\_\allowbreak{}simply\_\allowbreak{}typed}, Hindley's unification with an occurs check; a typed term normalizes, Tait 1967), and no loop or growth term may type. The rest stay open. The mass past L is counted with directed rounding, not run. Given L alone, it runs on the engine's record machine (\texttt{omega\_\allowbreak{}engine}), and every term the engine settled through 24 bits is run again by \texttt{omega\_\allowbreak{}run}, which must give the same fate at the same step and the same normal form. \texttt{knf\_\allowbreak{}identity} (24 September): the entropy history (M8) of \texttt{nbody\_\allowbreak{}lattice}'s law in a 64³ cube over 177 frames against spatial null permutations: the one-bit window rule, the 48 exact cube motions, the entangled entropy E against 19 free draws, the departure curves of two inverse nulls (inside b³ tiles, and whole tiles moved) at tiles 1 to 64 and over each body's box, the scene with no bodies, and alike-voxel controls (A13, \hyperref[chap:two-crystals]{two\_\allowbreak{}crystals.md}). \texttt{goodstein} (24 September; Doug: “perform tetration of omega”): Goodstein's sequences held as hereditary trees and never expanded, with b\textasciicircum{}L − 1 at a large L held as one run of b − 1 over the exponents 0 to L − 1; with the base read as ω each tree is an ordinal below ε₀, so the start 16 is ω\textasciicircum{}ω\textasciicircum{}ω and 65536 is ω↑↑4 (A13, \hyperref[chap:two-crystals]{two\_\allowbreak{}crystals.md}). \texttt{pi\_\allowbreak{}tower} (24 September; Doug: pi turning at the boundary, the idea behind tower recursion): the rotation by π − 3 on the circle, held as the integer turn nA mod 2\textasciicircum{}P with π bracketed by Machin's formula. Its closest returns are the floors, π's partial quotients. At 2\textasciicircum{}k cells, the step that fills every cell is read from the floors by the three-gap theorem, and each gap's question is a first hit answered by Euclid's descent. Whether the integer turn lands in the real turn's cell is itself a first hit, checked at every resolution. The resolutions come with the request (Doug: “you can go to 2\textasciicircum{}n arbitrarily in the tower it is one term”): \texttt{run.\allowbreak{}sh pi\_\allowbreak{}tower -\allowbreak{}-\allowbreak{} n} reads 2\textasciicircum{}n, \texttt{-\allowbreak{}-\allowbreak{} from to} reads 2\textasciicircum{}from to 2\textasciicircum{}to, and no argument reads 2\textasciicircum{}1 to 2\textasciicircum{}100. P is 3n + 64 rounded up to a word and at least 384, and a width that cannot hold 3(P + 32) + 64 bits is refused with the limbs it needs. Its deepest turns run on the engine (24 September; Doug: “I don't want anything less than 2\textasciicircum{}googol”, “the bbp program better use the engine to compute”). The turn at depth n is α's bit n, and the BBP formula (Bailey, Borwein and Plouffe, Math. Comp. 66, 1997) reads it where it stands. Each term i ≤ d is a lane of the record machine: 16\textasciicircum{}(d − i) mod 8i + j by base-16 powering, each square taken mod 8i + j, then its fraction floored to W bits. A tail program takes the terms past d, and a pair program adds two lanes mod 2\textasciicircum{}W, so rounds of it reduce a sweep on the record machine. keymath sizes every register, the scheduler lays them with reuse, a sweep is one \texttt{cycle\_\allowbreak{}record\_\allowbreak{}run} over as many lanes as the device keeps resident, and tessera admits the job. The error, 4(N + T) + 1 units over N terms and T tail terms, decides the leading bits, and only those are printed. A resolution past 2\textasciicircum{}20, whole or as base\textasciicircum{}exponent, is held as (2, n): P and the exact width a host turn would need print exact, and the engine sums the depth-n terms sweep by sweep, printing the terms done and the rate. Neither tower, T or T⁻¹ (A13), is in it. \texttt{SIM\_\allowbreak{}EXACT\_\allowbreak{}LIMBS} sets \texttt{ANCHOR\_\allowbreak{}EXACT\_\allowbreak{}LIMBS} for every object a sim links. Every sim that uses the device is one tessera job (\texttt{sim\_\allowbreak{}job.\allowbreak{}cu}, M14): before its first allocation it declares those bytes, under a signum over its name and arguments, and its admission and release are two of its checks. The files are in engine/README.md.
\item[{wants}] The engine run on \texttt{nbody\_\allowbreak{}lattice}'s lattice and graded against its truth (not yet run). F4 in bits (C1) from a sim's planted σ²: not computed, because log₂ of 2πeσ² is irrational and the sims carry no floats.
\item[{tried, and what it gave}] \textbf{nbody\_\allowbreak{}lattice}, 9 checks, 0 failed (11 with \texttt{-\allowbreak{}-\allowbreak{}out}): 24 frames of 20 × 96 × 96, 16 bodies, 10 founders and 3 divisions at frames 8, 12 and 16. \textbf{Proved:} the device truth equals the host's walk of each body's box; the signal equals the host's on all 4,423,680 lanes, 0 differ; no lane clipped; a noiseless camera renders the signal exactly; a second render is byte-identical. \textbf{Measured}, the camera law's two moments at gain 1, read variance 3, offset 100, fixed pattern 0 to 8: the residuals sum to 3,271 against a spread of the square root of 636,863,096 (between 25,236 and 25,237), and their squares sum to 635,890,703 against the law's 636,863,096, a ratio of 0.99847 (the check allows 1\%). \textbf{noise\_\allowbreak{}floor}, 11 checks, 0 failed. The Kolmogorov mock (C11): 4,096 streams of 512 bits, each from a linear generator of degree 8 to 64 and a drawn seed; Berlekamp–Massey on the device, one thread a stream, puts the complexity within the degree on 4,096 of 4,096 and regenerates each whole stream from its first L bits on 4,096 of 4,096 (\textbf{proved}). The description, 2L bits a stream, is 286,744 of 2,097,152 bits (13.67\%). Keyed draws of the same count read 252 to 262 against n/2 = 256. The camera lattice's bit planes, 8,640 streams of 512 consecutive lanes a plane, random meaning within 16 of n/2 (\textbf{measured}): planes 0 to 8 random on 8,640 of 8,640 (mean L/n 0.5004), plane 9 on 850 (0.0638), plane 10 on 29 (0.0029), planes 11 to 15 on none. The photon transfer curve (C3, C7) on d = I(t + 1) − I(t), 4,239,360 pairs, planted g = 1 and r² = 3 (\textbf{measured}): with the motion removed by the truth (S unchanged, 4,226,574 pairs), d² on S has slope 2.0046 against 2g² = 2 (the check allows 2\%) and intercept 5.643 against 2r² = 6 (the check allows 6 either side). From the data alone, d² on I(t) + I(t + 1) over every pair has slope 7.0460 against g = 1, and the neighbour coherence (C5), (d − d′)² on the four lanes' sum over 4,195,200 pairs, has slope 4.5292 against 1: sharp-edged moving bodies bias the data-only gain, and the neighbour coherence does not remove the bias. That is the measured gap between removing the motion by the truth and removing it from the data. The neighbour correlation of d (C6), Σdd′/Σd², is −71,727/1,161,503,631 on truth-static pairs (the check: within 5σ of 0) and 403,193,327/2,288,244,270 (0.1762) over every pair. \textbf{period\_\allowbreak{}power}, 599 checks, 0 failed under the fundamental and the per-axis null: from 64 e up x reads 5 on every volume at 19 draws, and the unplanted axes' false periods stay within 5σ of 1/(draws + 1) (M18, A12; the counts in \hyperref[chap:ledger]{ledger.md}). Before those rulings it failed one check on purpose, the one that tests A12's false-period rate: the strongest peak read a multiple, and the global shuffle read periods on the unplanted axes. \textbf{fixed\_\allowbreak{}pattern} (ART-4-005 ported), 10 checks, 0 failed: frame 8 × 8, 48 frames, swing 20, pattern 40, 8 draws, key 0xF17A. The detector reads period 64, the frame, at a dispersion ratio of 179.565 against its own shuffle's 0.972. The batch and Welford routes (reduced fractions, on the device) agree on 64 of 64 phases, the broken route splits from each on 61, and the residual equals the moving scene on 3,072 of 3,072 values, a reduction of exactly 100\% (\textbf{proved}). The shuffled stack reads period 11 at 2.161, and rejecting at the wrong period, 63, removes 0.2298\%. Against the band top of 5.241 over 8 shuffles, the matched pattern (179.565) is removed at 100\%, no pattern (1.477) is declined and returned untouched, and impulses (2.064) are declined. A static scene feature cannot be told from a fixed pattern and is removed with it: at depth 0, 5, 15 and 30 the reduction is 100.0000\%, 99.9263\%, 99.3370\% and 97.3480\%. It reads the raw integers, not the reference's byte view, and draws by key, not by Python's Mersenne Twister, so the numbers are not the reference's; the claims are. \textbf{classify\_\allowbreak{}reject\_\allowbreak{}recover} (ART-4-007 ported), 7 checks, 0 failed. A fixed pattern over a varying subject is found (179.565 against 5.241), the subject exact on 3,072 of 3,072, a reduction of 100\%, 0 of 64 pixels flagged. Incoherent noise of variance 1,600 over a varying subject is declined, but its best ratio, 3.523 at period 3, clears the band top of 3.294: it is declined only because the period is not the frame, a 1-in-9 event under the null. Impulses on a repeat: 0 flagged, exact on 64 of 64. Incoherent noise of variance 144 on a repeat: 55 of 64 flagged, and where the count and the median agree, exact on only 4 of 9 (\textbf{measured}). So “the routes agree, so a strict majority, so exact” does not hold: the mode and the median can agree on a wrong value. \textbf{root\_\allowbreak{}universal}, 1,268 checks, 0 failed (M20, A14). \textbf{ask\_\allowbreak{}state}, 33 checks, 0 failed (M21, A15). \textbf{ka\_\allowbreak{}psi}, 152 checks, 0 failed (A15): Sprecher's ψ descends, Köppen's is strictly increasing and continuous in both readings, and Braun and Griebel's separation Lemma 3.4 is false at k = 1 and restored by a corrected bound with a narrower ramp (\textbf{proved}). \textbf{chaitin\_\allowbreak{}omega}, 13 checks, 0 failed, at L = 30 on the host (\texttt{build/\allowbreak{}20260924\_\allowbreak{}004819\_\allowbreak{}sim\_\allowbreak{}chaitin\_\allowbreak{}omega}, 2,048 steps and 2,048 tokens a term; before the device port, \texttt{build/\allowbreak{}20260924\_\allowbreak{}001129\_\allowbreak{}sim\_\allowbreak{}chaitin\_\allowbreak{}omega}, 12 checks, 0 failed, the same fates). Of 647,463 closed terms, 645,942 halt, 934 are proven to loop, 213 are proven to grow forever, 0 are proven to halt by a type, and 374 stay open (86 out of steps, 288 outgrew the space). 0.12254295… ≤ Ω ≤ 0.12599214…. The lower bound is the halted mass plus the normal forms past L. The upper bound is the halted mass, the open mass and the closed mass past L, and it equals their exact sum. At 200,000 steps and 16,384 tokens (\texttt{build/\allowbreak{}20260924\_\allowbreak{}001207\_\allowbreak{}sim\_\allowbreak{}chaitin\_\allowbreak{}omega}, 73 s) the same 374 stay open (47 out of steps, 327 outgrew), and the bracket is the same. At L = 33 on the device (\texttt{build/\allowbreak{}20260924\_\allowbreak{}004740\_\allowbreak{}sim\_\allowbreak{}chaitin\_\allowbreak{}omega}, 14 checks, 0 failed), of 3,910,730 closed terms 3,900,966 halt, 5,663 loop, 1,180 grow forever, 0 type and 2,921 stay open; 0.12312353… ≤ Ω ≤ 0.12599095…. The table equals the host's L = 33 run (\texttt{build/\allowbreak{}20260924\_\allowbreak{}001712\_\allowbreak{}sim\_\allowbreak{}chaitin\_\allowbreak{}omega}) line for line. At L = 36 (\texttt{build/\allowbreak{}20260924\_\allowbreak{}004859\_\allowbreak{}sim\_\allowbreak{}chaitin\_\allowbreak{}omega}, 14 checks, 0 failed), 0.12358892… ≤ Ω ≤ 0.12598988…; at 256 steps and tokens (\texttt{build/\allowbreak{}20260924\_\allowbreak{}004941\_\allowbreak{}sim\_\allowbreak{}chaitin\_\allowbreak{}omega}, 13 checks, 0 failed, the strict busy beaver check dropped under 2,048) the bounds move by 2\textasciicircum{}−28.2 and 2\textasciicircum{}−25.5 against a gap of 2\textasciicircum{}−8.7. At L = 44 and 256 steps and tokens (\texttt{build/\allowbreak{}20260924\_\allowbreak{}005020\_\allowbreak{}sim\_\allowbreak{}chaitin\_\allowbreak{}omega}, 13 checks, 0 failed), 3,392,860,908 terms, the first typed halt, and 0.12443857… ≤ Ω ≤ 0.12598796…. Every bracket holds 1/8, so Ω = 0.00… in binary: 2 bits \textbf{proved}, and the third is not settled. The checks include the closed-term counts against a parse of every code through 22 bits, no loop or growth term typing, the device against the host through 30 bits, and the busiest normal form at each length against BusyBeaverWiki's BB λ (OEIS A333479) at every length through the smaller of L and 33, up to 1812 at 33: at most the published value and equal where every term halted under any budgets, and equal everywhere at 2,048 steps and tokens or more. An open length holds only a lower bound, so meeting the published maximum means the run reached the champion. A known defect, not fixed: the ledger keeps a line cut inside its last field once a rerun has ended it, and takes it over the whole line after it. At L = 31 that gave a wrong normal form champion with every check passing (\textbf{measured}; \hyperref[chap:ledger]{ledger.md}). The fates and the bracket cannot be hit. L = 44 at 256 steps and tokens on the idle device (\texttt{build/\allowbreak{}20260924\_\allowbreak{}005724\_\allowbreak{}sim\_\allowbreak{}chaitin\_\allowbreak{}omega}, RTX 3070): 28,080 ms, 120.8 million terms a second (\textbf{measured}). \textbf{As tessera jobs} (24 September, M14), each device sim was rerun in a scratch \texttt{TESSERA\_\allowbreak{}STATE}, and each count is its earlier count plus the job's two checks. Declared bytes, then the measured peak: period\_\allowbreak{}power 262,144 and 145,915,904 (601 checks, 0 failed); nbody\_\allowbreak{}lattice 26,542,080 and 175,276,032 (11, 0 failed); noise\_\allowbreak{}floor 9,142,272 and 187,858,944 (13, 0 failed); root\_\allowbreak{}universal 339,510 and 152,207,360 (1,270, 0 failed); fixed\_\allowbreak{}pattern 141,056 and 145,915,904 (12, 0 failed); classify\_\allowbreak{}reject\_\allowbreak{}recover 118,016 and 145,915,904 (9, 0 failed). \texttt{ask\_\allowbreak{}state} (33, 0 failed) and \texttt{ka\_\allowbreak{}psi} (152, 0 failed) run on the host and submit no job. \textbf{The Ω engine run} at L = 16 (\texttt{build/\allowbreak{}20260924\_\allowbreak{}024217\_\allowbreak{}sim\_\allowbreak{}chaitin\_\allowbreak{}omega}, 16 checks, 0 failed, declared 30,256 bytes, peak 286,425,088): 226 closed terms, all halt, 0.11664483… ≤ Ω ≤ 0.12600284…, 2 bits \textbf{proved}, BB λ equal to the published value at every length through 16, and every term the engine settled gives, under \texttt{omega\_\allowbreak{}run}, the same fate at the same step and the same normal form. The host path at L = 16 (\texttt{build/\allowbreak{}20260924\_\allowbreak{}024434\_\allowbreak{}sim\_\allowbreak{}chaitin\_\allowbreak{}omega}, 13 checks, 0 failed) gives the same table line for line. \textbf{knf\_\allowbreak{}identity}, 23 checks, 0 failed (e4eae72; a tessera job, admitted with 404,752,392 bytes reserved, peak 370,311,168). \textbf{Proved:} the host's walk equals the engine's history at all 262,144 voxels; of 8,192 single flipped bits 1,500 leave it unchanged, and the window rule names all 8,192; the 48 exact motions carry it, with E and the cloud unchanged, on 48 of 48; the whole is its tiles plus its seams at all 7 tile sizes, and a rigid move keeps every tile's inside on 133 of 133 draws; alike voxels, 4 of 64 volumes identified, within 5σ of 1/20. \textbf{Measured:} E is 44.626 times the free null's mean, above all 19 draws; the inside and between departure shares cross between tiles of 2 and 4 and sum to 0.590 at 4; each moving body's curve runs 0.112 to 0.196 at 2 up to 0.947 to 1.019 at 32; with no bodies E is −10.675 times the null's mean, not identified (the table in \hyperref[chap:two-crystals]{two\_\allowbreak{}crystals.md}). \textbf{goodstein}, 7 checks, 0 failed (c633988, \texttt{build/\allowbreak{}20260924\_\allowbreak{}191006\_\allowbreak{}sim\_\allowbreak{}goodstein}). \textbf{Proved:} the bump keeps the tree on 80,256 of 80,256 values below 20,000 at bases 2 to 6; 2↑↑k in hereditary base 2 is ω↑↑k for k = 1 to 4; the starts 1, 2 and 3 reach 0 in 1, 3 and 5 steps, each value the numeric sequence's; the starts 4, 16 and 65536 each run 1,000,000 steps to base 1,000,002 with the ordinal falling on every step, and every value that fits 64 bits equals the numeric sequence's (all 10⁶ steps for 4). \textbf{Measured:} 65536's top run holds 7,625,597,484,984 terms as one block, and all three starts end on the same tail, ω² + ω·8 + 572861 (an observation). \textbf{pi\_\allowbreak{}tower}, 8 checks, 0 failed (0b10199, \texttt{build/\allowbreak{}20260924\_\allowbreak{}192322\_\allowbreak{}sim\_\allowbreak{}pi\_\allowbreak{}tower}). \textbf{Proved:} floor(π·2\textasciicircum{}384) is one integer at both ends of Machin's bracket (416 bits, 116 terms), and after the point π begins 0x243F6A8885A308D3; the bracket certifies 109 partial quotients, the first 21 the published [3; 7, 15, 1, 292, …, 1, 84]; the turn's 66 record closest returns to 2\textasciicircum{}112 steps are the convergent denominators q\_\allowbreak{}j in order; the integer turn lands in the real cell on every step searched at 2\textasciicircum{}1 to 2\textasciicircum{}100 cells; the fill step and last cell read from the floors equal a walk of every step at 2\textasciicircum{}1 to 2\textasciicircum{}24; at every resolution the last cell's first touch is the fill step. \textbf{Measured:} 2\textasciicircum{}64 cells fill at step 81,856,329,715,838,968,403 (8.186e19) and 2\textasciicircum{}100 at 1.593e30 (q\_\allowbreak{}57 − 1). The fill takes 1 to 5 times the cells except inside a large floor: up to 71.9 inside 292 (2\textasciicircum{}8), 20.9 inside 84 (2\textasciicircum{}34) and 23.8 inside 99 (2\textasciicircum{}61). The last cell filled is often just behind the start, at 1 − α, 1 − 2α, 1 − 3α (an observation). Since then, 17 checks, 0 failed (7ed7fc6, \texttt{build/\allowbreak{}20260924\_\allowbreak{}194636\_\allowbreak{}sim\_\allowbreak{}pi\_\allowbreak{}tower}), for Doug's arc and balance. \textbf{Proved:} the helix over our disk (radius 1/(2π), rising 1/π a turn), in ℚ(π) at four quarter turns, has κ² = 4π⁶/(π²+1)², τ = 2π²/(π²+1) and τ/κ = 1/π, the tangent of its plane's tilt; its Darboux vector lies along our axis; the torque about the axis is zero and L\_\allowbreak{}z² = 1/(4(π²+1)). From the bracket to 12 places: κ = 5.705134342735, τ = 1.816000663299, tilt 17.656787151412°, L\_\allowbreak{}z = 0.151657235526. In the square billiard from the corner at slope π, unfolded, the segment in cell (i, j) carries L = (−1)\textasciicircum{}(i+j)((j + ½) − π(i + ½)) about the centre. Over 2\textasciicircum{}24 columns walked from the integer turn with every floor certain, no wall hit is a corner, no L is zero, and the record near-corners are exactly the q\_\allowbreak{}j. \textbf{Measured:} 52,707,178 floor-wall hits to 16,777,216 side-wall hits; the lead changes hands 35,929,962 times in 69,484,394 segments and is never held longer than 2; the nearest tie, |L| = 1.910e-8, is at π ≈ 5419351/1725033, a convergent. Then 18 checks, 0 failed (b0b6da6), for Doug's second run. At every resolution, H is the first step after the fill in cell 0 and C the first step after H in the fill's last cell, each one first hit, with the certainty check run to C. Both equal a walk at 2\textasciicircum{}1 to 2\textasciicircum{}24. \textbf{Measured:} H lands on a record return or a sum of them, and the wait H − T is either a few steps or about one floor's q. The fill is within 4 steps of the next record return at 64 of 100 resolutions. The second run C − H has the first run's period at 13 of 100, and is shorter at 86 by a floor combination (a record return at 80, c·q\_\allowbreak{}j at 6) and longer at 1 (2\textasciicircum{}6). Then 23 checks, 0 failed (c26b6a7), for Doug's residue, his golden helix and any 2\textasciicircum{}n; the earlier numbers at 2\textasciicircum{}1 to 2\textasciicircum{}100 are unchanged. \textbf{Proved}, the residue: on every floor with q\_\allowbreak{}j ≤ 2\textasciicircum{}21 (j = 0 to 11), π's first q\_\allowbreak{}j steps have the whole parts of the permutation n ↦ n·p\_\allowbreak{}j mod q\_\allowbreak{}j, and at step q\_\allowbreak{}j the walk stands δ\_\allowbreak{}j = q\_\allowbreak{}jπ − P\_\allowbreak{}j off the whole, with q\_\allowbreak{}j|δ\_\allowbreak{}j| < 1 certified. The identities are 3/1, 22/7, 333/106, 355/113, 103993/33102, … 5419351/1725033, with the residues +1.415e-1, −8.851e-3, +8.821e-3, −3.014e-5, +1.912e-5, … −3.820e-8. \textbf{Proved}, the golden helix, with every residue bracketed from the turn on floors q\_\allowbreak{}j ≤ 2\textasciicircum{}112 (j = 0 to 65): the residues flip sign every floor; q\_\allowbreak{}j ≥ F\_\allowbreak{}\{j+1\} and |δ\_\allowbreak{}j| < 1/q\_\allowbreak{}\{j+1\}; and the shrink |δ\_\allowbreak{}j|/|δ\_\allowbreak{}\{j−1\}| lies between ½ and 1 exactly where a\_\allowbreak{}\{j+1\} = 1, each side of 1/φ decided at both ends. \textbf{Measured:} of 65 shrinks, 45 fall below 1/φ and 20 above (every one above on a floor of 1, while 8 floors of 1 fall below), and 28 lie between ½ and 1. The growth a floor, q\_\allowbreak{}65\textasciicircum{}\{1/65\}, is 3.210576 against φ = 1.618033. \textbf{Proved}, any 2\textasciicircum{}n: at 2\textasciicircum{}1000 (\texttt{SIM\_\allowbreak{}EXACT\_\allowbreak{}LIMBS=512}, 16,384 bits; P = 3,072, Machin at 3,104 bits and 864 terms, 888 floors certified), all 23 checks hold. The default width refuses 2\textasciicircum{}1000 by name, “the turn needs an exact width of 9376 bits: run with SIM\_\allowbreak{}EXACT\_\allowbreak{}LIMBS=512”. \textbf{Measured} at 2\textasciicircum{}1000: the fill is step 1.288e302, 12.02 times the cells, on floor 600 (a = 106), and the second run is not the first's period. Then 41 checks, 0 failed (e340d08), for the engine. \textbf{Proved} on the engine's record machine: π's first 64 bits after the point, 0x243F6A8885A308D3; the hex digits at Bailey's published positions (D. H. Bailey, “The BBP Algorithm for Pi”, 2006, table 1 and its text), 26C65E52CB4593 at 10\textasciicircum{}6, 6C65E52CB459350050E4BB1 at 10\textasciicircum{}6 + 1, 17AF5863EFED8D at 10\textasciicircum{}7 and ECB840E21926EC at 10\textasciicircum{}8; and the engine's cell at 2\textasciicircum{}100, its low 64 bits, equals the exact turn's, 0x5a308d313198a2e0. \textbf{Measured} (RTX 3070, 70,656 lanes a sweep): 10\textasciicircum{}8 terms in 1,416 sweeps and 3.887 s. At 2\textasciicircum{}(2\textasciicircum{}30) cells (hex position 268,435,440), 268,435,441 terms in 3,800 sweeps and 10.638 s, 2.52e7 terms a second, 127 bits certified, the cell's low 64 bits 0x8293097a8232c37f. At 2\textasciicircum{}googol cells (n = 10\textasciicircum{}100, hex position 2.5e99), 46,844,928 of 2.5e99 terms in 663 sweeps and 64.0 s, 7.31e5 terms a second, about 1.08e86 years at that rate; the run was stopped, and no digit at that depth was printed. The job declared 3,413,540 bytes and peaked at 221,413,376: the declaration counts the sim's own buffers, not the record kernel's per-thread file. The runs are in \hyperref[chap:ledger]{ledger.md}.
\item[{status}] built; all twelve sims pass every check, the device sims as tessera jobs, Ω runs on the engine, and π's BBP digits run on the engine
\item[{next}] The engine on \texttt{nbody\_\allowbreak{}lattice}'s lattice against its truth.
\end{description}

\tablerecord{\textbf{M20. The root universal}}
\begin{description}
\item[{the algebra it holds to}] The tower is the additive skeleton of the Kolmogorov–Arnold form: sums, collapsed in order onto the floor. The lookup table is its one-variable edge. Laid between the floors, the edges make the collapse an outer sum with nonlinear edges (Doug, 23 September: “The tower is its own object, the lut is its own object, if we integrate them both we get a root universal”). Inside the reversible stream an edge must be a bijection, so it permutes the low bits of a coefficient and passes the high bits. Any other function rides widened, (x, y) → (x, y ⊕ f(x)), the Bennett embedding: a dimensional expansion (Doug, 23 September), so the bound is one of width, not of kind (A14).
\item[{does today}] \texttt{tower\_\allowbreak{}lift} and \texttt{tower\_\allowbreak{}lower} take edges (23 September; Doug chose bijective edges in the stream over read-only taps). \texttt{Tower\allowbreak{}Edge} \{floor, index\_\allowbreak{}bits, forward\} permutes the low index\_\allowbreak{}bits of every approximation coefficient on floor k through \texttt{forward}, a permutation of [0, 2\textasciicircum{}index\_\allowbreak{}bits). Floor k is the approximation just before level k's lifting; floor = floors is the collapsed floor. Lift lays E0, L0, E1, …, L(N−1), EN, and lower undoes EN⁻¹, L(N−1)⁻¹, …, L0⁻¹, E0⁻¹, so the edges on one floor unwind as a stack. index\_\allowbreak{}bits runs 1 to 20. Every edge is checked before any device work (a permutation, the width in range, the floor at most floors) and refused otherwise. With no edges the tower is unchanged. The crystal path in \texttt{engine.\allowbreak{}cu} passes no edges.
\item[{wants}] The edges in the crystal and seal path, with the \texttt{.\allowbreak{}kcr} recording the edge program (Doug's call). Which functions to lay, without fitting (item 3 of \hyperref[chap:kolmogorov-arnold]{kolmogorov\_\allowbreak{}arnold.md}). A fitted edge's price.
\item[{tried, and what it gave}] \texttt{tower\_\allowbreak{}edge\_\allowbreak{}test}, 25 checks, 0 failed (\textbf{proved}; the 20 before the Bennett edges, then 5 for them, A14): one edge, stacked edges on one floor, and edges on every floor including the collapsed one each rebuild the exact lanes; an edge changes the crystal; a non-permutation, a width of 0, a width of 21 and floor 50 are each refused; a clean edge after the refusals still round-trips; a function that is not a bijection rides as a Bennett edge, (x, y) → (x, y ⊕ f(x)), at the price of width: |x| and x ≥ 100 as 16-bit edges round-trip on 64 of 64 lattices and read out exactly, and |x| laid bare is refused (A14). \texttt{root\_\allowbreak{}universal} (M19's camera law), 1,268 checks, 0 failed. \textbf{Proved}, the root stays exact: 48 volumes, 24 of 4 × 12 × 48 × 48 and 24 of 3 × 13 × 37 × 41, 6 floors each, carry 171 random edges (1 to 6 a volume, widths 1 to 12 bits, and one 20-bit edge per extent) through lift, code, wipe of the held coefficients, decode and lower: plain exact on 48 of 48, edged exact on 48 of 48, the crystal changed on 48 of 48. \textbf{Proved}, the fold keeps order: on floor 1, 8-bit edges a then b lay the same crystal as the one table b(a(i)) on 8 of 8; b then a lays a different crystal on 8 of 8; a on floor 1 and b on floor 2 lay a different crystal on 8 of 8, so a floor between two edges blocks the fold. \textbf{Measured}, the price of an unfitted edge (one keyed random permutation a floor, coded bits a voxel, 8 volumes of 4 × 12 × 48 × 48): 6.802 with no edge; at widths 2, 4, 8 and 12, floor 0 (110,592 coefficients) 6.810, 6.908, 8.719 and 12.608, floor 3 (72) 6.802, 6.802, 6.804 and 6.817, the collapsed floor 6 (1) 6.802 at every width, every floor 6.808, 6.911, 8.742 and 12.714.
\item[{status}] built and proved in \texttt{tower}; not in the crystal path
\item[{next}] Doug's ruling on the \texttt{.\allowbreak{}kcr} recording the edge program; a fitted edge's price against the unfitted.
\end{description}

\tablerecord{\textbf{M21. The ask and the state}}
\begin{description}
\item[{the algebra it holds to}] The terms (Doug, 23 September). An \textbf{ask} (a measurement, formally a POVM) is finitely many outcomes with positive semidefinite operators E\_\allowbreak{}i summing to the identity. A \textbf{state} ρ (the quantum functional) is a positive semidefinite matrix of trace 1; for a qubit ρ = (I + r·σ)/2, |r| ≤ 1. The \textbf{answer distribution} is p\_\allowbreak{}i = Tr(ρ E\_\allowbreak{}i): “the order falling out of an ask” (Doug), not an order imposed. One basis's ask returns the \textbf{magnitude}, the squared size of each amplitude; the \textbf{phase}, the relative angles between the amplitudes, it cannot see, and other asks can. A \textbf{complete ask} (informationally complete) is one whose answers determine the state, one to one; it needs at least d² outcomes. The \textbf{valid set} is the distributions some state gives, a proper subset of all of them. The \textbf{crossing rule} is the fixed linear map from a complete ask's answers to any other ask's; it is not classical total probability, and the difference is interference (A15).
\item[{does today}] \texttt{engine/\allowbreak{}sims/\allowbreak{}ask\_\allowbreak{}state} (23 September), on the host, no GPU work: a qubit carried exactly as the answers of a complete ask. Rationals on the exact integer, reduced by a word gcd when both parts fit a word and carried exact, unreduced, when they do not (Doug: “why would you use 64 bit?”). The SIC's numbers in Q(√3), a + b√3, multiplied with (√3)² = 3, each sign decided exactly by comparing a² with 3b² (Doug: “represent square roots as their real number with their square operation”). Two asks, both E\_\allowbreak{}i = (I + v\_\allowbreak{}i·σ)/4: a rational tetrahedral ask, v = (±1, ±1, ±1)/2 with an even number of minus signs, and the SIC, v = (±1, ±1, ±1)/√3.
\item[{wants}] Doug's two statements, not built (A15): the sparse form of ψ's terms, and the countable numbers held exactly.
\item[{tried, and what it gave}] \texttt{ask\_\allowbreak{}state}, 33 checks, 0 failed. \textbf{Proved} for each ask: the operators sum to the identity and are positive, and the frame is c·I, so the ask is complete (the rational ask: |v|² = 3/4, not rank one, frame I, r = 4 Σ p\_\allowbreak{}i v\_\allowbreak{}i; the SIC: |v|² = 1, rank one, frame (4/3)I, r = 3 Σ p\_\allowbreak{}i v\_\allowbreak{}i). 25 states, 21 pure (14 on the equator from rational t, ((1 − t²)/(1 + t²), 2t/(1 + t²), 0), and rational points of the sphere) and 4 mixed, go state → answers → state exactly, 25 of 25; the pure lie on the boundary, |r|² = 1, 21 of 21, and the mixed strictly inside, 4 of 4. The crossings onto the +x, +y and +z asks equal (1 + r\_\allowbreak{}a)/2, 75 of 75. The phase falls out of the ask: the 14 equator states all answer ½, ½ to the 0/1 ask, yet the complete ask tells apart all 91 of their pairs. The crossing weights onto +x are 3/2, 3/2, −1/2, −1/2 (the rational ask) and ½ + ½√3, ½ + ½√3, ½ − ½√3, ½ − ½√3 (the SIC): a weight is negative, so the crossing is not classical total probability. (1, 0, 0, 0) is refused, no state, and crosses onto +x at 3/2 and at ½ + ½√3, values that are no probabilities. The SIC's answers cross into the rational ask's exactly, 25 of 25. \textbf{Measured:} of the 455 distributions with denominator 12, 55 are states for the rational ask and 87 for the SIC; of the others, 104 and 72 still cross to probabilities on the x, y and z asks, so the valid set is stricter than “the three axis crossings are probabilities” (both counts recounted independently in exact fractions).
\item[{status}] built as a sim; proved exact on the host
\item[{next}] Not set (Doug's).
\end{description}

\section{\texorpdfstring{The machine's algebra}{The machine's algebra}}

Every line is exact. No line assumes a scale: orders, windows, widths and spacings are symbols the request fills.

\subsection{\texorpdfstring{A0. Exact numbers (M1): \texttt{decimal\_\allowbreak{}double}}{A0. Exact numbers (M1): decimal\_double}}

A decimal d = (−1)\textasciicircum{}neg · digits · 10\textasciicircum{}ex is 2\textasciicircum{}\{ex\} · 5\textasciicircum{}\{ex\} · digits. Its binary value's leading W bits are found by multiplying or dividing by the widest power of 5 that fits a limb, chunk by chunk, with the fives stripped from a fraction first: the result is the exact integer ⌊|d| · 2\textasciicircum{}\{−e2\}⌋ with e2 chosen so the integer has W bits, and \texttt{terminates} says whether the rest is zero. When the value needs more than W bits, \texttt{needed\_\allowbreak{}bits} says how many and \texttt{fits} = 0. Nothing is rounded.

\subsection{\texorpdfstring{A1. The residual (M2): \texttt{residual\_\allowbreak{}program}, \texttt{keymath}, \texttt{cycle}, \texttt{unit\_\allowbreak{}sweep}}{A1. The residual (M2): residual\_program, keymath, cycle, unit\_sweep}}

Per axis a, with [1 1] the unit step:

L\_\allowbreak{}n = [1 1]\textasciicircum{}n ∗ I,  B\_\allowbreak{}s ∗ B\_\allowbreak{}b = B\_\allowbreak{}\{s+b\},  R = 2\textasciicircum{}\{|b|\} L\_\allowbreak{}s − L\_\allowbreak{}\{s+b\}

\begin{itemize}
\item \textbf{The two terms are two rungs of one ladder.} B\_\allowbreak{}b applied to B\_\allowbreak{}s I is B\_\allowbreak{}\{s+b\} I, so R is the difference of rungs s and s + b of one binomial ladder of the same image. Every residual for every (s, b) is a difference of two rungs, so one sweep that keeps its rungs gives all of them. Orders per request, and a sweep over candidate periods, cost one pass (theory).
\item \textbf{The kernel sums to zero.} B\_\allowbreak{}s sums to 2\textasciicircum{}\{|s|\}, so 2\textasciicircum{}\{|b|\} B\_\allowbreak{}s and B\_\allowbreak{}\{s+b\} both sum to 2\textasciicircum{}\{|s|+|b|\}. A constant added to I leaves R unchanged: the medium's mean cancels exactly.
\item \textbf{The heat identity.} In centred steps H = [1 2 1] = [1 1]², with m = b/2 per axis:

4\textasciicircum{}m δ − H\textasciicircum{}m = Σ\_\allowbreak{}\{k=0\}\textasciicircum{}\{m−1\} 4\textasciicircum{}\{m−1−k\} H\textasciicircum{}k (4δ − H),  and 4δ − H = [−1 2 −1] = −Δ²

so, per axis, R = Σ\_\allowbreak{}\{k=0\}\textasciicircum{}\{m−1\} 4\textasciicircum{}\{m−1−k\} · (−Δ² S\_\allowbreak{}k), with S\_\allowbreak{}k = H\textasciicircum{}k S and S = H\textasciicircum{}\{s/2\} I. H/4 is one explicit step of the discrete heat equation, so R is 4\textasciicircum{}m times what m heat steps remove from the smoothed field: the residual is exactly a weighted sum of negative discrete Laplacians down the ladder. In three axes the same telescoping runs axis by axis: 2\textasciicircum{}\{G\}δ − B\_\allowbreak{}\{b\_\allowbreak{}z\}⊗B\_\allowbreak{}\{b\_\allowbreak{}y\}⊗B\_\allowbreak{}\{b\_\allowbreak{}x\} = (2\textasciicircum{}\{b\_\allowbreak{}z\}δ − B\_\allowbreak{}\{b\_\allowbreak{}z\})⊗2\textasciicircum{}\{b\_\allowbreak{}y+b\_\allowbreak{}x\}δ + B\_\allowbreak{}\{b\_\allowbreak{}z\}⊗(2\textasciicircum{}\{b\_\allowbreak{}y\}δ − B\_\allowbreak{}\{b\_\allowbreak{}y\})⊗2\textasciicircum{}\{b\_\allowbreak{}x\}δ + B\_\allowbreak{}\{b\_\allowbreak{}z\}⊗B\_\allowbreak{}\{b\_\allowbreak{}y\}⊗(2\textasciicircum{}\{b\_\allowbreak{}x\}δ − B\_\allowbreak{}\{b\_\allowbreak{}x\}) (theory: exact algebra, not built as a pass).
\item \textbf{Parity.} 2δ − [1 1] = [1 −1], a first difference, centred half a voxel off. So an odd background order puts the subtracted term half a voxel from the kept one, and an odd smooth order moves both by the same half voxel. That is ruling (a): refuse odd b, hold odd s with its offset.
\end{itemize}

The positive set is P(x) = [R(x) > 0].

\subsection{\texorpdfstring{A2. The component tree (M3): \texttt{max\_\allowbreak{}tree}}{A2. The component tree (M3): max\_tree}}

\begin{itemize}
\item c(x) is the dense rank of R(x) among the admitted voxels. Each tree node n has a level λ(n); ν(x) is x's own node, the node at level c(x) that holds x. The component at level r holding x is the ancestor a of ν(x) with λ(a) ≥ r > λ(parent(a)). Its voxels are exactly those x′ with ν(x′) in a's subtree.
\item Number the nodes in DFS order; a's subtree is one range [in(a), out(a)).
\item \textbf{Probes by LCA.} Two voxels share a component at level r exactly when r ≤ λ(LCA(ν(x), ν(x′))). Sort k probes by in(ν). The partition at every level at once is given by the k − 1 levels λ(LCA(p\_\allowbreak{}i, p\_\allowbreak{}\{i+1\})) of neighbours in that order: at level r, neighbours i, i + 1 are together exactly when r ≤ that level, and a probe whose own level is below r is absent. The slide's bisection over levels becomes k − 1 LCA queries (theory).
\item \textbf{The overlap at every level from one count.} For a lag v, the own-node pair table J(v)[n, n′] = |\{ x : ν\_\allowbreak{}t(x) = n, P\_\allowbreak{}\{t+1\}(x + v), ν\_\allowbreak{}\{t+1\}(x + v) = n′ \}|. Then for components A at level r of t and B at level r′ of t + 1:

\texttt{C\_\allowbreak{}\{r,r′\}[A, B](v)} = Σ\_\allowbreak{}\{n ∈ sub(A)\} Σ\_\allowbreak{}\{n′ ∈ sub(B)\} J(v)[n, n′]

a rectangle sum over the two DFS ranges, so one summed-area table of J answers every pair of levels. The census (one partner, unpaired, mutual, forked) is the degrees of the support of C at the chosen levels. Today's overlap recounts per sampled cut (theory).
\end{itemize}

\subsection{\texorpdfstring{A3. The moments (M4)}{A3. The moments (M4)}}

With own-node moments μ(n) = Σ\_\allowbreak{}\{ν(x)=n\} (1, x, x xᵀ), every node's moments are Σ\_\allowbreak{}\{n′ ∈ sub(n)\} μ(n′): one prefix over the DFS order gives every component at every level. The centred second moment is K = m M − S Sᵀ, carried with its mass, never divided.

\subsection{\texorpdfstring{A4. The correlation operator (M5)}{A4. The correlation operator (M5)}}

For frames t, t′, a shift v, a label a of t and a label b of t′ (or ∅, a positive voxel under no label):

\texttt{C\_\allowbreak{}\{t,t′\}[a, b](v)} = Σ\_\allowbreak{}x [L\_\allowbreak{}t(x) = a] · [P\_\allowbreak{}\{t′\}(x + v)] · [L\_\allowbreak{}\{t′\}(x + v) = b]

With a weight φ(x) ∈ \{1, x, x xᵀ\} at v = 0, t′ = t and b = a, the same sum gives the moments. Its reductions:

\begin{itemize}
\item the argmax over v of the sum over all labels (\texttt{shift\_\allowbreak{}agreement});
\item a local argmax per label from a start lag (\texttt{climb\_\allowbreak{}machine});
\item a transpose: \texttt{C\_\allowbreak{}\{t′,t\}[b, a](−v)} = \texttt{C\_\allowbreak{}\{t,t′\}[a, b](v)};
\item box sums over nested widths around a lag, from one summed-area table over v (the core: each width is the box on top of the one below);
\item the autocorrelation of a box, a tent, over one frame's own C (the contact form).
\end{itemize}

The labels may be any level's components: with A2, C at every pair of levels is a rectangle sum of J. So the level is one more axis of the same box, and the whole of the tree's frame-to-frame reading is C over (v, r, r′), counted once.

\subsection{\texorpdfstring{A5. The marginal (M6): \texttt{marginal}}{A5. The marginal (M6): marginal}}

For sources S and targets T with integer weights ω(s, t) and a null weight ω(s, ∅) per source, an arrangement α maps each source to a target or ∅, no two sources to one target:

Z = Σ\_\allowbreak{}α Π\_\allowbreak{}s ω(s, α(s)),  Z\_\allowbreak{}\{s→t\} the same over α(s) = t,  P(s → t) = Z\_\allowbreak{}\{s→t\} / Z

Z is the permanent of the |S| × (|T| + |S|) matrix [ω(s, t) | diag ω(s, ∅)]. \textbf{It factors over the connected components of the support graph} (s \textasciitilde{} t where ω(s, t) ≠ 0): the only coupling between sources is a shared target, so the sum over all arrangements is the product of the sums over components, and Z\_\allowbreak{}\{s→t\} / Z over s's component is the global ratio exactly. A local set smaller than the component is exact only when it is the component. The cost lives in the component's size; a permanent is hard in general, so the component's size is what to measure and report (theory).

\subsection{\texorpdfstring{A6. Matching (M7): \texttt{heaviest\_\allowbreak{}matching}}{A6. Matching (M7): heaviest\_matching}}

The one-to-one set of largest total weight on exact integer costs; no weight is scaled or compared as a fraction.

\subsection{\texorpdfstring{A7. The entropy history (M8)}{A7. The entropy history (M8)}}

For voxel x, bit j and frames 0 … F − 1, with f\_\allowbreak{}j(x) the number of transitions where bit j flips: f\_\allowbreak{}j(x) ≡ bit j of (I\_\allowbreak{}0(x) ⊕ I\_\allowbreak{}\{F−1\}(x)) (mod 2). A voxel where this fails means the history was not taken whole; it is the pass's own proof. The windows partition the transitions; their width is a symbol (see the audit).

\subsection{\texorpdfstring{A8. The files (M9)}{A8. The files (M9)}}

Integer lifting (predict, then update) takes each step's integer quotient by a power of two, and the inverse step subtracts the identical quotient, computed from the same integers, so each step is a bijection on integers and nothing is lost: the tower rebuilds every voxel. The proof is not the algebra but the read-back: every sample is rebuilt from the file on the device and compared voxel for voxel, and again from a second read of the source pixel for pixel.

\subsection{\texorpdfstring{A9. The golden ladder (M11)}{A9. The golden ladder (M11)}}

β(n) = the number of Fibonacci rungs at or below n after the first, over the rungs that fit 64 bits.

\subsection{\texorpdfstring{A10. The seal (M13): \texttt{obsignatio}}{A10. The seal (M13): obsignatio}}

Each level ℓ keys its nodes: K\_\allowbreak{}ℓ = BLAKE3-derive-key-context(“obsignatio aeterna 2026-09-23 “ ‖ ℓ), H\_\allowbreak{}ℓ(m) = BLAKE3-keyed(K\_\allowbreak{}ℓ, m). For a sample of extent (T, Z, Y, X):

ρ(t, z, y) = H\_\allowbreak{}row(I[t, z, y, ·]),  π(t, z) = H\_\allowbreak{}plane(ρ(t, z, ·)),  υ(t) = H\_\allowbreak{}volume(π(t, ·)),  Λ = H\_\allowbreak{}lanes(υ(·)) κ(c) = H\_\allowbreak{}chunk(LE64(ℓ\_\allowbreak{}c) ‖ bits\_\allowbreak{}c),  Σ = H\_\allowbreak{}stream(offsets ‖ B ‖ κ(·)) σ = H\_\allowbreak{}side(H\_\allowbreak{}side stored(deflated) ‖ H\_\allowbreak{}side inflated(inflated)),  μ = H\_\allowbreak{}members(tables ‖ names) Ω = H\_\allowbreak{}sample(head ‖ Λ ‖ Σ ‖ σ ‖ μ),  Θ = H\_\allowbreak{}set(Ω\_\allowbreak{}1 ‖ … ‖ Ω\_\allowbreak{}n)

\begin{itemize}
\item \textbf{Locality.} A change at (t, z, y, x) changes exactly ρ(t, z, y), π(t, z), υ(t), Λ, Ω and Θ. A change in chunk c changes exactly κ(c), Σ, Ω and Θ. Any other node agrees except with probability 2\textasciicircum{}−256.
\item \textbf{Size is strength.} A k-bit check misses a random corruption with probability 2\textasciicircum{}−k, whatever the function. The DAG's gain over CRC-64 is structure: nonlinearity, locality and composition, not a smaller word.
\item \textbf{The tree is BLAKE3's.} Pairing each level and carrying the odd node up is BLAKE3's left-heavy tree exactly (its left subtree is the largest power of two of chunks below the whole), so the level-parallel device fold equals the host's stack fold.
\item \textbf{Where it stops.} The tower lifts across all four axes, so a decoded row can be checked only after a whole decode. The stored bytes check alone.
\item \textbf{The shared shape.} The seal folds the lattice x, y, z, t, as the tower's floors lift it and the moments' prefix sums it (A3). The residual, the tree, C and the seal are each one lattice read once with a different combining operator. The seal's operator alone is not associative, so its tree shape is fixed, and that fixed shape is what makes a mismatch locatable (observation).
\end{itemize}

\subsection{\texorpdfstring{A11. The scheduler (M14): \texttt{tessera}}{A11. The scheduler (M14): tessera}}

The device holds C bytes; U(τ) is in use. Each admitted job j has a declaration d\_\allowbreak{}j, a reservation r\_\allowbreak{}j (from d\_\allowbreak{}j, only growing) and a measured use u\_\allowbreak{}j(τ) by pid:

O(τ) = U(τ) − Σ\_\allowbreak{}j u\_\allowbreak{}j(τ),  H(τ) = C − O(τ) − Σ\_\allowbreak{}j max(r\_\allowbreak{}j, u\_\allowbreak{}j(τ))

\begin{itemize}
\item \textbf{Admission:} k is admitted when d\_\allowbreak{}k ≤ H(τ), so admission alone keeps H ≥ 0. Doug, 24 September: “We reserve what they ask for and then if it cost less we remember that for next time”. The code in the working tree reserves max(d\_\allowbreak{}k, p\_\allowbreak{}σ) instead (\texttt{tessera\_\allowbreak{}ledger\_\allowbreak{}wants}), which Doug did not approve (M14).
\item \textbf{Growth:} u\_\allowbreak{}j(τ) > r\_\allowbreak{}j sets r\_\allowbreak{}j = u\_\allowbreak{}j(τ) and warns. Only growth and outside processes drive H negative, and both are measured, so the books equal C − H at every sweep.
\item \textbf{Memory of a request:} p\_\allowbreak{}σ = max over a run of u\_\allowbreak{}j(τ), kept under the request's signum σ. It is a measurement at the sweep's resolution, not a byte-exact constant: identical requests have read peaks 4,194,304 bytes apart (M14). The engine is deterministic in its request and its result, not in the bytes a sweep samples. A new declaration d > p\_\allowbreak{}σ is over budget: asked, held, overridable.
\item \textbf{Liveness:} a ticket is (pid, start time), so a reused pid is a different process. A closed socket or a gone process releases r\_\allowbreak{}j at once.
\item \textbf{Backfill with the head kept:} the history keeps (p\_\allowbreak{}σ, t\_\allowbreak{}σ). The head's shadow S is the least time with H + Σ\_\allowbreak{}\{e\_\allowbreak{}j ≤ S\} max(r\_\allowbreak{}j, u\_\allowbreak{}j) ≥ d\_\allowbreak{}h, and the spare is X = that sum − d\_\allowbreak{}h. A job k behind the head goes now if d\_\allowbreak{}k ≤ H and (now + t\_\allowbreak{}σk ≤ S or d\_\allowbreak{}k ≤ X). A signum with no history is never backfilled.
\item \textbf{Time:} every deadline sits in one min-heap keyed by time to change; the root is the next event, handled in order.
\end{itemize}

\subsection{\texorpdfstring{A12. The period reading (M18): \texttt{period\_\allowbreak{}read}}{A12. The period reading (M18): period\_read}}

On an axis of extent n with the other axes holding M = N/n voxels, take L = ⌊n/2⌋ lags and the start positions whose coordinate on the axis is below n − L. Every lag ℓ ≤ L then compares the same Q = (n − L)·M pairs:

a(ℓ) = \#\{x : I(x) = I(x + ℓ·e)\},  r(ℓ) = a(ℓ)/Q

\begin{itemize}
\item \textbf{The local peak} (Doug, 23 September). A candidate p must stand above both its neighbour lags, and so must 2p:

a(p) > max(a(p − 1), a(p + 1))  and  a(2p) > max(a(2p − 1), a(2p + 1))

Its height is the smaller rise, h(p) = min(a(p) − max(a(p − 1), a(p + 1)), a(2p) − max(a(2p − 1), a(2p + 1))), an integer count over the axis's Q pairs, and its margin is the rational h(p)/Q. The candidates run over 2 ≤ p with 2p + 1 ≤ L, so P ≤ ⌊(L − 1)/2⌋: 8 on an axis of 34 slices (L = 17), 5 on 24, 4 on 20.
\item \textbf{The floor: a drawn null per axis} (Doug, 23 September). The request names d draws. For axis a, draw k permutes each line along a on its own by π\_\allowbreak{}k, so the shuffled lattice I∘π\_\allowbreak{}k holds exactly the same values on every line, its histogram is the original's, and every other axis keeps its structure; only a's own order is destroyed. Each draw reads a's strongest height, h\_\allowbreak{}k = max\_\allowbreak{}p h\_\allowbreak{}\{I∘π\_\allowbreak{}k\}(p); the draws that reach a peak make a's band, sorted exactly. A read of rank r makes r·d shuffles, one per axis per draw.
\item \textbf{The choice: the fundamental} (Doug, 23 September). P is the smallest candidate whose height clears the band's top:

P = min \{ p : h(p)/Q > max\_\allowbreak{}k h\_\allowbreak{}k/Q \}

with each comparison a 128-bit cross-multiply. A period P also peaks at 2P and 3P with heights matching P's, so the strongest peak picked a multiple; the smallest peak above the top picks P. An empty band, where no draw reached a peak, or a given top of 0 lets any peak through, so the smallest peak is P. When none clears, P = 0 and the reported candidate is the strongest peak (the largest height, the smallest p on a tie). A lag that makes no peak at p and 2p is never a candidate. Nothing is rounded, and no constant is written in.
\item \textbf{The fundamental leaves the rate alone.} Some peak clears the top exactly when the strongest peak does, so an axis reads a period under the fundamental exactly when it did under the strongest peak; the rule changes which period is read, not whether one is (theory).
\item \textbf{The set-wide band} (Doug, 23 September). The knee's first ruling was 73,158 draws a series, so that fewer than one false period is expected over the set's 24,386 series × 3 axes = 73,158 axes. \textbf{Measured:} 152 ms a draw on 34 × 960 × 960 (RTX 3070), about 3.1 hours a series and about 75,000 GPU-hours for the set. So one band serves the whole set. The 152 ms was the global shuffle's, one shuffle and one count a draw. The per-axis null makes one of each per axis, and it is \textbf{measured} (23 September, RTX 3070, on 34 × 960 × 960 = 31,334,400 voxels, the same kind of lattice: a smooth ramp x + y + 2z plus uniform noise 0 to 63). \texttt{period\_\allowbreak{}draw} takes 455.9 ms a call, the mean of 8 after a warm-up; a call is one count of the original, then a shuffle and a count on each of the 3 axes. \texttt{period\_\allowbreak{}read} takes 1,934.2 ms at 4 draws and 9,136.5 ms at 20, so a marginal draw costs 450.1 ms and the fixed cost is 133.7 ms (each as the run printed it from the unrounded times). That is 2.96 times the global shuffle's 152 ms, as three shuffle-and-count passes against one would give. From the measured draw, not a measured set run: about 9.1 hours a series (3.1 before) and about 223,000 GPU-hours for the 24,386 series (75,000 before). So one band serves the whole set the more. Its limits: one run on one lattice; host wall-clock with the syncs; and \texttt{period\_\allowbreak{}draw} allocates its buffers on every call, so part of the 455.9 ms may be allocation (not separated). The timing program was a scratch program, not a repo test. \texttt{period\_\allowbreak{}draw} shuffles each axis once along its own lines, keyed by the content, the draw number and the axis, and returns each axis's strongest height (0 over Q where no peak is reached). The program calls it across its series, keeps the top per axis, and passes those tops to \texttt{period\_\allowbreak{}read} as \texttt{null\_\allowbreak{}top} with no draws of its own; a given top with draws is a request error. P is the smallest peak above the given top, and a top of 0 lets any peak through. The tops compare exactly across shapes, because the cross-multiply carries each Q. Whether one band is fair across shapes is Doug's call.
\item \textbf{The false-period rate.} With no structure along axis a, each line's order along a is one more arrangement of its values, whatever the other axes carry, and if the draws behave as uniform shuffles the d + 1 heights are exchangeable (theory). The lattice's own strongest height then stands strictly above all d others with probability at most 1/(d + 1), exactly that when no two tie. The caller sets the rate with d: 1/9 at 8 draws, 1\% at 99. \textbf{Measured,} 5 of 192 random axes at 8 draws (2.6\%, M18). With a period 5 planted along x at 1,024 e, the unplanted z and y read a period on 5 of 64 at 8 draws (against 64/9) and 2 of 64 at 19 draws (against 64/20) (\texttt{period\_\allowbreak{}power}, M19): the per-axis null holds the bound when another axis carries a strong period.
\item \textbf{Measured: the fundamental is read} (\texttt{period\_\allowbreak{}power}, M19). A period 5 is planted along x of 32 × 64 × 64 under the camera law, at amplitudes 0 to 1,024 e, 32 volumes each. At 8 draws x reads 5, a multiple (10 or 15), or another lag: 10, 6, 0 at 16 e; 26, 3, 0 at 32 e; 32, 0, 0 at 64 and 128 e; 31, 1, 0 at 256 e; 32, 0, 0 at 1,024 e. At 19 draws it reads 5 on 32 of 32 from 64 e up.
\item \textbf{The strongest peak (superseded, Doug, 23 September): it read a multiple.} The candidate was the largest height. Once the plant was found, the reading landed on a multiple, 10 or 15, 103 times in 156 at 8 draws from 32 e up, and at 32 e read 5 on 8 of 32 and a multiple on 20 (\textbf{measured}, \texttt{period\_\allowbreak{}power}). The multiples' peaks are as high as P's, so the noise picked among them.
\item \textbf{The global shuffle (superseded, Doug, 23 September): it read a period where another axis carried one.} One Feistel shuffle of the whole index per draw served every axis. With the plant at 1,024 e along x, the unplanted z and y read a period on 23 of 64 at 8 draws (36\%, against 1/9) and on 10 of 64 at 19 draws (15.6\%, against 1/20) (\textbf{measured}, \texttt{period\_\allowbreak{}power}). The explanation, not measured (the agreement levels were not printed): the global shuffle destroyed x's structure too, so the null's agreement fell below the live lattice's at the lags of z and y, and the exchangeability the false-period rate rests on failed there.
\item \textbf{The margin rule (superseded, Doug, 23 September).} For a candidate p with 2p ≤ L, take the lags that are not multiples of p: B of them, with A = Σ a(ℓ) over them (the total less the multiples). The margin was m(p) = (a(p) + a(2p))/(2Q) − A/(B·Q) = (B·(a(p) + a(2p)) − 2A) / (2·B·Q), the largest positive margin the candidate. A smooth lattice agrees more at every short lag than at the long ones, so the rule read smoothness as a period: a smooth ramp with noise read P = 2 or 3 on every axis (the 34 × 960 × 960 timing run, \textbf{measured}). Under the drawn null it passed 21 of 192 random axes at 8 draws (10.9\%).
\item \textbf{Why chance could not do this.} The old floor was m ≥ Σc²/N². Σc² is unchanged when the values are rearranged, so a floor built from it cannot tell structure from a rearrangement of the same values, and uniform random data passed 63 of 192 axes (33\%). The drawn null keeps the histogram and Σc² exactly and changes only the arrangement, which is the thing a period is. Σc² is still reported, as \texttt{collisions}.
\item \textbf{The shuffle.} A keyed Fisher–Yates on each line along the axis, one thread a line: for step s from n − 1 down to 1, the element at s swaps with the one at h mod (s + 1), h a 32-bit hash of the line index and s under four keys (\texttt{PERIOD\_\allowbreak{}ROUNDS}). Each key is the content signum's words mixed with c · 4 + round, for the counter c = draw · rank + axis. The seed is the lattice's content, so the same content draws the same band. Superseded: a Feistel network on the whole index, 4 rounds, with cycle-walking to stay below N (the global shuffle).
\end{itemize}

\subsection{\texorpdfstring{A13. The record machine (M10): \texttt{keymath}, \texttt{key\_\allowbreak{}schedule}, \texttt{cycle}}{A13. The record machine (M10): keymath, key\_schedule, cycle}}

A record program is a list of steps s\_\allowbreak{}1, …, s\_\allowbreak{}n over registers of limbs, each step reading earlier registers and writing its own, run for every record in one sweep.

\begin{itemize}
\item \textbf{The steps are n} (Doug, 23 September: “increase the scheduler steps to n”; 24 September: stack the floors, “make it so”).

\begin{itemize}
\item The step count has no bound. \texttt{ENGINE\_\allowbreak{}RECORD\_\allowbreak{}STEPS\_\allowbreak{}MAX} is gone: it sized a per-thread array holding one sign per step, and each sign now sits beside its register in the file.
\item The binding resource is the file: \texttt{ENGINE\_\allowbreak{}RECORD\_\allowbreak{}LIMBS\_\allowbreak{}MOST} = 256 limbs of live registers.
\item A stack of floors runs as one program in one launch. Each floor is a round that reads the floor below and leaves its state, often wrapped, for the next.
\item The layout frees each register in a single pass, bucketed by its last reader. It frees the same registers at the same steps as the per-step scan it replaced.
\item The algebra of the stack, what it compresses and what it leaves, with the stacked-against-chained measurement (13 to 22 times, records equal): \hyperref[chap:vertical-time-compression]{vertical\_\allowbreak{}time\_\allowbreak{}compression.md}.
\item \textbf{Proved:} 700 floors of a round over four 32-bit words (xor, and, sum, xor with a constant, a 32-bit wrap), 4,204 steps, run in an 8-limb file. On 1,024 lanes the device equals the host word for word, and both equal the same rounds on the CPU's own 64-bit two's complement at every tapped floor (\texttt{test/\allowbreak{}record\_\allowbreak{}bitwise\_\allowbreak{}test}).
\end{itemize}
\item \textbf{Xor, and, and the two's complement wrap} (Doug, 24 September).

\begin{itemize}
\item \texttt{ENGINE\_\allowbreak{}RECORD\_\allowbreak{}XOR} and \texttt{\_\allowbreak{}AND} work bitwise on each operand's two's complement, sign-extended without end, one bit wider than the wider operand. The imprint narrows that where it can show a register is never negative. An and with such a register is no wider than it, and an xor of two such is no wider than the wider. The result's sign is read from the operands' signs.
\item \texttt{ENGINE\_\allowbreak{}RECORD\_\allowbreak{}WRAP} takes left modulo 2\textasciicircum{}w, read back signed, with w ≥ 4 in \texttt{right}.
\item \textbf{Proved:}

\begin{itemize}
\item Every output bit is rebuilt from the raw fields one bit at a time, and each sign still holds 64 bits past its written width. This holds on 4,096 lanes in the 64-limb file and 512 in the 256-limb file.
\item Hand-worked answers match.
\item The narrowed widths are exactly the rule's: 32, 20, 20, 32, 41, 41, 32 and 32 over unsigned 32- and 20-bit fields and a signed 40-bit field. All 4,096 lanes equal the oracle past every narrowed width.
\item A wrap below 4 bits is refused at imprint (\texttt{test/\allowbreak{}record\_\allowbreak{}bitwise\_\allowbreak{}test}, 24 checks, 0 failed).
\end{itemize}
\end{itemize}
\item \textbf{Register reuse} (opt-in, 23 September). Register r is live from the step that writes it to its last reader. With reuse on, a liveness free-list returns r to the file once its last reader has run, so a chain needs only the registers live at once, not one per step. \textbf{Proved} on a 21-step additive chain: the same output from 3 limbs as from 21. Off by default, so a program that does not ask keeps one register per step: each register placed at the running total of the limbs before it, the same bump as before, so a program with no tables keeps the same layout as before, by construction (not measured against the old binary).
\item \textbf{The lookup table} (Doug, 23 September: “implement the lut”). A table step reads the low b bits of one register's magnitude, b ≤ 32, and returns T[|x| mod 2\textasciicircum{}b], an \texttt{out\_\allowbreak{}bits}-wide value, from 2\textasciicircum{}b rows of ⌈out\_\allowbreak{}bits/32⌉ limbs. The file holds a magnitude and a sign: a negative x reads the row of |x|, and x and −x read the same row. A register that can be negative is made never negative before it indexes a table. It is a general one-variable function on the lane's alphabet: 2\textasciicircum{}16 entries for a u16 lane, built once by pushing every value through. \textbf{Proved:} filled by running the ordinary operations over the whole 16-bit alphabet, the table equals them on 4,096 lanes, on the device and the host (M10).
\item \textbf{Composition keeps the order.} Reading g through f gives the table (g∘f)[x] = g[f[x]], one table for the chain: the floor is 2-D, but the operations need not be, for a chain of them “gets collapsed in order onto it” (Doug, 23 September). Composition regroups, (h∘g)∘f = h∘(g∘f), but does not reorder: g∘f ≠ f∘g. \textbf{Proved} with f(x) = 65535 − x and g(y) = |y − 30000|. Regrouping: g read through f equals the ops for |35535 − x|, and one table composed on the host equals the two. No reordering: f∘g differs on some lane.
\item \textbf{The two towers} (Doug, 24 September: a second tower over the first one's boundary, inverted, and the two collapse together).

\begin{itemize}
\item Floor division is a record floor: ⌊v / 2\textasciicircum{}k⌋ = EXACT\_\allowbreak{}QUOTIENT(v − AND(v, 2\textasciicircum{}k − 1), 2\textasciicircum{}k), rounding toward −∞ for every integer v, as \texttt{tower\_\allowbreak{}floor\_\allowbreak{}shift} does.
\item With it, the crystal's 5/3 lifting T (A14) and its inverse are sums, differences, constants and that division: stacks of record floors. An operation program F over the crystal and T⁻¹ compose into one stack F ∘ T⁻¹ by the regrouping law.
\item \textbf{Proved:} one 5/3 level over 8 signed 16-bit samples and its inverse, 129 steps in one program and a 12-limb file. On 4,096 lanes the device equals the host word for word, the forward floor equals tower.cu's formulas, and the inverse returns every sample exactly (\texttt{test/\allowbreak{}record\_\allowbreak{}bitwise\_\allowbreak{}test}).
\item Derived in \hyperref[chap:vertical-time-compression]{vertical\_\allowbreak{}time\_\allowbreak{}compression.md}:

\begin{itemize}
\item The file holds live registers only. A field counts against it only while its register is live; the atom stays in the record.
\item keymath folds no constant, drops no dead step and merges no duplicate step.
\item A record is addressable to 2\textasciicircum{}32 bits for each member; past that, device memory bounds it.
\item A lifting cone spans (2\textasciicircum{}\{L+2\} − 3)\textasciicircum{}D samples, yet one level over a 4-D cone of 625 samples runs in a 15-limb file.
\item keymath's widths grow every register 1 bit a level, the first level 2, where the lows' values grow 0.585 bits and the highs' 1 (from 25 September, the linear forms of A16; E4). Before, a low grew 4 bits a level, and 7 before the constant-divisor narrowing (the narrowing is \textbf{proved}, \texttt{test/\allowbreak{}record\_\allowbreak{}divide\_\allowbreak{}test}, 19 checks, 0 failed).
\item T is a bijection; an operation G between the towers runs as T⁻¹GT, and the levels above G cancel past it.
\item The heap (the bits a floor's values occupy) mirrors exactly across the crystal. The ring (the imprint's widths) does not, until a wrap at the mirror makes it symmetric to one bit a register. The crystal with the ring's framing is a code length bounding K(x) from above, and by counting no lens shortens most lanes.
\end{itemize}
\item Open: the neighbor gather in D axes (the halo in the record, or one axis per sweep), F ∘ T⁻¹'s error in D axes, which operations commute with T, the heap and the wrap at the mirror as committed tests (both measured in scratch runs; \textbf{proved} since by \texttt{test/\allowbreak{}record\_\allowbreak{}boundary\_\allowbreak{}test}), the oval (Doug's call), and the whole crystal as one stack (\hyperref[chap:vertical-time-compression]{vertical\_\allowbreak{}time\_\allowbreak{}compression.md}).
\item \textbf{Where the two towers stand in the engine} (24 September; Doug: “You didn't build the inverse tower into the engine yet”).

\begin{itemize}
\item In the engine, T and T⁻¹ are \texttt{tower.\allowbreak{}cu}'s own kernels (\texttt{tower\_\allowbreak{}forward\_\allowbreak{}kernel} and \texttt{tower\_\allowbreak{}inverse\_\allowbreak{}kernel}, behind \texttt{tower\_\allowbreak{}lift} and \texttt{tower\_\allowbreak{}lower}). The crystal path (\texttt{engine.\allowbreak{}cu}) and \texttt{root\_\allowbreak{}universal} call them. They are not record programs there.
\item As record floors, T and T⁻¹ exist only in the tests: \texttt{test/\allowbreak{}record\_\allowbreak{}bitwise\_\allowbreak{}test} (one level over 8 samples) and \texttt{test/\allowbreak{}record\_\allowbreak{}boundary\_\allowbreak{}test} (64 samples at 4 levels). No engine module emits them.
\item No engine path composes an operation with T⁻¹ in one record program. F ∘ T⁻¹ and T⁻¹GT are proved on those test programs and derived beyond them, and not built.
\item The record programs the engine runs today (the cell program's, Ω's reduct in \texttt{chaitin\_\allowbreak{}omega}, π's BBP terms in \texttt{pi\_\allowbreak{}tower}) read their own fields and use neither tower.
\item \textbf{Built:} the towers as kernels. \textbf{Not built:} the towers as record floors in the engine, and F ∘ T⁻¹ as the path a program runs on the crystal.
\end{itemize}
\end{itemize}
\item \textbf{The two crystals} (Doug, 24 September: the finite towers and their limit, with an infinite delta between them).

\begin{itemize}
\item A register is a 2-adic integer, and a wrap to w is the projection ℤ₂ → ℤ/2\textasciicircum{}w. Sum, difference, product, xor and and commute with every projection. An exact quotient by an odd c is the product by c⁻¹ in ℤ₂.
\item \textbf{Proved} (\texttt{test/\allowbreak{}record\_\allowbreak{}coherence\_\allowbreak{}test}, 14 checks, 0 failed): 393,216 of 393,216 lane-widths agree three ways over 16 random programs, 73,728 of 73,728 for the odd divisors, and a quotient and a comparison break the agreement. Mod 3\textasciicircum{}5, 3\textasciicircum{}10, 5\textasciicircum{}4 and 7\textasciicircum{}3, 262,144 of 262,144 lane-moduli agree three ways over 16 ring-only programs, the CRT join with the 8-bit window equals the exact run on all of them, and xor breaks mod 3.
\item \textbf{Proved} (\texttt{test/\allowbreak{}record\_\allowbreak{}boundary\_\allowbreak{}test}, 48 checks, 0 failed, 24b2785): the crystal of 64 samples at 4 levels as a boundary. T's reach is 3ℓ at the level-ℓ lows and 3ℓ − 2 at its highs and T⁻¹'s is L + 2, each met on the device; arbitrary crystals return through T⁻¹ then T; constants and 2\textasciicircum{}\{3L\}ℤ\textasciicircum{}n moves pass through T, negation and doubling do not; the heap mirrors, the ring is ring\_\allowbreak{}0 + 6(n − n/2\textasciicircum{}ℓ); the quotient toward zero keeps order and a sum within one; T and T⁻¹ keep Haar measure, counted; det M = ±1 exactly, its sign +1 in the run. T against 8 null shuffles a lane identifies each structured class past the null's 1/(d + 1) bound and noise within it, and the crystal is a one-to-one ID: a draw changes it exactly when it moves a value, and every single flipped bit changes the image.
\item Derived in \hyperref[chap:two-crystals]{two\_\allowbreak{}crystals.md}: what does not factor through the projection, the lifting as a homeomorphism of ℤ₂\textasciicircum{}n whose level-L lows read the samples to w + 3L bits, the windows' union ℤ and the projections' limit ℤ₂ with the solenoid between them, what passes to a limit, the top projection's limit ℝ, the odd crystals and the places of ℚ, the count as coset counting, and Doug's posits bounded (the anchors' complexity open; the fingerprint of T against null permutations, proved; the knf's identity by spatial null permutation, built; the pairwise agreement between sections, open).
\item \textbf{Proved} (\texttt{knf\_\allowbreak{}identity}, 23 checks, 0 failed, e4eae72): the knf of the nbody lattice (64³, 177 frames) against spatial null permutations. The host's walk equals the engine's knf at every voxel, and a single flipped bit leaves it unchanged exactly where the window rule says (1,500 of 8,192): the knf is many-to-one. The 48 exact cube motions carry the knf and keep its entangled entropy E and its cloud. The whole is its tiles plus its seams at every tile size. Alike-voxel controls are identified within 5σ of 1/20. \textbf{Measured:} E is 44.6 times the free null's mean, the inside and between departure curves cross between tiles of 2 and 4 and sum to 0.590 at 4, and each body's box has its own curve. Written in \hyperref[chap:two-crystals]{two\_\allowbreak{}crystals.md}.
\end{itemize}
\item \textbf{The full Kolmogorov–Arnold form.} With linear steps only, the machine held the theorem's skeleton with linear edges; the table gives it general one-variable edges, which is item 1 of \hyperref[chap:kolmogorov-arnold]{kolmogorov\_\allowbreak{}arnold.md} and the LUT of \hyperref[chap:noise-sieve-tower]{noise\_\allowbreak{}sieve\_\allowbreak{}tower.md} §2. Which functions to tabulate, without fitting, is still open (item 3 there).
\end{itemize}

\subsection{\texorpdfstring{A14. The root universal (M20): \texttt{tower}}{A14. The root universal (M20): tower}}

The tower and the table are separate objects; laid together they are one (Doug, 23 September). Let W\_\allowbreak{}k be the approximation on floor k, L\_\allowbreak{}k level k's lifting, and E\_\allowbreak{}k the edges laid on floor k.

\begin{itemize}
\item \textbf{The edge is a bijection.} E(w) = (w with its low b bits x replaced by π(x)), π a permutation of [0, 2\textasciicircum{}b), the high bits passed through. So E is a bijection on the coefficient word and E⁻¹ uses π⁻¹ (proved by construction; the round trips below are the check).
\item \textbf{The order.} Lift is EN ∘ L(N−1) ∘ … ∘ E1 ∘ L0 ∘ E0, and lower is its inverse, E0⁻¹ ∘ L0⁻¹ ∘ … ∘ EN⁻¹. The edges on one floor are a stack, undone last laid first. \textbf{Proved:} every edge program in \texttt{tower\_\allowbreak{}edge\_\allowbreak{}test} and \texttt{root\_\allowbreak{}universal} rebuilds the exact lanes, 171 random edges on 48 volumes through the whole crystal path.
\item \textbf{The fold.} Two edges on one floor are one edge: π\_\allowbreak{}b ∘ π\_\allowbreak{}a, laid as one table, lays the same crystal. It keeps the order, since π\_\allowbreak{}a ∘ π\_\allowbreak{}b lays another, and a lifting level between two edges blocks the fold. \textbf{Proved} on 8 volumes each.
\item \textbf{What it is} (theory, Doug's framing). The tower is the Kolmogorov–Arnold form's additive skeleton, the table its one-variable edge, and with edges between the floors the collapse is an outer sum with nonlinear edges.
\item \textbf{Its two bounds} (theory; kept). It is universal over the finite coefficient alphabet: every function of the indexed bits is one table, exactly, not approximately. It is not the continuous-function theorem of Kolmogorov and Arnold. And inside the reversible stream every edge is a permutation, so the second bound is one of width, not of kind (below).
\item \textbf{The dimensional expansion} (Doug, 23 September: “reaching a boundary here means we need a dimensional expansion to realize the transform”). A function f that is not a bijection (|x|, a comparison) rides the stream through the Bennett embedding, (x, y) → (x, y ⊕ f(x)): a permutation of the pair for any f, since y comes back as the carrier ⊕ f(x) (proved by construction). The price is width: a permutation takes b\_\allowbreak{}x bits, any f takes b\_\allowbreak{}x + b\_\allowbreak{}y, within the 20-bit limit. No engine change; the tables are built on the host. \textbf{Proved} by \texttt{tower\_\allowbreak{}edge\_\allowbreak{}test}, 25 checks, 0 failed. |x| of an 8-bit two's complement field and x ≥ 100, each a 16-bit Bennett edge (8 bits of x, 8 of carrier) on the collapsed floor of 1 × 8 × 8 × 8, alternate over 64 lattices: they round-trip on 64 of 64; f(x) reads out of the carrier (x kept, the carrier moved by exactly f(x), the high bits passed) on 64 of 64; the edge touches only its one coefficient on 64 of 64. The two edges on floors 1 and 2, where lifting mixes their outputs, round-trip. |x| laid bare on 16 bits, not embedded, is refused, since x and −x collide. The quantum forms of the same expansion are standard theorems (theory here): Naimark (an ask is a projective measurement in a larger space), Stinespring (a transform is a unitary in a larger space, then a discard) and purification (a mixed state is part of a pure state in dimension d²); the ask is A15.
\item \textbf{The price} (measured, M20). An unfitted edge is exact but not free. On floor 0, where every coefficient is, a random 12-bit permutation raises the coded bits from 6.802 to 12.608 a voxel; on the collapsed floor, one coefficient, the cost does not show at three decimals at any width. The reading, not measured: a random permutation of the low bits sends small magnitudes to large ones, and the magnitude coder pays for each, so the cost scales with the coefficients the edge touches. Whether a fitted or data-named edge lowers the bits, not only raises them, is not measured.
\end{itemize}

\subsection{\texorpdfstring{A15. The ask and the state (M21): \texttt{ask\_\allowbreak{}state}}{A15. The ask and the state (M21): ask\_state}}

The terms are M21's (Doug, 23 September). The idea is standard: QBism (Fuchs and Schack) and the SIC-POVMs.

\begin{itemize}
\item \textbf{The ask returns the state.} For a qubit and E\_\allowbreak{}i = (I + v\_\allowbreak{}i·σ)/4, the answers are p\_\allowbreak{}i = (1 + r·v\_\allowbreak{}i)/4. With Σ v\_\allowbreak{}i = 0 and Σ v\_\allowbreak{}i v\_\allowbreak{}iᵀ = c·I, Σ p\_\allowbreak{}i v\_\allowbreak{}i = (c/4) r, so r = (4/c) Σ p\_\allowbreak{}i v\_\allowbreak{}i and the state is read back from the answers alone (theory; \textbf{proved} on 25 states for each ask, c = 1 and c = 4/3). A complete ask needs at least d² outcomes, here 4 (theory, standard).
\item \textbf{The phase falls out of the ask.} The equator states differ only in phase: each answers ½, ½ to the 0/1 ask, which sees magnitudes only. The complete ask tells all 91 pairs of the 14 apart (\textbf{proved}). So the phase is not lost; one basis's ask cannot see it, and a complete ask holds it.
\item \textbf{The crossing rule.} The chance of + on axis a is (1 + r\_\allowbreak{}a)/2 = Σ p\_\allowbreak{}i (½ + (2/c) v\_\allowbreak{}\{i,a\}), a fixed linear map from the complete ask's answers (theory; \textbf{proved}, 75 of 75 for each ask). Its weights onto +x are 3/2, 3/2, −1/2, −1/2 for the rational ask and ½ ± ½√3 for the SIC. A negative weight is what classical total probability never has, and the difference is interference.
\item \textbf{The valid set.} A distribution is a state exactly when |r|² ≤ 1. (1, 0, 0, 0) gives |r|² = 12 for the rational ask and 9 for the SIC, so it is refused, and it crosses onto +x at 3/2 and ½ + ½√3, no probabilities (\textbf{proved}). On the 455 distributions with denominator 12 the valid set is 55 (rational) and 87 (SIC), and 104 and 72 more pass the three axis crossings without being states (\textbf{measured}): the valid set is stricter than the axis crossings.
\item \textbf{Exact numbers.} Every value is exact: rationals on the exact integer, and the SIC's √3 carried by its relation, (a + b√3)(c + d√3) = (ac + 3bd) + (ad + bc)√3, with the sign of a + b√3 read from a² against 3b², never equal unless both are 0, since √3 is irrational (proved by construction). The SIC's answers cross into the rational ask's exactly, 25 of 25 (\textbf{proved}).
\item \textbf{The dimensional expansion} (theory, standard theorems). Naimark: an ask is a projective measurement in a larger space. Stinespring: a transform is a unitary in a larger space, then a discard. Purification: a mixed state is part of a pure state in dimension d². The engine's form is the Bennett embedding of A14, proved.
\item \textbf{Doug's statements} (his; the first built in \texttt{ka\_\allowbreak{}psi}, the rest theory). “The exp is symbolic after the second nth exp”: ψ's terms γ\textasciicircum{}−β(r), with β(r) = (nʳ − 1)/(n − 1), are carried as a digit and an exact integer exponent, a sparse form whose cost is polynomial in the scale, not the expanded digits. \texttt{ka\_\allowbreak{}psi} carries ψ this way to β(60) = 1,152,921,504,606,846,975 (built; 152 checks, 0 failed). ψ is the inner function of Sprecher's version of the superposition theorem, as J. Braun and M. Griebel state it (“On a constructive proof of Kolmogorov's superposition theorem”, Constructive Approximation 30(3) (2009) 653–675, doi:10.1007/s00365-009-9054-2, the record checked at Crossref; the preprint read, pages 1 to 5). Their Theorem 2.1 takes integers n ≥ 2, m ≥ 2n, γ ≥ m + 2 and a = [γ(γ − 1)]⁻¹, and writes f(x) = Σ\_\allowbreak{}\{q=0\}\textasciicircum{}\{m\} Φ\_\allowbreak{}q ∘ ξ(x\_\allowbreak{}q), with ξ(x\_\allowbreak{}q) = Σ\_\allowbreak{}\{p=1\}\textasciicircum{}\{n\} α\_\allowbreak{}p ψ(x\_\allowbreak{}p + qa), α\_\allowbreak{}1 = 1 and α\_\allowbreak{}p = Σ\_\allowbreak{}\{r≥1\} γ\textasciicircum{}−(p−1)β(r). ψ is defined on the terminating base-γ rationals D\_\allowbreak{}k = \{ Σ\_\allowbreak{}\{r≤k\} i\_\allowbreak{}r γ\textasciicircum{}−r \} (their 2.2), and in Köppen's form (their 2.7) level k adds its digit at the weight γ\textasciicircum{}−β(k). So γ\textasciicircum{}−β(r) is the γ-adic weight of level r. What ψ is, strictly increasing and continuous in Köppen's form and not so in Sprecher's, is proved exactly by the sim \texttt{ka\_\allowbreak{}psi} (M19), to depth 60 with the exponent never expanded, in “Kolmogorov's inner function, exactly” of \hyperref[chap:kolmogorov-arnold]{kolmogorov\_\allowbreak{}arnold.md}. Its “Separation” part finds the paper's Lemma 3.4 false at k = 1 and restores Lemmas 3.7 and 3.8 at every k with a corrected bound and a ramp narrowed by γ². The outer functions' contraction is not built. “If we can represent a two-sqrt product as two ints and their words of math operations then that is countable”: the numbers named by finitely many integers and a finite word of operations (the algebraic, the computable, π as its program) are countable, and the exact storage holds each of them. And “Pi compresses to 1kb”: π as a spigot program.
\end{itemize}

\subsection{\texorpdfstring{A16. The Gaussian step (M10, M19): \texttt{test/\allowbreak{}record\_\allowbreak{}gaussian\_\allowbreak{}test}}{A16. The Gaussian step (M10, M19): test/record\_gaussian\_test}}

BBP's base is a power of a Gaussian prime (Doug, 24 September: “The shape of the algorithm itself is dictating the geometric ruleset we must enforce”, then “if the geometry looks absurd we adjust our algorithm until it's fluid”). Each rule carries one label: \textbf{cited} (a published or standard theorem), \textbf{derived} (exact algebra worked here and checked in floating point by a scratch script, the ledger's 24 September; theory, not built), \textbf{proved} (the run named) or \textbf{open} (not known).

\begin{itemize}
\item \textbf{The ring and its prime} (cited, standard). ℤ[i] = \{a + bi : a, b ∈ ℤ\} is Euclidean under the norm N(a + bi) = a² + b². 2 = −i(1 + i)², so 1 + i is the one prime over 2, ramified, with N(1 + i) = 2. The completion there is ℤ₂[i], with 1 + i its uniformizer: v(1 + i) = 1 and v(2) = 2.
\item \textbf{Eight steps are sixteen} (derived; proved below). (1 + i)² = 2i, so (1 + i)\textasciicircum{}8 = (2i)\textasciicircum{}4 = 16. The powers for k = 1 to 8 are 1 + i, 2i, −2 + 2i, −4, −4 − 4i, −8i, 8 − 8i and 16.
\item \textbf{The step turns and scales} (derived). 1 + i = √2·e\textasciicircum{}\{iπ/4\}, so z ↦ z(1 + i) turns z by π/4 and scales it by √2, and N doubles exactly. Eight steps are a whole turn and a scale of 16 = 2\textasciicircum{}4. So BBP's digit shift, ×16, one hex digit, is eight Gaussian steps.
\item \textbf{The record floor} (derived; proved below). On z = a + bi the step is (a, b) ↦ (a − b, a + b): one DIFFERENCE and one SUM, the matrix [[1, −1], [1, 1]], determinant 2.

\begin{itemize}
\item Its image is the pairs of equal parity, which is the ideal (1 + i), of index 2. So each floor drops one bit, and after k floors the image is (1 + i)\textasciicircum{}k ℤ[i], of index 2\textasciicircum{}k. v rises by 1 a floor and N doubles: the bit leaves the archimedean side for the (1 + i)-adic one.
\item The inverse floor is ((c + d)/2, (d − c)/2): two EXACT\_\allowbreak{}QUOTIENTs by 2, exact on the image. A pair of mixed parity is exactly a lane the machine refuses (derived, not built).
\item SUM and DIFFERENCE commute with every wrap (E1, proved by \texttt{test/\allowbreak{}record\_\allowbreak{}coherence\_\allowbreak{}test}), so the floor runs exactly in any window. In a window of w bits each floor keeps an index-2 part of the pairs, and 2w floors send every pair to 0, since (1 + i)\textasciicircum{}\{2w\} = (2i)\textasciicircum{}w ≡ 0 mod 2\textasciicircum{}w (derived).
\end{itemize}
\item \textbf{The product formula over ℚ(i)} (the formula cited, standard; the values derived). Take |x|\_\allowbreak{}ℂ = x·x̄ and |x|\_\allowbreak{}\{1+i\} = 2\textasciicircum{}\{−v(x)\}; every other place gives 1 on these values. Then |16|\_\allowbreak{}ℂ·|16|\_\allowbreak{}\{1+i\} = 256·2\textasciicircum{}\{−8\} = 1. A floor multiplies |·|\_\allowbreak{}ℂ by 2 and |·|\_\allowbreak{}\{1+i\} by ½, so the product holds at 1 on every floor: the floor moves a bit from one place to the other and loses none.
\item \textbf{π from the step} (the series cited; the rest derived).

\begin{itemize}
\item BBP (Bailey, Borwein and Plouffe, Math. Comp. 66, 1997): π = Σ\_\allowbreak{}k 16\textasciicircum{}\{−k\}(4/(8k + 1) − 2/(8k + 4) − 1/(8k + 5) − 1/(8k + 6)).
\item Expanding 1/(1 − x⁸) term by term gives Σ\_\allowbreak{}k 16\textasciicircum{}\{−k\}/(8k + j) = 2\textasciicircum{}\{j/2\} ∫\_\allowbreak{}0\textasciicircum{}\{1/√2\} x\textasciicircum{}\{j−1\}/(1 − x⁸) dx. So π = ∫\_\allowbreak{}0\textasciicircum{}\{1/√2\} (4√2 − 8x³ − 4√2x⁴ − 8x⁵)/(1 − x⁸) dx, the integral their proof runs through. The upper limit is 1/√2 = |1/(1 + i)|, and (1/√2)\textasciicircum{}8 = 1/16.
\item \textbf{Adjusted: four of the eight roots of unity are poles, not all eight.} 1 − x⁸ vanishes at the eighth roots of unity, e\textasciicircum{}\{ikπ/4\} = ((1 + i)/√2)\textasciicircum{}k. The ruleset as first stated put a pole at each. The numerator is −8(1 + x²)(x − 1/√2)(x² + √2x + 1), which cancels ±i and e\textasciicircum{}\{±3iπ/4\}. What is left is 8(1/√2 − x)/((1 − x²)(x² − √2x + 1)), with poles at ±1 and e\textasciicircum{}\{±iπ/4\} = (1 ± i)/√2 only, and it vanishes at the upper limit.
\item With y = √2x, the form BBP's proof reaches: π = ∫\_\allowbreak{}0\textasciicircum{}1 16(y − 1)/((y² − 2)(y² − 2y + 2)) dy. y² − 2y + 2 = (y − ρ)(y − ρ̄) with ρ = 1 + i, so the poles are the Gaussian prime, its conjugate, and ±√2 = ±|ρ|. By partial fractions the integrand is 2/(y − √2) + 2/(y + √2) − 2ρ/(y − ρ) − 2ρ̄/(y − ρ̄). The residue at the prime is −2ρ, and the four residues sum to 0.
\item Over [0, 1] the poles at ±√2 give −2 ln 2 and the real part at ρ gives +2 ln 2, so they cancel. What is left is 4·arctan 1 = 4·(π/4), and π/4 is the angle [0, 1] subtends at ρ, arg(−i) − arg(−1 − i). So π is four times one Gaussian step's turn, and BBP's logarithms, the step's scale (ln N(ρ) = ln 2), cancel exactly: the series keeps the turn and drops the scale.
\end{itemize}
\item \textbf{The width, made fluid} (Doug, 25 September: “do the keymath rule”; built, and proved at every floor).

\begin{itemize}
\item Until 25 September keymath widened a SUM or DIFFERENCE to max + 1, so floor k carried w + k bits, 25 to 32 at w = 24.
\item Floor k's values are (a·x\_\allowbreak{}k − b·y\_\allowbreak{}k, a·y\_\allowbreak{}k + b·x\_\allowbreak{}k) with (1 + i)\textasciicircum{}k = x\_\allowbreak{}k + i·y\_\allowbreak{}k, and |x\_\allowbreak{}k| + |y\_\allowbreak{}k| = 2\textasciicircum{}\{⌈k/2⌉\}. So they grow half a bit a floor, against the old rule's whole bit. The 24 September script's corner figures (25, 26, 27, 27, 28, 28, 28 and 28) are two's complement widths, sign included (checked by hand at floors 2, 3, 7 and 8). The register holds the magnitude with its sign beside it, and there floor k needs 24 + ⌈k/2⌉.
\item The cause: the old rule bounded each register alone. An odd floor's values fill a turned square, |c| + |d| ≤ 2\textasciicircum{}\{⌈k/2⌉\}·2\textasciicircum{}\{w−1\}, and a per-register box keeps its bounding square, which the next floor turns into twice the reach. Over 8n floors the imprint grew 8n bits and the values 4n.
\item The rule built (\texttt{keymath.\allowbreak{}cu}): every register is a linear form over atoms, integer coefficients plus a constant. A sum or difference adds its operands' forms, a product by a constant scales the other's, a constant is its value, and a wrap that passes its register through keeps that register's form. Every other register is an atom. The width is the fewer of the operation's own rule and the form's bound, |x| ≤ |c| + Σ |c\_\allowbreak{}i|(2\textasciicircum{}\{b\_\allowbreak{}i\} − 1). A coefficient past 2\textasciicircum{}62 makes the register an atom. On the Gaussian floor the coefficients are the parts of (1 + i)\textasciicircum{}k, so the bound is 2\textasciicircum{}\{⌈k/2⌉\}·2\textasciicircum{}w and the width is w + ⌈k/2⌉.
\item Read through Mathai and Thiang (arXiv:1505.05250): T-duality turns the bulk-boundary map ∂ into a restriction ι*. Here the width, a boundary quantity, is read by restricting the bulk form to its atoms' widths, not carried register by register (an analogy, not a theorem about keymath).
\item The wrap at w + ⌈k/2⌉ + 1 bits a floor is not needed.
\end{itemize}
\item \textbf{Proved} (\texttt{test/\allowbreak{}record\_\allowbreak{}gaussian\_\allowbreak{}test}, 11 checks, 0 failed, \texttt{build/\allowbreak{}20260925\_\allowbreak{}012737\_\allowbreak{}record\_\allowbreak{}gaussian\_\allowbreak{}test}, a job on tessera; the first run, under the old rule, was 8 checks, 0 failed, \texttt{build/\allowbreak{}20260924\_\allowbreak{}222402\_\allowbreak{}record\_\allowbreak{}gaussian\_\allowbreak{}test}). The program is two signed 24-bit fields and eight floors of one DIFFERENCE and one SUM, 18 steps with every floor an output. It ran on 4,096 lanes: every pair of the five edge values (0, −1, 1, −2\textasciicircum{}23 and 2\textasciicircum{}23 − 1), and keyed draws for the rest. Its checks:

\begin{itemize}
\item the program imprints, lays and loads, and floor k is 24 + ⌈k/2⌉ bits wide: 25, 25, 26, 26, 27, 27, 28 and 28;
\item some lane fills every floor's width, |value| ≥ 2\textasciicircum{}\{bits − 1\}, so no width is loose;
\item the device equals the host word for word;
\item every floor equals z(1 + i)\textasciicircum{}k, taken on the CPU from the powers above;
\item floor 4 is −4z and floor 8 is 16z on every lane;
\item each floor's two registers share their parity;
\item floor 8's values lie in [−2\textasciicircum{}27, 2\textasciicircum{}27), 28 bits, and reach −2\textasciicircum{}27;
\item tessera admits the job and it releases.
\end{itemize}
\item \textbf{The rule across the other record tests} (25 September, each a job on tessera, run one at a time): guide 12, table 16, divide 20, bitwise 25, coherence 15 and order 18 checks, 0 failed; boundary 49 checks, 1 failed. The failure is its ring law ring\_\allowbreak{}ℓ = ring\_\allowbreak{}0 + 6n(1 − 2\textasciicircum{}\{−ℓ\}), which the old rule's widths made; the ring measured is in E4's widths. With the law set to ring\_\allowbreak{}0 + n + 2n(1 − 2\textasciicircum{}\{−ℓ\}) (Doug, 25 September: “yes”), boundary is 49 checks, 0 failed (\texttt{build/\allowbreak{}20260925\_\allowbreak{}020556\_\allowbreak{}record\_\allowbreak{}boundary\_\allowbreak{}test}).
\item \textbf{Where it stands in the engine.} The width rule is in \texttt{keymath.\allowbreak{}cu}, so every record program takes it. The Gaussian step itself is built only in the test. \texttt{pi\_\allowbreak{}tower} runs 16\textasciicircum{}\{d−i\} mod m a nibble at a time (E13) and does not use the step.
\item \textbf{The bridge} (open).

\begin{itemize}
\item The digit at d is \{16\textasciicircum{}d π\} = \{4S₁ − 2S₄ − S₅ − S₆\}, where S\_\allowbreak{}j = Σ\_\allowbreak{}\{i≤d\} (16\textasciicircum{}\{d−i\} mod m\_\allowbreak{}i)/m\_\allowbreak{}i plus a tail, with m\_\allowbreak{}i = 8i + j.
\item Each term's power reduces by Carmichael's λ (cited, standard): on m\_\allowbreak{}i's odd part, the exponent is taken mod λ; on its power of 2, the term is 0 once 4(d − i) reaches that power (derived). Still, d + 1 terms remain.
\item The Gaussian form does not shorten them. 16 is a rational integer, so (1 + i)\textasciicircum{}\{8(d−i)\} mod m\_\allowbreak{}i in ℤ[i]/(m\_\allowbreak{}i) is 16\textasciicircum{}\{d−i\} mod m\_\allowbreak{}i (derived).
\item What would close it is a sum over i taken at once: a modular form or Hecke eigenvalue whose coefficient is Σ\_\allowbreak{}i r\_\allowbreak{}i/m\_\allowbreak{}i mod 1. The Hecke recursion, T\_\allowbreak{}\{p\textasciicircum{}\{n+1\}\} = T\_\allowbreak{}p·T\_\allowbreak{}\{p\textasciicircum{}n\} − p\textasciicircum{}\{k−1\}·T\_\allowbreak{}\{p\textasciicircum{}\{n−1\}\} on level-1 forms of weight k (cited, standard), makes a coefficient at 2\textasciicircum{}\{4d\} a power of a 2 × 2 matrix, about log d steps (derived).
\item No form whose coefficients give BBP's sum is known. No method is known to us that finds π's hex digit at d in fewer than about d steps (open).
\end{itemize}
\end{itemize}

\section{\texorpdfstring{The scale audit: every number in the machine that fixes a scale}{The scale audit: every number in the machine that fixes a scale}}

A number in this table is a debt unless it is a word's width or a format's specification. Each is to come from the request or the data, or to be measured and reported, never tuned to a program.

{\footnotesize
\setlength{\tabcolsep}{6pt}
\begin{longtable}{>{\RaggedRight\arraybackslash}p{\dimexpr0.2582\linewidth-2\tabcolsep\relax}>{\RaggedRight\arraybackslash}p{\dimexpr0.2435\linewidth-2\tabcolsep\relax}>{\RaggedRight\arraybackslash}p{\dimexpr0.2133\linewidth-2\tabcolsep\relax}>{\RaggedRight\arraybackslash}p{\dimexpr0.2851\linewidth-2\tabcolsep\relax}}
\toprule
number & where & what it fixes & why it is a debt \\
\midrule
\endhead
the even-order guard & \texttt{residual.\allowbreak{}cu:24}, \texttt{keymath.\allowbreak{}cu:118}, \texttt{unit\_\allowbreak{}sweep.\allowbreak{}cu:254} & refuses every odd order & stricter than the algebra (A1); ruling (a) replaces it \\
\texttt{ENGINE\_\allowbreak{}RESIDUAL\_\allowbreak{}LIMBS} 9 & engine\_\allowbreak{}config.h & the residual's width: 16 input bits plus the orders' growth & sized for one program's orders; the width should follow the request's orders (each step adds a bit). The planes path derives it (input bits + Σ orders + 1) and returns it, since 25 September (M2); the 16-bit path and the key still use 9 \\
\texttt{UNIT\_\allowbreak{}SWEEP\_\allowbreak{}INPUT\_\allowbreak{}BITS} 16, unsigned 8 and 16 bit input only & unit\_\allowbreak{}sweep.cu, \texttt{engine\_\allowbreak{}source\_\allowbreak{}read} & the input's depth & a source deeper or signed needs a representation, not a refusal (the RSNA ruling: signed stored + 2\textasciicircum{}15 into the u16 lane, Doug to confirm). Since 25 September the unit sweep takes planes of any declared width (M2); its volume path and \texttt{engine\_\allowbreak{}source\_\allowbreak{}read} stay at 16 \\
\texttt{ENGINE\_\allowbreak{}AXES} 3 & engine\_\allowbreak{}config.h & three spatial axes & a word of the problem, but the residual, the tree and C are separable per axis \\
\texttt{ENGINE\_\allowbreak{}HISTORY\_\allowbreak{}WINDOW} 11, \texttt{ENGINE\_\allowbreak{}HISTORY\_\allowbreak{}WINDOWS\_\allowbreak{}MAX} 16 & engine\_\allowbreak{}config.h & the entropy window and how many fit & a window is a scale; it belongs in the request \\
\texttt{ENGINE\_\allowbreak{}COEFFICIENT\_\allowbreak{}LIMIT} 2\textasciicircum{}30 & engine\_\allowbreak{}config.h, read by tower & the largest lifted coefficient & a word's width; an overflow is a request error, not a silent clip \\
\texttt{COMPRESSION\_\allowbreak{}BLOCK} 64, 64 blocks a chunk, \texttt{COMPRESSION\_\allowbreak{}K\_\allowbreak{}BITS} 5, \texttt{COMPRESSION\_\allowbreak{}ESCAPE} 24 & compression.cu & the Rice coder's block and chunk & chosen; measure them against the data, or let the file carry them \\
\texttt{CRC\_\allowbreak{}SEGMENT\_\allowbreak{}PIXELS} 64 & crc.h & the CRC's parallel segment & no longer read by the tower (the seal replaced its CRC); crc.h goes with stage B \\
free / 3 · 2 & schedule.cu & the share of free device memory the schedule plans against & chosen (flagged 23 September by the RSNA session); ingest, prove and flatten consult no free memory at all, and nothing arbitrates the device between processes \\
extents ≤ 2\textasciicircum{}20 and lanes ≤ 2\textasciicircum{}40 & apxrep history read, \texttt{apxrep\_\allowbreak{}lanes}, \texttt{entry\_\allowbreak{}lanes} & the largest sample & chosen ceilings; the words hold more \\
\texttt{MARGINAL\_\allowbreak{}ARRANGEMENTS\_\allowbreak{}MOST} 2\textasciicircum{}16, \texttt{MARGINAL\_\allowbreak{}TARGETS\_\allowbreak{}MOST} 64, \texttt{MARGINAL\_\allowbreak{}LIMBS} 16 & marginal & the largest local set & OrganoidTracker's bound, not measured here; A5 says report the component's size instead \\
\texttt{SHIFT\_\allowbreak{}AGREEMENT\_\allowbreak{}LONGEST\_\allowbreak{}AXIS} 2\textasciicircum{}23 & shift\_\allowbreak{}agreement & the NTT's longest padded axis & the prime's transform length; a word's width, but a larger view needs a second prime, not a refusal \\
\texttt{ENGINE\_\allowbreak{}RECORD\_\allowbreak{}LIMBS\_\allowbreak{}MOST} 256, \texttt{CYCLE\_\allowbreak{}RECORD\_\allowbreak{}BLOCKS\_\allowbreak{}MOST} 4096, \texttt{ENGINE\_\allowbreak{}RECORD\_\allowbreak{}MEMBERS\_\allowbreak{}MAX} 3 & engine\_\allowbreak{}config.h, cycle.h & the record machine's reach; 256 limbs of live registers is the binding resource (A13) & not measured \\
\texttt{ENGINE\_\allowbreak{}RECORD\_\allowbreak{}WRAP\_\allowbreak{}BITS\_\allowbreak{}LEAST} 4 & engine\_\allowbreak{}config.h & the narrowest two's complement wrap, a nibble (A13) & a floor the request sets, not a resource; the step count \texttt{ENGINE\_\allowbreak{}RECORD\_\allowbreak{}STEPS\_\allowbreak{}MAX} it sat beside is gone, removed 24 September, since each sign now sits beside its register \\
\texttt{ENGINE\_\allowbreak{}RECORD\_\allowbreak{}TABLE\_\allowbreak{}INDEX\_\allowbreak{}BITS\_\allowbreak{}MOST} 32 & engine\_\allowbreak{}config.h & a table's index, the low bits of one register (A13) & a word's width: the index fits one limb \\
\texttt{TOWER\_\allowbreak{}EDGE\_\allowbreak{}INDEX\_\allowbreak{}BITS\_\allowbreak{}MOST} 20 & tower.h & a tower edge's widest index, 2\textasciicircum{}20 entries of 4 bytes, 4 MiB a table (A14) & not measured \\
\texttt{BODY\_\allowbreak{}OVERLAP\_\allowbreak{}CHUNK} 256, \texttt{CLIMB\_\allowbreak{}MACHINE\_\allowbreak{}BUNDLE} 32, \texttt{MAX\_\allowbreak{}TREE\_\allowbreak{}BUCKETS} 256, \texttt{UNIT\_\allowbreak{}SWEEP\_\allowbreak{}SHARED\_\allowbreak{}BYTES\_\allowbreak{}MOST} 48 KiB & kernels & device grains & the device's, not the data's; measure per device \\
\texttt{STACK\_\allowbreak{}UNCACHED\_\allowbreak{}BLOCK} 64 MiB, \texttt{NPY\_\allowbreak{}FORTRAN\_\allowbreak{}BLOCK} 16 MiB & stack, npy & IO grains & size by experiment to the disk's throughput (build plan 22); \texttt{NPY\_\allowbreak{}HEADER\_\allowbreak{}ROOM} and \texttt{NRRD\_\allowbreak{}HEADER\_\allowbreak{}ROOM} (1 MiB caps on a header) removed 23 September: the npy header is bounded by its own length field and the file, the nrrd header is read in doubling pieces until its blank line \\
flatten's slab, half the free device memory & flatten & frames per residual pass & measured 107.0 s against 9.0 s one frame at a time; a fraction of memory is not a scale of the problem \\
\texttt{SCRIPTURA\_\allowbreak{}DIRECT\_\allowbreak{}BYTES} 32 & scriptura.c & below it a memory call skips the arm table for portable & the portable short path's reach (four words), not the data's; measured only on this part \\
\texttt{SCRIPTURA\_\allowbreak{}AVX2\_\allowbreak{}STRING\_\allowbreak{}BYTES} 16 KiB & scriptura\_\allowbreak{}arm\_\allowbreak{}avx2.c & where copy and fill switch to \texttt{rep movsb}/\texttt{stosb} & the part's, not the data's: the measured crossover on a Haswell-E with ERMS (4 KiB lost 1.28×, 16 KiB and up at 1.00× the runtime); other parts move it, so measure per part \\
\texttt{PERIOD\_\allowbreak{}THREADS} 256, \texttt{PERIOD\_\allowbreak{}GRID\_\allowbreak{}ROWS\_\allowbreak{}MOST} 65,535 & period.cu & a block's width and the grid's row limit & the device's, not the data's: the columns are read from the device (multiprocessors × resident blocks), and the lags beyond the row limit are walked by stride \\
\texttt{PERIOD\_\allowbreak{}ROUNDS} 4 & period.cu & the keys of the line shuffle's hash in each draw of the null (A12; the Feistel shuffle's rounds before it) & a mixing depth, not a scale of the data \\
fewer than 2\textasciicircum{}32 voxels a reading & period.cu & the largest lattice a period is read on & a word's width: N² and Σc² must fit 64 bits; a larger lattice needs wider counts, not a refusal \\
\texttt{ANCHOR\_\allowbreak{}EXACT\_\allowbreak{}LIMBS} 256 in our build (anchor\_\allowbreak{}sift's default is 128) & \texttt{cell\_\allowbreak{}tracking/\allowbreak{}build\_\allowbreak{}engine.\allowbreak{}sh} and \texttt{build\_\allowbreak{}driver.\allowbreak{}sh}, fitted from \texttt{ENGINE\_\allowbreak{}RECORD\_\allowbreak{}LIMBS\_\allowbreak{}MOST} by a define, with no change to the source & 8,192 bits, 1,024 digits & a word's width, derived rather than chosen; \texttt{decimal\_\allowbreak{}double} reports \texttt{needed\_\allowbreak{}bits} past it \\
\bottomrule
\end{longtable}
}

Already ripped out: the npy zip directory cap (64 MiB), 23 September.

\section{\texorpdfstring{What the engine offers programs, and what it asks}{What the engine offers programs, and what it asks}}

{\small
\setlength{\tabcolsep}{6pt}
\begin{longtable}{>{\RaggedRight\arraybackslash}p{\dimexpr0.7441\linewidth-2\tabcolsep\relax}>{\RaggedRight\arraybackslash}p{\dimexpr0.2559\linewidth-2\tabcolsep\relax}}
\toprule
the engine offers & through \\
\midrule
\endhead
exact numbers from decimal strings, with \texttt{needed\_\allowbreak{}bits} & M1 \\
the residual at any orders and any input width the request names, its width returned & M2 \\
the component tree at every level: cuts, probe partitions, the bipartite census over probe links & M3 \\
moments per component & M4 \\
one correlation C and its reductions & M5 \\
the exact marginal and one-to-one matching on the caller's integer weights & M6, M7 \\
the entropy history and its cloud & M8 \\
proved files, and every source format read exactly & M9 \\
a seal on every crystal and set, and the place of any mismatch & M13 \\
one device shared between processes by memory, each job measured by pid & M14 \\
the sift: a pattern found by anchors whose count the corpus's coherence sets & M15 \\
a step's state drawn as a sheet or a volume, host and device alike & M16 \\
the record machine for the program's own verdicts & M10 \\
reversible lookup edges between the tower's floors, the root universal & M20 \\
a stateful error on every failure & M12 \\
length, find, copy, move, compare and fill at each instruction set's speed & M17 \\
\bottomrule
\end{longtable}
}

What the engine asks of a program: every scale (orders, windows, widths, spacings, lags, probes) in the request, and nothing about what the program's bodies are. Programs on it today: cell tracking (cell\_\allowbreak{}tracking\_\allowbreak{}table.md) and RSNA knee MRI.

\section{\texorpdfstring{What to work toward, in order}{What to work toward, in order}}

As of 23 September (the build plan, \hyperref[chap:build-plan]{build\_\allowbreak{}plan.md}):

\begin{enumerate}
\item \textbf{The seal} (M13): stage B (history, bodies, flattened, key, schedule), then the universal root and its anchor. The RSNA converter is done (M9); the export follows.
\item \textbf{Errors in every module} (M12), and its two fail-closed defects.
\item \textbf{The scale audit}, row by row: the entropy window into the request, ruling (a) at the three guards, the residual's width from the orders, the marginal's limits reported rather than set.
\item \textbf{Orders from the per-axis coherence reading} (M2), now built as M18 (\texttt{period\_\allowbreak{}read}, A12), so orders come from the data (order = period); the knee pulls it.
\item \textbf{The compression floor, by noise functionals} (M8, M9, M5) is its own theory: its plan and order of work are the compression book's, compression\_\allowbreak{}table.md.
\item \textbf{The algebra as passes}: the ladder-keeping sweep (A1), the LCA probe partition and the DFS-range overlap (A2), moments at every level (A3), C as the one pass (A4), the marginal over the whole component (A5).
\item \textbf{The port into anchor\_\allowbreak{}sift} (Doug, 23 September: close the gaps before anything moves). anchor\_\allowbreak{}sift's review found these gaps:

\begin{itemize}
\item (a) The layout is pinned to a commit hash before they re-check it.
\item (b) Narrowed 23 September. Our \texttt{base/\allowbreak{}no\_\allowbreak{}rounding} arms, \texttt{render} and \texttt{nbody/\allowbreak{}anchor\_\allowbreak{}sift} now carry their full doc text and are byte-identical to their files (their code was already identical). Our \texttt{exact\_\allowbreak{}integer.\allowbreak{}\{c,h\}} is ahead of theirs: it has the open width, the multiplication ladder, the division and Newton's division (M1), and theirs multiplies by long multiplication alone and has no division. This tree is the source their engine is to be copied from (anchor\_\allowbreak{}sift). The root \texttt{bench/\allowbreak{}} still differs from theirs in every file.
\item (c) Closed 23 September: \texttt{build\_\allowbreak{}engine.\allowbreak{}sh} and \texttt{build\_\allowbreak{}driver.\allowbreak{}sh} compile the exact integer from this tree's \texttt{base/\allowbreak{}no\_\allowbreak{}rounding} by default, and \texttt{ANCHOR\_\allowbreak{}EXACT\_\allowbreak{}ROOT} only overrides, so moving their \texttt{no\_\allowbreak{}rounding} no longer breaks our build.
\item (d) \texttt{engine/\allowbreak{}CMake\allowbreak{}Lists.\allowbreak{}txt} was their CMake verbatim and named paths this tree did not have. Fixed 23 September: every path is this tree's, the C-only warnings are kept off the CUDA sources, the CUDA arms join where a CUDA compiler is found, and \texttt{bench\_\allowbreak{}lattice} is built only where the compiler has C99 \texttt{\_\allowbreak{}Complex}. It builds and runs every bench on four configurations (M15).
\item (e) Their C builds host-only (CMake, MinGW, a Pi 5, WSL), but obsignatio, tessera and most modules are \texttt{.\allowbreak{}cu} only. Each module needs a host-only path, or CUDA optional per module.
\item (f) Their \texttt{src/\allowbreak{}engine/\allowbreak{}n\_\allowbreak{}body} holds older copies of modules we have (and some we retired). One spelling, and the old copies retired in the same change.
\item (g) The sift and the exact integer are outside M12's stateful errors: wrap at the boundary or wire them.
\item (h) Moved to \texttt{examples/\allowbreak{}anchor\_\allowbreak{}sift/\allowbreak{}} (Doug, 23 September): the Python fork of their \texttt{representation} package (games, particles, sound, text, proteins). A session copied it on 21 September, and nothing imports it. The engine holds no domain code.
\item Their proposed order: new modules land as siblings under their \texttt{src/\allowbreak{}engine/\allowbreak{}c}; the nbody modules then replace their \texttt{n\_\allowbreak{}body} copies; their own base/nbody/sims split comes last. Since gap (c) closed, that split no longer has to land with a change to our build.
\end{itemize}
\item \textbf{The two towers in the engine} (added 24 September; Doug: “You didn't build the inverse tower into the engine yet”). T and T⁻¹ emitted as record floors by the engine itself, not only by the tests, so that a program runs F ∘ T⁻¹ on the crystal as one stack, checked against \texttt{tower.\allowbreak{}cu}'s kernels lane for lane (A13; \hyperref[chap:vertical-time-compression]{vertical\_\allowbreak{}time\_\allowbreak{}compression.md}).
\end{enumerate}
